% concatenated from main.tex
\documentclass{article}
\usepackage{graphicx}
\usepackage{fullpage}
\usepackage{placeins}

\usepackage{amsthm}
\usepackage{amsmath}
\usepackage{amssymb}
\usepackage{amsfonts}
\usepackage{bm}
\usepackage{url}
\usepackage[top=3cm,bottom=3cm,left=2.5cm,right=2.5cm]{geometry}
\usepackage{enumitem}
\usepackage[dvipsnames]{xcolor}
\usepackage{graphicx}
\usepackage{multirow}
\usepackage{booktabs}
\usepackage{soul}
\usepackage{subfigure}
\usepackage{rotating}
\usepackage{longtable}
\usepackage{dsfont}
\usepackage[normalem]{ulem}
\usepackage{xparse}
\usepackage{algorithm}
\usepackage{algpseudocode}
\usepackage{booktabs}
\usepackage{hyperref}
\usepackage{cleveref}
\usepackage{authblk}
\usepackage{array}
\usepackage{cases}


\NewDocumentCommand{\cg}{ O{L} O{L_1} O{L_2} m m m }{C^{#1 #4}_{#2 #5 #3 #6}}
\newcommand{\vrho}[2]{\varrho^{#1}_{#2}}
\newcommand{\tvrho}{\tilde\varrho^i_j}
\newcommand{\rep}[2]{\rho^{#1}_{#2}}
\newcommand{\blamb}[2]{\bm\lambda^{#1}_{#2}}
\newcommand{\bphi}[2]{\bar\phi^{#1}_{#2}}
\newcommand{\bpsi}[2]{\bar\psi^{#1}_{#2}}
\newcommand{\bV}{{\bar V}}
\newcommand{\tblamb}{\bm{\tilde\lambda}^{i}_{\bbm}}
\newcommand{\dd}{{d}}
\newcommand{\card}[1]{|{#1}|}
\newcommand{\dimge}[2]{\dim(V^{#1,#2})}
\newcommand{\dimpi}[2]{\dim(\bV^{#1,#2})}
\newcommand{\dimgeest}[2]{\dim^{\textup{est}}(V^{#1,#2})}
\newcommand{\dimpiest}[2]{\dim^{\textup{est}}(\bV^{#1,#2})}

\newtheorem{theorem}{Theorem}[section]
\newtheorem{proposition}[theorem]{Proposition}
\newtheorem{assumption}[theorem]{Assumption}
\newtheorem{definition}[theorem]{Definition}
\newtheorem{corollary}[theorem]{Corollary}
\newtheorem{remark}[theorem]{Remark}
\newtheorem{example}[theorem]{Example}
\newtheorem{lemma}[theorem]{Lemma}
\newtheorem{conjecture}[theorem]{Conjecture}

\numberwithin{equation}{section}

\newcommand{\cgd}[1]{ \textcolor{red}{ \texttt{GD: #1}}}
\newcommand{\gd}[1]{ \textcolor{red}{{#1}}}
\newcommand{\lz}[1]{ \textcolor{cyan}{{#1}}}
\newcommand{\clz}[1]{ \textcolor{cyan}{\texttt{[LZ: #1]}}}
\newcommand{\ch}[1]{ \textcolor{orange}{ \texttt{[CH: #1]}}}
\newcommand{\cht}[1]{ \textcolor{orange}{{#1}}}


\def\N{\mathbb{N}}
\def\K{\mathbb{K}}
\def\Z{\mathbb{Z}}
\def\R{\mathbb{R}}
\def\C{\mathbb{C}}
\def\F{\mathbb{F}}
\def\bbC{\mathbb{C}}
\def\bbS{\mathbb{S}}

\def\equ{{K}}
\def\Nb{{N_{\rm blocks}}}
\def\bc{{\bm c}}
\def\btc{\tilde{\bm c}}
\def\br{{\bm r}}
\def\bk{{\bm k}}
\def\bl{{\bm l}}
\def\bbm{{\bm m}}
\def\bu{{\bm u}}
\def\bx{{\bm x}}
\def\bz{{\bm z}}
\def\bp{{\bm p}}
\def\bq{{\bm q}}
\def\bmu{{\bm \mu}}
\def\bn{{\bm n}}
\def\bnu{{\bm \nu}}
\def\bL{{\bm L}}
\def\bll{{{\bm l}_1}}
\def\blr{{{\bm l}_2}}
\def\bnl{{{\bm n}_1}}
\def\bnr{{{\bm n}_2}}
\def\bbml{{{\bm m}_1}}
\def\bbmr{{{\bm m}_2}}
\def\bLl{{{\bm L}_1}}
\def\bLr{{{\bm L}_2}}
\def\bR{{\bm R}}
\def\equl{{\equ_1}}
\def\equr{{\equ_2}}
\def\obm{{\overline\bbm}}
\def\obmu{{\overline\bmu}}



\def\calD{\mathcal{D}}
\def\calB{\mathcal{B}}
\def\calU{\mathcal{U}}
\def\calP{\mathcal{P}}
\def\calN{\mathcal{N}}
\def\calR{\mathcal{R}}
\def\calM{\mathcal{M}}
\def\calE{\mathcal{E}}
\def\calV{\mathcal{V}}
\def\calK{\mathcal{K}}
\def\calF{\mathcal{F}}
\def\calC{\mathcal{C}}
\def\calL{\mathcal{L}}
\def\calS{\mathcal{S}}
\def\calLG{\mathcal{\calL_{\rm G}}}
\def\Ml{\calM_\bl}
\def\Mnl{\calM_{\bn\bl}}
\def\oMl{\overline{\calM}_\bl}
\def\oMlK{\overline{\calM}_{\bl,K}}
\def\oMnl{\overline\Mnl}
\def\MlL{{M}}
\def\Ll{\calL_\bl}
\def\calKL{\calK_{\bl,L}}

\def\intD{\bar{D}}
\def\diagD{{\bm D}}

\def\Ib{I_b}

\def\Asym{\mathcal{A}}

\def\rcut{r_{\rm cut}}
\def\ro{{r_0}}
\def\fcut{f_{\rm cut}}
\def\rnn{r_{\rm nn}}

\def\<{\langle}
\def\>{\rangle}

\def\b{\big}
\def\B{\Big}
\def\bg{\bigg}
\def\calA{\mathcal{A}}
\def\calI{\mathcal{I}}
\def\RI{{\rm RI}}
\def\Nat{{N}}
\def\MS{{\textup{MS}(\Omega)}}
\def\dim{{{\rm dim}}}
\def\Mpm{M^{\pm}}
\def\Mp{M^{+}}
\def\Mm{M^{-}}
\def\Apm{A^{\pm}}
\def\Ap{A^{+}}
\def\Am{A^{-}}
\def\Bpm{B^{\pm}}
\def\Bp{B^{+}}
\def\Bm{B^{-}}
\def\Mup{{\MlL}^{\textup{up}}}
\def\Mlow{{\MlL}^{\textup{down}}}


\def\Spec{{\bf Spec}}
\def\SpecV{{\bf Spec}^V}

\def\sep{\,|\,}
\def\bsep{{\,\b|\,}}
\def\Bsep{\,\B|\,}

\title{Efficient construction and explicit dimensionality of Lie group-equivariant and permutation-invariant spaces}


\author[1]{Eloïse Barthelemy}
\affil[1]{Université Marie et Louis Pasteur, CNRS, LmB (UMR 6623), F-25000 Besançon, France}

\author[1]{Genevi\`eve Dusson\footnote{genevieve.dusson@math.cnrs.fr}}

\author[1]{Camille Hernandez}
  
\author[2]{Liwei Zhang}
\affil[2]{Institut f\"ur Geometrie und Praktische Mathematik, RWTH Aachen University, Im Süsterfeld 2, 52072 Aachen, Germany }

\date{}

\begin{document}

\maketitle

\begin{abstract}
We introduce a practical construction of group-equivariant and 
permutation-invariant functions of $\Nat$ variables 
given a finite-dimensional space stable with respect to the group action.
The construction applies to any connected linear Lie group
and relies on leveraging the Lie algebra to build a matrix $M$
whose kernel is in one-to-one correspondence with the subspace with desired equivariance and invariance properties, removing the need for prior knowledge of Clebsch--Gordan coefficients. 
A similar construction is proposed for group-equivariant functions alone, without imposing permutation-invariance.
For the groups $SO(3)$ and $SU(2)$, we further exploit the structure of the Lie algebra to demonstrate the sparsity pattern and rank of the matrix $M$, which yields the exact dimension of the group-equivariant and permutation-invariant space, as well as the dimension of the group-equivariant space alone.
We demonstrate analytically and verify numerically that the proposed method scales linearly with respect to the dimensionality of the basis, 
offering a high computational gain compared to existing methods in the literature which typically scale exponentially.
We finally perform a dimensionality comparison, showing that for large values of~$\Nat$, the dimension of group-equivariant and permutation-invariant spaces is of comparable order as the dimension of permutation-invariant spaces, while pre-asymptotically,
the first dimensionality is orders of magnitude lower than the second. Hence a substantial computational gain can be achieved by explicitly enforcing group-equivariance on top of permutation-invariance when approximating such functions.
\end{abstract}

\vspace{1em}
\noindent\textbf{Keywords} Permutation-invariance,  Lie group-equivariance, Generalized Clebsch-Gordan coefficients, Numerical linear algebra
\\

\noindent\textbf{Mathematics Subject Classification} 65D15, 65Y20, 46N50, 81R05

\tableofcontents



\section{Introduction}




In this work, we are interested in $\Nat \in \N = \{1,2,\ldots\}$ variable functions satisfying either an equivariance property with respect to some Lie group action (group-equivariance; GE), or an invariance with respect to the permutation of the variables on top of the Lie group equivariance (group-equivariance and permutation-invariance; GE-PI). 
Such functions are fundamental across diverse scientific fields, from physics and chemistry to materials science. They are especially crucial for studying interacting particle systems, where permutation-invariance arises from considering identical particles.
Often Lie group equivariance naturally arises from the physical setting, e.g. $O(3)$ in atomistic modeling~\cite{Batatia2022-zi}, $SO^+(1,3)$ in high-energy physics~\cite{Munoz2022-ur}, or $SU(n)$ in quantum mechanics~\cite{Lagrave2021-cw,Favoni2022-zb}.
When approximating such functions, constructing a parametrization that respects the inherent symmetries is essential. 
This approach typically reduces the number of required parameters, making the approximation more efficient~\cite{Bachmayr2024-pv}.
As a toy example, the dimension of the space of $N$-variable polynomials with 
given total degree is asymptotically $\Nat!$ times larger than the corresponding number of totally symmetric polynomials.
While permutation-invariant functions are well-understood and can be efficiently constructed with symmetric polynomials and variations thereof, 
efficiently constructing functions satisfying both permutation-invariance and group-equivariance is in general difficult; this is the main motivation for this article.


To do so, we consider a  connected linear
Lie group $G$ with neutral element $e$ and Lie algebra $\mathfrak g$.
We then consider  a vector space
$\F$
over a field $\K$, typically $\R$ or $\C$ in the examples below, $\rho$ a continuous representation of $G$ of dimension $K$ on $\K$, $\Omega$ a spatial domain, and $\cdot$ a group action for $G$ on $\Omega$:
\[
    \left\{ 
    \begin{array}{lc}
        \forall \br\in \Omega, g\in G,  &  g\cdot \br \in \Omega, \\
       \forall \br\in \Omega,  & e \cdot \br = \br, \\
       \forall g_1,g_2\in G, & 
       (g_1 g_2) \cdot \br = g_1 \cdot (g_2 \cdot \br).
    \end{array}
    \right.
\]
We define the corresponding group action over $\Omega^\Nat$, which is still denoted by $\cdot$, 
as 
\[
    \forall \bR = (\br_1,\ldots,\br_\Nat) \in \Omega^\Nat, \; 
     g \in G, \quad 
    g \cdot \bR = 
    (g \cdot \br_1, 
    \ldots, 
    g \cdot \br_\Nat).
\]
GE  and GE-PI functions are therefore defined as follows.
\begin{definition}[Group-equivariant (GE) function]
    A function $F: \Omega^\Nat \rightarrow \F^K$ is group-equivariant for the representation $\rho$ if 
\begin{equation}\label{eq:GE}
  \forall \bR \in \Omega^\Nat, \;  g \in G,  \quad   F(g \cdot \bR) = \rho(g) \; F(\bR).
\end{equation}
\end{definition}


\begin{definition}[Group-equivariant and permutation-invariant (GE-PI) function]
\label{def:GEPI}
A function $F: \Omega^\Nat \rightarrow \F^K$ is group-equivariant for the representation $\rho$
and permutation-invariant if
\begin{equation}\label{eq:GE-PI}
  \forall \bR \in \Omega^\Nat, \; g \in G, \; \sigma \in S_\Nat, \quad   F(g \cdot \bR_\sigma) = \rho(g) \; F(\bR),
\end{equation}
where $S_\Nat$ denotes the symmetric group of degree $\Nat$, and
\[
    \forall \bR = (\br_1,\ldots,\br_\Nat) \in \Omega^\Nat, 
    \quad 
    \bR_\sigma = (\br_{\sigma_1},\ldots,\br_{\sigma_\Nat}).
\]
\end{definition}

In some contexts, e.g., for functions defined on multisets, the permutation invariance of the functions is hidden in the domain of definition. Precisely, denoting by $\textup{MS}(\Omega)$ the union of multisets of $\Omega$ with arbitrary multiplicity, that is 
\[
    \textup{MS}(\Omega) 
    = \bigcup_{n=1}^{+\infty} 
    \Big\{ 
    \left\{ 
    \br_1, \br_2, \ldots, \br_n
    \right\}, 
    \br_1,\br_2,\ldots, \br_n \in \Omega
    \Big\},
\]
a function $F:\textup{MS}(\Omega)\to\F^K$ satisfying
\[
    \forall \bR \in \textup{MS}(\Omega), \;  g \in G,  \quad   F(g \cdot \bR) = \rho(g) \; F(\bR),
\]
for the representation $\rho$ with dimension $K$ is also said to be GE-PI.


The aim of this manuscript is therefore the following: \textit{Given a finite-dimensional space that is stable with respect to the group action, efficiently determine the subspace of GE or GE-PI functions in this space, and provide its dimension.}


To begin with, we consider one-variable function spaces, sometimes also called one-particle function spaces~\cite{Dusson2022-ie,Batatia2023-us}, stable with respect to the group action, and denoted by
\[
    V^l =  {\rm Span}\{ \phi^l_m, \; m\in  \calM_l \}, 
\]
 where the index $l$ lies in an index set $\calL$ which is typically countable and totally ordered (e.g. $\N_0 = \{0,1,\ldots\}$ for the group $SO(3)$, see Section~\ref{sec:compatible_bases} for more examples). Here $\calM_l$ denotes a finite totally ordered index set depending on $l$, typically $\{ 1,2, \ldots, \card{\calM_l} \}$. The functions $\phi^l_m:\Omega \rightarrow \F$ are defined on a spatial domain $\Omega$ in a space endowed with a scalar product denoted by $\langle \cdot, \cdot \rangle$, typically $\langle \cdot, \cdot \rangle_{L^2(\Omega)}$,
and
are assumed for simplicity to be orthonormal, i.e.,
\begin{equation}
\forall l,l' \in \calL, \;
     m,m' \in \calM_l,
\quad 
\langle\phi^l_m,\phi^{l'}_{m'}\rangle
=\delta_{ll'}\delta_{mm'}.
\label{eq:orthogonality}
\end{equation}
The space $V^l$ being stable with respect to the group action means that
\[
    \forall g \in G,
    \quad \br \mapsto \phi^l_m(g \cdot \br) \in V^l.
\]
This implies that there exists a representation of $G$ of dimension $\card{\calM_l}$, denoted by $\rep{l}{}$, such that
\begin{equation}\label{eq:1b-basis}
   \forall g\in G, \;  \br \in \Omega, \quad  \phi^l_m(g \cdot \br) = \sum_{m' \in \calM_l} \rep{l}{m,m'}(g) \phi^l_{m'}(\br).
\end{equation}
Indeed, by the uniqueness of the coordinates in the basis 
$\{ \phi^l_m \}_{m\in \calM_l}$, we have
$\rep{l}{}(e) = I_{\calM_l}$, and 
\[
    \forall g_1,g_2\in G, \quad \rep{l}{}(g_1 g_2) = \rep{l}{}(g_1) \rep{l}{}(g_2).
\]
We call $\{\phi^l_m\}_{m\in\calM_l}$ a compatible basis of $V^l$ for the representation $\rep{l}{}$. 

We then consider for $\bl \in \calL^\Nat$ an $\Nat$-variable function space as the
tensor product space 
\[
    V^{\bl} = \bigotimes_{i=1}^\Nat V^{l_i},
\]
which due to~\eqref{eq:1b-basis} is stable by group action. 
An orthonormal basis for $V^{\bl}$ is 
\begin{equation}
\label{eq:calMbl}
       \left\{  \phi^\bl_\bbm := \prod_{i=1}^\Nat \phi^{l_i}_{m_i}, \quad \bbm \in \calM_\bl \right\}, \quad  \calM_\bl := \bigotimes_{i=1}^\Nat \calM_{l_i},
\end{equation}
where the corresponding scalar product on $V := \bigoplus_{\bl\in\calL^\Nat} V^{\bl}$ is defined as 
\[
  \forall  \phi^\bl_\bbm, \; \phi^{\bl'}_{\bbm'} \in V,
  \quad \langle \phi^\bl_\bbm, \phi^{\bl'}_{\bbm'} \rangle 
    = \prod_{i=1}^\Nat \langle \phi^{l_i}_{m_i}, \phi^{l'_i}_{m'_i} \rangle =\delta_{\bl\bl'}\delta_{\bbm\bbm'}.
\]
To characterize GE (respectively GE-PI) functions whose components are in $V^{\bl}$, we derive a linear system based on~\eqref{eq:GE} or~\eqref{eq:GE-PI} that has to be satisfied for any element of the group $G$, similarly to~\cite{Batatia2023-us}. 
Differentiating this linear system and evaluating it at the generators of the Lie algebra $\mathfrak g$ gives rise to another, way simpler linear system $M \bc = 0$: 
a basis of the kernel of the matrix $M$ is in one-to-one correspondence with a basis of the space of GE (resp. GE-PI) functions having components in $V^{\bl}$.
Indeed since $G$ is a  connected linear Lie group, the continuous representations $\rep{l}{}$ are differentiable and the derivative of $\rep{l}{}$ at the neutral element $e\in G$ is defined as (see~\cite[Theorem 3.28]{Hall2013-af})
\begin{equation}\label{eq:rep-deriv}
    \forall X\in\mathfrak g,\quad  d\rep{l}{}(X) = \left.\frac{d}{dt}\right|_{t=0} \rep{l}{}(\exp(t  X)).
\end{equation}
The linear system therefore involves terms depending on 
the derivatives of $\rep{l}{}$ at a set of generators of the Lie algebra $X_\dd\in \mathfrak g$ 
for $l \in \calL$ and $d = 1, \ldots, N_{\rm{dim}}$, denoted by
\[
    \vrho{l,\dd}{}=[\vrho{l,\dd}{m,m'}]_{m,m' \in \calM_l}:= [d\rep{l}{}(X_\dd)]_{m,m' \in \calM_l}.
\]
We then prove that reciprocally, if the linear system $M\bc =0$ is satisfied, the corresponding functions are GE or GE-PI according to the setting, guaranteeing that a basis of GE or GE-PI space is obtained. 
Moreover, the matrix $M$ is usually very sparse (c.f. equations~\eqref{eq:M_GE} and~\eqref{eq:M_GE_PI}), making the computation of its kernel inexpensive. In particular, a major numerical gain can be expected for the GE-PI case, as the matrix is expressed independently of permutations, which significantly reduces the computational cost compared to existing literature as detailed below.

\paragraph{Main contributions} The main results of this article are the following.
\begin{enumerate}
    \item We provide in Proposition~\ref{prop:GE_basis} sufficient and necessary conditions to obtain a basis of GE functions with components in a given finite-dimensional space stable with respect to the group action.
    \item We present in Theorem~\ref{thm:GEPI_basis} similar conditions on the more challenging case of GE-PI functions which completely avoids explicitly symmetrizing over permutations.
    \item We propose an alternative recursive construction for GE and GE-PI spaces in Propositions~\ref{prop:GE_rec} and~\ref{prop:GEPI_rec}, the latter 
    similarly
    avoiding an explicit symmetrization over permutations.
    \item For the groups $SU(2)$ and $SO(3)$, we simplify the linear systems of  Proposition~\ref{prop:GE_basis} and Theorem~\ref{thm:GEPI_basis}. We show that the corresponding matrix is extremely sparse and characterize the dimensionality of the kernel a priori, providing an explicit dimensionality of GE and GE-PI spaces in Proposition~\ref{prop:ge-su2} and Theorem~\ref{thm:gepi-su2}.
    We also provide asymptotic dimensionalities in Proposition~\ref{prop:asym-dim-ge}, and Proposition~\ref{prop:asym-dim-gepi}.
    \item We carry out a theoretical complexity analysis showing that the numerical method of constructing GE-PI bases scales linearly with respect to the number of basis functions in Section~\ref{sec:efficiency}, overcoming the exponential complexity of existing methods, and we numerically observe  the linear complexity.
    \item We benchmark our method against several packages from the literature (see Section~\ref{sec:num}) and show that the proposed method is substantially faster and can deal with a large number of variables $N$ where other methods are too expensive to even run the calculations.
    \item Finally, we provide a comparison between the dimensionalities of non-symmetric, GE, PI, and GE-PI spaces in Section~\ref{sec:dimensionalities}.
\end{enumerate}


\paragraph{Prior work}

GE and GE-PI functions are 
used in different contexts. Related to the groups $SO(3)$ and $O(3)$, such functions appear in e.g. the construction of interatomic potentials~\cite{Schutt2018-dw, Drautz2019-mn, Batzner2022-az, Batatia2022-zi}, Hamiltonians~\cite{Zhang2022-ay,Li2022-cj,Zhong2023-ar,Suman2025-qz,Qian2025-du,gong2023general}, density matrices~\cite{Wetherell2020-us,Hazra2024-cy,Zhang2025-ve}, wavefunctions~\cite{zhou2024-wf,unke2021se}, friction tensors in Langevin heatbaths~\cite{Sachs2025}, or protein structure~\cite{lu2024dynamicbind,eismann2021hierarchical}. 
GE-PI functions related to $SO^+(1,3)$ are used for the classification of high-energy particles and the regression of their properties (e.g. mass)~\cite{Bogatskiy2020-jt,Munoz2022-ur,Bogatskiy2024-oc}.
The corresponding proposed parametrizations vary across the literature, and are typically either linear in a chosen symmetry-adapted basis~\cite{Munoz2022-ur,Witt2023-rm,Sachs2025,Zhang2025-ve}, or rely on the use of equivariant neural networks~\cite{Batatia2022-zi,geiger2022e3nn,lu2024dynamicbind}.

Most of the works cited above build on the construction of so-called many-body equivariant features, which consist of linearly independent basis functions equivariant 
for specific group representations.
When considering functions of $\Nat$ variables satisfying GE but not PI, linearly independent equivariant features correspond to so-called generalized Clebsch--Gordan coefficients, already defined in 1962~\cite[Chapter 2]{Yutsis1962-wv}. A practical method for computing these coefficients was, e.g., proposed in~\cite{Dusson2022-ie} within the framework of the Atomic Cluster Expansion (ACE).
However, an extra difficulty arises when permutation-invariance is enforced, as the systematic construction of generalized Clebsch--Gordan coefficients in this case typically involves evaluating a Gramian matrix to remove linear dependencies, 
which requires performing a sum over permutations leading to an exploding cost when the number of variables $\Nat$ grows~\cite{Dusson2022-ie}. Alternative methods have been proposed such as~\cite{Nigam2020-zj,Batatia2022-zi,
Goff2024-sw,geiger2022e3nn,geiger2024accelerate} for groups like $SO(3)/SU(2)$ and~\cite{Batatia2023-us} in the larger framework of reductive Lie groups but to the best of our knowledge, the computational cost always scales exponentially with respect to the space dimension and therefore the construction becomes costly and even unfeasible for large values of $\Nat$. 
Note that, compared to other methods, such as~\cite{Nigam2020-zj,Dusson2022-ie}, we do not need Clebsch--Gordan coefficients which characterize 2-variable GE functions as a prerequisite to compute generalized Clebsch--Gordan coefficients. In the numerical Section~\ref{sec:num} we benchmark our method against four of the existing packages mentioned above~\cite{geiger2022e3nn,Batatia2022-zi,Batatia2023-us,geiger2024accelerate}.


Regarding the explicit dimensionalities obtained in this work for GE and GE-PI in the case of rotation groups, previous results were obtained by one of the authors for GE spaces in~\cite{Dusson2022-ie} in a less explicit way, see Remark~\ref{rem:ace}. For GE-PI functions, few results on the dimensionality are presented in~\cite{Zhao2005-zx,Talmi2005-gu,Pain2019-ft} for specific representations, but without the construction of the corresponding GE-PI basis, mostly relying on characters.
As a direct outcome of exact dimensionality and dimensionality estimates, we can compare the cardinality
of bases with different symmetries (no symmetry, PI, GE, GE-PI), complementing the asymptotic results presented in~\cite{Bachmayr2024-pv} restricted to the PI case.
Indeed asymptotically the number of PI basis functions is comparable to the number of GE-PI basis functions while preasymptotically, which is of practical interest, the respective numbers of PI and GE basis functions are much larger than the number of GE-PI basis functions.
This is 
particularly useful to contribute to the growing discussion whether known symmetries need to be put in the model as a constraint or if data-augmentation is sufficient and can be used as a replacement. For example, for $SO(3)$ the work~\cite{Bigi2026-yx} discusses how symmetry-unconstrained models can still be accurate, while the article~\cite{Klinteback2026-xg} shows the high cost of data augmentation. There is a similar discussion for the Lorentz group in~\cite{Li2024-bm}.



\paragraph{Outline}

The outline of the article is as follows.
In Section~\ref{sec:equivariant_lie_algebra} we present two methods for obtaining a basis of GE and GE-PI functions: a direct and a recursive one. We apply the method to $SO(3)$ and $SU(2)$ groups and extend it to $O(3)$, and provide explicit dimensionality as well as asymptotics in this case in Section~\ref{sec:appl}. We then provide numerical results in Section~\ref{sec:num}, showcasing that the proposed method is more efficient than existing methods from the literature. 
We also provide a complexity analysis and comparison of GE and GE-PI dimensionalities. 
Finally we present concluding remarks in Section~\ref{sec:conclusion}
and gather the proofs of the main results in Section~\ref{sec:proofs}.

\section{Generating equivariant bases using the Lie Algebra}
\label{sec:equivariant_lie_algebra}

In this section we determine group-equivariant basis functions (both GE and GE-PI)
via the kernel of a specific matrix 
based on the Lie algebra of the group.

\subsection{Group-equivariant basis}

In the following we provide necessary and sufficient conditions for characterizing a basis of GE functions for a representation $\rep{L}{}$ in the space $[V^\bl]^{\dim(\rep{L}{})}$ for a given $\bl\in\calL^\Nat$.
Unless otherwise specified, we assume $L\in\calL$, so that $\calM_L$ is well-defined and $\dim(\rep{L}{})=\card{\calM_L}$. We denote the space of desired GE functions by $V^{\bl,L}$, and characterize a basis of this space in the following proposition.


\begin{proposition}[Basis of the GE space $V^{\bl,L}$]
    \label{prop:GE_basis}
    Let $\bl\in\calL^\Nat$, $L\in\calL$, and let
    \begin{equation}
     \label{eq:matM}
         M^{\bl,L} = \begin{pmatrix}
         M_1 \\
         M_2 \\
         \vdots \\
         M_{N_{\rm{dim}}}
     \end{pmatrix},
     \end{equation}
     where for $\dd = 1,\ldots, N_{\rm{dim}}$,
     the elements of the submatrices 
     $M_\dd$ are defined for $(\bbm,k),\; (\bbm',k') \in \calM_\bl\times\calM_L$ as
    \begin{equation}
        M_\dd[(\bbm,k),(\bbm',k')] = 
        \begin{cases}
        \displaystyle\sum_{j=1}^\Nat
        \vrho{l_j,\dd}{m_j,m_j} - \vrho{L,\dd}{k,k}, \qquad &\textup{if}\;\; \bbm=\bbm',  \; \; k=k', \\
         -\vrho{L,\dd}{k,k'}, \qquad &\textup{if}\;\; 
         \bbm = \bbm',  \; \;  k\ne k', \\
         \vrho{l_j,\dd}{m'_j,m_j}, \qquad &\textup{if}\;\; 
         \forall s\in \{1,\ldots,\Nat\} \neq j,\; m_s = m'_s, \;\; m_j\ne m'_j, \;\;  k = k',\\
        0, \qquad &\textup{otherwise.} \\
    \end{cases}
    \label{eq:M_GE}
    \end{equation}
    Then, a basis of the GE space $V^{\bl,L}$
    is given by
    \begin{equation*}      
        b^{\bl,L}_i = 
        \left[ 
            \sum_{\bbm \in \calM_\bl}c_{\bbm,k,i}^{\bl,L}\phi^\bl_\bbm 
        \right]_{k\in \calM_L},\;\; i=1,\ldots,\dim\left(\ker(M^{\bl,L})\right),
    \end{equation*}
    where $\bc^{\bl,L}_i = \{c_{\bbm,k,i}^{\bl,L}\}_{\bbm\in\calM_\bl,k\in\calM_L},\; i=1,\ldots,\dim\big(\ker(M^{\bl,L})\big)$ is a basis of $\ker(M^{\bl,L})$. 
    Thus,  
    \[
        \dimge{\bl}{L}=\dim\big(\ker(M^{\bl,L})\big).
    \]
\end{proposition}

We sometimes refer to the coefficients $\bc^{\bl,L}_i$ as coupling coefficients. The proof of Proposition~\ref{prop:GE_basis} can be found in Section~\ref{sec:proof_equivariant_lie_algebra}. For ease of notation, the superscript is sometimes omitted from the matrix $M^{\bl,L}$ when the associated $\bl$ and $L$ are clear from the context.
Using the definition of the matrix $\MlL$~\eqref{eq:matM}, its kernel is the intersection of the kernels of its submatrices, that is,
\begin{equation}
\label{eq:ker_inclusion}
    {\rm Ker}(M) = \bigcap_{d=1}^{N_{\rm dim}} {\rm Ker}(M_d).
\end{equation}
Note that $M_\dd$ can be written as a sum of Kronecker products, 
\begin{equation}
     M_\dd = \sum_{j=1}^\Nat 
    \left[ 
    \bigotimes_{s=1}^{j-1}I_{\card{\calM_{l_s}}}\otimes (\vrho{l_j,\dd}{})^T\bigotimes_{s=j+1}^{\Nat}I_{\card{\calM_{l_s}}}
    \right]
    \bigotimes I_{\card{\calM_L}} - \bigotimes_{j=1}^\Nat I_{\card{\calM_{l_j}}}\otimes \vrho{L,\dd}{},
    \label{eq:Mp}
\end{equation}
and its kernel can be related to the tensor products of eigenvectors of matrices $\vrho{l,\dd}{}$, as presented in the next proposition. 
Note also that the linear system of Proposition~\ref{prop:GE_basis} can be expressed as $N_{\rm dim}$ Sylvester equations.

\begin{proposition}
\label{prop:kernel_sum}
    For any $l\in \calL$, $d \in \{ 1, \ldots, N_{\rm dim} \}$, let $(\epsilon^{l,d}_j, u^{l,d}_j, v^{l,d}_j) \in \F \times \F^{|\calM_l|}\times \F^{|\calM_l|}$ be respectively the eigenvalues, left eigenvectors, and right eigenvectors  of the matrices $\vrho{l,\dd}{}$ for $j \in \calM_{l}$. Then the kernel of the matrix $M_d$ defined in~\eqref{eq:M_GE} is spanned by all
    \[
        u^{l_1,d}_{j_1} \otimes 
        u^{l_2,d}_{j_2} \otimes 
        \ldots 
        \otimes 
        u^{l_N,d}_{j_N} 
        \otimes v^{L,d}_{j},
    \]
    such that 
    \[
    \epsilon^{l_1,d}_{j_1} + \epsilon^{l_2,d}_{j_2} + \ldots 
    + \epsilon^{l_\Nat,d}_{j_\Nat}
    -  \epsilon^{L,d}_{j} = 0.
    \]
\end{proposition}


\begin{proof}
By using the expression of the matrices $M_d$ in~\eqref{eq:Mp}, the eigenvectors of the matrices $M_d$
are $u^{l_1,d}_{j_1} \otimes 
        u^{l_2,d}_{j_2} \otimes 
        \ldots 
        \otimes 
        u^{l_N,d}_{j_N} 
        \otimes v^{L,d}_{j}$
        for all possible indices $j_i \in \calM_{l_i}$ for $i=1,\ldots,\Nat$ and $j\in \calM_{L}$,
        with corresponding eigenvalues 
        $\epsilon^{l_1,d}_{j_1} + \epsilon^{l_2,d}_{j_2} + \ldots 
    + \epsilon^{l_\Nat,d}_{j_\Nat}
    -  \epsilon^{L,d}_{j}$.
    Therefore the kernel of $M_d$ corresponds to the space spanned by eigenvectors with zero eigenvalues.
\end{proof}

In principle, one could determine the kernel either using Proposition~\ref{prop:GE_basis}
or Proposition~\ref{prop:kernel_sum}. While we mostly use Proposition~\ref{prop:GE_basis} to numerically determine the coupling coefficients in practice, Proposition~\ref{prop:kernel_sum} could be useful to obtain some explicit dimensionality of GE spaces.

We end this section with a remark providing a link with Clebsch--Gordan coefficients.

\begin{remark}
\label{rem:CG}
    In the case where $\bl = (l_1,l_2)$, the coupling coefficients correspond to the well-known Clebsch--Gordan coefficients. For larger $\bl$'s, these coupling coefficients are sometimes called generalized Clebsch--Gordan coefficients~\cite{Dusson2022-ie,Batatia2023-us}. 
\end{remark}


\subsection{Group-equivariant and permutation-invariant basis}\label{sec:ge-pi}

We now turn to characterize GE-PI functions (see Definition~\ref{def:GEPI}).
This is essential in practice, since incorporating the permutation invariance into modeling can significantly reduce the dimension of the targeted function space, especially when high-dimensional problems are considered. 
As a first result, we provide a basis for PI functions in $V^\bl$ for some $\bl \in \calL^\Nat$.
To this end, we define
\[
    \calS_{\bl}=\{\pi \in S_\Nat \; | \;
    \pi\bl=\bl\}\subset S_\Nat,
\]
and 
\begin{equation}\label{eq:GE_class}
    \oMl:=\{\overline{\bbm}\;|\;\bbm\in \calM_{\bl}\},
\end{equation}
where $\overline{\bbm}$ stands for the equivalent class of $\bbm$ with the equivalent relation $\sim$ being defined as 
\[
    \bbm\sim\bbm^{\prime}  \Leftrightarrow \exists \; \pi\in \calS_{\bl}, \quad \bbm^{\prime}=\pi\bbm.
\]
For simplicity of notation, and as it is often clear from the context, we still denote the classes in $\oMl$ by $\bbm$. 
\begin{proposition}[Basis of PI functions in $V^\bl$]\label{prop:basis-pi}
    Let $\bl\in\calL^\Nat$. 
    A basis of PI functions belonging to $V^{\bl}$ consists of 
\begin{equation}
\label{eq:PI_Nbody}
        \left\{ \displaystyle\bpsi{\bl}{\bbm}: =\sum_{\sigma\in S_\Nat}\phi^\bl_\bbm\circ\sigma, \quad \bbm\in \oMl \right\}.
\end{equation}
In particular, the dimension of the space of PI functions in $V^\bl$ is $\card{\oMl}$.
\end{proposition}

The proof can be found in Section~\ref{sec:proof_equivariant_lie_algebra}.

Since the evaluation of basis functions defined in~\eqref{eq:PI_Nbody} scales exponentially with $N$ 
due to the sum over permutations, 
a density projection trick was proposed~\cite{Drautz2019-mn,Dusson2022-ie}, which relies on using another set of $N$-variable PI functions
\begin{equation}
      \forall l\in \calL, \; m\in \calM_l, \quad  
  \bphi{l}{m} : \bR = (\br_1,\ldots,\br_\Nat) \in \Omega^\Nat \mapsto \sum_{i=1}^\Nat \phi^l_m(\br_i).
  \label{eq:sumphi}
\end{equation}
The considered PI function space is then 
\[
    \bV^l = {\rm Span} \{
    \bphi{l}{m}, \; 
    m \in \calM_l
    \}.
\]
Similar to $V^{\bl}$, the space of PI functions $\bV^{\bl}$ for $\bl\in \calL^\Nat$ can be defined by forming tensor products of $\bV^{l_i}$, as 
\[
  \bV^{\bl} := \bigotimes_{i=1}^\Nat \bV^{l_i} = {\rm Span} \left\{  \bphi{\bl}{\bbm} := \prod_{i=1}^\Nat \bphi{l_i}{m_i}, \; 
  \bbm \in \calM_\bl \right\}, \quad  \calM_\bl := \bigotimes_{i=1}^\Nat \calM_{l_i}.
\]
As indicated in~\cite{Dusson2022-ie}, the evaluation cost of the $\bphi{\bl}{\bbm}$ scales linearly with respect to $\Nat$.
Similarly as with $\bpsi{\bl}{\bbm}$, some functions in this set are identical. We therefore restrict $\bbm\in\oMl$, and obtain the following result.


\begin{proposition}[Basis of PI space $\bV^\bl$]\label{prop:basis-pi2}
    Let $\bl\in\calL^\Nat$. A basis of $\bV^{\bl}$ consists of the functions
\[
    \left\{ \bphi{\bl}{\bbm},\; \bbm\in\oMl\right\}.
\]
In particular, the dimension of $\bV^\bl$ is $\card{\oMl}$.
\end{proposition}

We defer the proof of this proposition to Section~\ref{sec:proof_equivariant_lie_algebra}. Note that Definition~\eqref{eq:sumphi} can naturally be extended to an arbitrary number of variables by defining $\bphi{l}{m} : \bR \in \textup{MS}(\Omega) \mapsto \sum_{\br \in \bR } \phi^l_m(\br)$. 
We can therefore also account for multiset functions within the same framework, although Proposition~\ref{prop:basis-pi2} may no longer hold in the multiset setting.

\begin{remark}[Group action]
\label{rem:group_action}
Note that $\bV^\bl$ is stable with respect to the group action, and the latter acts in the same way over the functions $\phi^\bl_\bbm$, 
 $\bpsi{\bl}{\bbm}$ and $\bphi{\bl}{\bbm}$. Consider $\rep{\bl}{}=\displaystyle\bigotimes_{i=1}^\Nat\rep{l_i}{}$, which is a representation as tensor product of representations. Since
    \[
        \bphi{\bl}{\bbm}(\bR)=\sum_{1\le n_1,\ldots,n_\Nat\le\Nat}\phi^\bl_\bbm(\br_{n_1},\ldots,\br_{n_\Nat}),
    \]
    we have that for any $g\in G$
    \begin{align*}
          \bphi{\bl}{\bbm}(g \cdot \bR)&=\sum_{1\le  n_1,\ldots,n_\Nat\le\Nat}\phi^\bl_\bbm(g \cdot (\br_{n_1},\ldots,\br_{n_\Nat}))\\
          &= \sum_{\bbm'\in \calM_\bl}\rep{\bl}{\bbm,\bbm'}(g)\sum_{1\le n_1,\ldots,n_\Nat\le\Nat}\phi^\bl_{\bbm'}(\br_{n_1},\ldots,\br_{n_\Nat})\\
          &= \sum_{\bbm'\in \calM_\bl}\rep{\bl}{\bbm,\bbm'}(g)\bphi{\bl}{\bbm'}(\bR).
    \end{align*}
It can be easily checked that the same property holds for $\phi^\bl_\bbm$ and $\bpsi{\bl}{\bbm}$.
\end{remark}

{\begin{definition}[Scalar product on $\bV^\bl$]
    As $\left\{  \bphi{\bl}{\bbm},\; \bbm \in \oMl \right\}$ is a basis of $\bV^{\bl}$, we can define an inner product on $\bV^{\bl}$ by the relation
    \[
        \langle\bphi{\bl}{\bbm},\bphi{\bl}{\bbm'}\rangle=\delta_{\bbm,\bbm'}.
    \]
\end{definition}
We now turn to characterizing GE-PI functions with components in $\bV^\bl$, starting by introducing some notation.
For $\bl$'s containing a single value $l$, i.e., 
    $\bl = (l,l,\ldots,l)\in\calL^N,$
the classes defined in~\eqref{eq:GE_class} are in one-to-one correspondence with ordered $\bbm\in\calM_\bl$, which we use by default as a representative for the corresponding class. 
Furthermore, the classes $\bbm\in\oMl$ are also in one-to-one correspondence with their Parikh vectors or count vectors 
$\blamb{}{\bbm}$ (see e.g.~\cite[Supplementary lecture H]{Kozen2007-ot}), which are defined as the vectors containing the number of occurrences of the elements in $\calM_l$. 
We give an explicit example to better illustrate this concept.
\begin{example}
    Assume $l=1$, $\calM_l = \{-1,0,1\}$ and $\bl=(l,l,l,l,l)$, then the class $\bbm=\overline{(-1,-1,0,1,1)}\in\oMl$ has a count vector 
    $\blamb{}{\bbm}=(2,1,2)$, whose components stand for the two occurrences of $-1\in\calM_l$, one occurrence of $0\in\calM_l$ and two occurrences of $1\in\calM_l$, respectively. 
\end{example}

Note that the sum of the elements in the count vector is equal to $\Nat$, the length of $\bl$, and we can identify $\oMl$ with the set
\begin{equation}
\label{eq:lambda}
    \Lambda_{\bl} := 
    \{\blamb{}{}\in \N^{\calM_l}:\sum_{i\in\calM_l}\blamb{}{i} = \Nat\}.
\end{equation}

For a general $\bl$ possibly containing different values, we assume that it is ordered for notation simplicity, but without loss of generality. The index $\bl$ can then be divided into a set of minimal length $\Nb$ of sub-indices $\{\bl^{(j)}\}_{j=1}^\Nb$ of length $|\bl^{(j)}|$, so that all elements in $\bl^{(j)}$ are identical, that we call minimal partition as defined below. Specifically, we say that $\bl^{(j)}$ and $\bl^{(j')}$ do not intersect if $\bl^{(j)}_{p}\neq\bl^{(j')}_{q}$ for all possible $p,q$, and we write $\bl^{(j)} \cap \bl^{(j')} = \varnothing$ in this case.

\begin{definition}[Minimal partition]
    Let $\bl\in\calL^\Nat$ be ordered, we say that $(\bl^{(j)})_{j=1}^{\Nb}$ is the minimal partition of $\bl$ if each block $\bl^{(j)}$ has a single repeated value and that for all $j,j'\in\{1,\ldots,\Nb\},j\neq j', \; \bl^{(j)}\cap\bl^{(j')}=\varnothing$.
\end{definition}

Then it immediately follows that 
\[
     \calS_{\bl} = \bigotimes_{j=1}^\Nb S_{|\bl^{(j)}|},
\quad 
\text{and}
\quad
     \oMl = \bigotimes_{j=1}^\Nb \overline{\calM}_{\bl^{(j)}}.
\]
For $\obm=(\obm^{(1)}, \ldots, \obm^{(\Nb)})\in\oMl$, we use $\bbm = (\bbm^{(1)}, \ldots, \bbm^{(\Nb)})$ as its representative, where $\bbm^{(j)}$ are ordered, for all $j\in\{1,2,\ldots,\Nb\}$. 
Thus, for a general ordered $\bl$, $\oMl$ is in one-to-one correspondence with 
$\displaystyle \bigotimes_{j=1}^\Nb \Lambda_{\bl^{(j)}}$, 
with $\Lambda_{\bl^{(j)}}$ defined in~\eqref{eq:lambda}. For simplicity, we also denote this set by~$\Lambda_{\bl}$, which will not cause confusion since when $\bl$ contains only identical values, this definition is consistent with~\eqref{eq:lambda}.
For any $\bbm\in\oMl$, we denote its corresponding element in $\Lambda_{\bl}$ by $\blamb{}{\bbm}=(\blamb{}{\bbm^{(1)}},\blamb{}{\bbm^{(2)}},\ldots,\blamb{}{\bbm^{(\Nb)}})$. Similarly, for any $\blamb{}{} \in \Lambda_\bl$, we denote the associated element in $\oMl$ by $\bbm_{\blamb{}{}}$. 

We now define the set of interacting classes as follows, which will be key in constructing a GE-PI basis.

\begin{definition}[Interacting classes]
    For a given $\bbm\in\oMl$, we define the set of its interacting classes $\calN_\bbm\subset\oMl$ as 
    \begin{equation}
        \label{eq:interacting_classes}
         \begin{split}
     \calN_\bbm := 
     \Bigg\{\bbm^{}_{(j,p,q)}:= &\;\;\bbm_{\blamb{}{(j,p,q)}},
     \quad 
     j\in\{1,2,\ldots,\Nb\}, \;\; p,q\in\calM_{\bl^{(j)}_1},\;\; p\ne q, \\
     &
     \textup{where } \quad 
     \blamb{}{(j,p,q)}
     =  (\blamb{}{\bbm^{(1)}},\blamb{}{\bbm^{(2)}},\ldots,\blamb{}{\bbm^{(j-1)}},\blamb{}{\bbm^{(j)}} - {\bf e}_p + {\bf e}_q,\blamb{}{\bbm^{(j+1)}},\ldots,\blamb{}{\bbm^{(\Nb)}})\\
     & \textup{if $\blamb{}{\bbm^{(j)},p}> 0$, otherwise   $ \blamb{}{(j,p,q)}$ is discarded }
     \Bigg\},
 \end{split}
    \end{equation}
where ${\bf e}_k$ denotes the $k$-th vector of the canonical basis of $\R^{\calM_{\bl^{(j)}_1}}$, $\bl^{(j)}_1$ is the first component of $\bl^{(j)}$ and $\blamb{}{\bbm^{(j)},p}$ the $p$-th component of $\blamb{}{\bbm^{(j)}}$. 
\end{definition}
The following theorem provides a basis for GE-PI functions for the representation $\rep{L}{}$ in the space $\left[ \bV^{\bl} \right]^{\card{\calM_L}}$, whose span we denote by $\bV^{\bl,L}$, by analogy with the GE case.


\begin{theorem}[Basis of the GE-PI space $\bV^{\bl,L}$]
\label{thm:GEPI_basis}
Let $\bl\in\calL^\Nat$, $L\in\calL$, and let 
\begin{equation}\label{eq:matM-PI}
         M^{\bl,L} = \begin{pmatrix}
         M_1 \\
         M_2 \\
         \vdots \\
         M_{N_{\rm{dim}}}
    \end{pmatrix},
    \end{equation}
    where for $\dd = 1,\ldots, N_{\rm{dim}}$,
    the elements of the matrices $M_\dd$ read for $(\bbm,k), (\bbm',k') \in \oMl\times \calM_L$ as
    \begin{equation}
    M_\dd[(\bbm,k),(\bbm',k')] = 
        \begin{cases}
        \displaystyle\sum_{j=1}^\Nat
        \vrho{\bl^{(j)}_1,\dd}{m_j,m_j} - \vrho{L,\dd}{k,k}, \qquad &\textup{if}\;\; \bbm=\bbm',  \; \; k=k', \\
         -\vrho{L,\dd}{k,k'}, \qquad &\textup{if}\;\; 
         \bbm = \bbm',  \; \;  k\ne k', \\
         \blamb{}{{\bbm'}^{(j)},q}
         \;\vrho{\bl^{(j)}_1,\dd}{q,p},  \qquad &\textup{if}\;\; 
         \bbm'=\bbm_{(j,p,q)}\in\calN_{\bbm}, \;\;  k = k',\\
        0, \qquad &\textup{otherwise,} \\
    \end{cases}
    \label{eq:M_GE_PI}
    \end{equation}
    with $\calN_\bbm$ defined in~\eqref{eq:interacting_classes}.
    Then a basis of $\bV^{\bl,L}$
    is given by 
    \[
    b^{\bl,L}_i = 
    \left[ 
    \sum_{\bbm \in \oMl}c_{\bbm,k,i}^{\bl,L}\bphi{\bl}{\bbm} \right]_{k\in \calM_L},\;\; i=1,\ldots,\dim\big(\ker(M^{\bl,L})\big),
    \]
     where 
     $\bc^{\bl,L}_i = \{c_{\bbm,k,i}^{\bl,L}\}_{\bbm\in\oMl,k\in\calM_L},\; i=1,\ldots,\dim\big(\ker(M^{\bl,L})\big)$
     is a basis of $\ker(M^{\bl,L})$. 
    Thus, 
    \[
        \dimpi{\bl}{L}=\dim\big(\ker(M^{\bl,L})\big).
    \]
\end{theorem}

The proof of this theorem can be found in Section~\ref{sec:proof_equivariant_lie_algebra}. 
As with the GE case, the superscript of the matrix $M^{\bl,L}$ is sometimes omitted when the associated $\bl$ and $L$ are clear. Compared to Proposition~\ref{prop:GE_basis} that deals with the GE case, the only two differences are: (1) the set $\calM_\bl$ being replaced by the set of classes $\oMl$, an even smaller set compared with $\calM_\bl$; and (2) the presence of the
 $
    \blamb{}{{\bbm'}^{(j)},q} \;\;
         \vrho{\bl^{(j)}_1,\dd}{q,p},  $
         instead of $\vrho{l_j,\dd}{m'_j,m_j}$.
Therefore, our construction of the GE-PI basis is almost as efficient as that of the GE basis. To the best of our knowledge, this has not been the case in any existing work, where the presence of permutation invariance often incurs substantial additional costs~\cite{Batatia2023-us,Dusson2022-ie,geiger2022e3nn}. 


We now provide a similar result when the basis $\bpsi{\bl}{\bbm}$ is used.
\begin{corollary}
    Let $\bl\in\calL^\Nat$, $L\in\calL$.
    Let $\bc^{\bl,L}_i = \{c_{\bbm,k,i}^{\bl,L}\}_{\bbm\in\oMl,k\in\calM_L}, i=1,\ldots, \dim(\ker(M^{\bl,L}))$ be a basis of $\ker(M^{\bl,L})$. 
    A basis of GE-PI
    functions for the representation $\rep{L}{}$ in the space $\left[V^\bl\right]^{\card{\calM_L}}$ is
    given by
    \[
    b^{\bl,L}_i = 
    \left[ 
    \sum_{\bbm \in \oMl}c_{\bbm,k,i}^{\bl,L} \bpsi{\bl}{\bbm}\right]_{k\in \calM_L},\;\; i=1,\ldots,\dim\big(\ker(M^{\bl,L})\big).
    \]
\end{corollary}

\begin{proof}
    Since the group action on $G$ acts the same over the functions $\bpsi{\bl}{\bbm}$ and $\bphi{\bl}{\bbm}$, Proposition~\ref{prop:basis-pi} combined with an adapted proof of Theorem~\ref{thm:GEPI_basis} easily provides the result.
\end{proof}



\subsection{Recursive construction of GE and GE-PI bases}


In this section, we show how the GE and GE-PI bases can be constructed recursively. This also offers a way to obtain an exact recursive formula for the dimensionality of the considered GE and GE-PI spaces. 
Indeed, a priori it is not so clear whether the tensor product of two GE sub-bases is a basis of the final space or an overcomplete set of functions. To obtain these results, we need a few additional assumptions on the considered group and the representations. We assume that the group $G$ is compact, so that the Haar measure is well-defined~\cite[Chapter 4]{Hall2013-af}, and we consider a complete set of mutually inequivalent irreducible representations $\rep{l}{}$ for $l\in \calL$ with corresponding spaces $V^l$.
Under these assumptions, the spaces $V^\bl$ and $\bV^\bl$ for $\bl\in\calL^N$ can respectively be decomposed into the direct sum of GE and GE-PI spaces.


\begin{proposition}[Direct sum of spaces with GE functions]\label{prop:direct_sum}
    Suppose $G$ is a 
    compact Lie group,
    $\bl 
    \in \calL^\Nat$ and 
    $V^\bl$ is stable by group action. 
    There holds 
    \begin{equation}
    \label{eq:direct_sum_GE}
        V^{\bl} = \bigoplus_{L \in \calL}  
        \;
        \bigoplus_{k\in \calM_L} [V^{\bl,L}]_k, 
    \end{equation}
    where $[V^{\bl,L}]_k$ denotes the space spanned by the $k$-th component of the functions in $V^{\bl,L}$.
\end{proposition}

\begin{proof}
    Since the space $V^\bl$ is stable by group action, the group action defines a representation $\rep{\bl}{}$, which can be blockdiagonalized with irreducible representations $\rep{L}{}$, $L\in\calL$ and corresponding basis functions $v^L_{i,k}$, $i \in \{1, \ldots, d_L\}$ and $k\in \calM_L$, where $d_L$ is the number of copies of $\rep{L}{}$ in the blockdiagonalization.
    Noting that 
    $V^{\bl,L} = {\rm Span} \{ [v^L_{i,k}]_{k\in \calM_L}, i\in \{1, \ldots, d_L\}  \}$, the result follows.
\end{proof}


\begin{proposition}[Direct sum of spaces with GE-PI functions]
    Suppose $G$ is a compact Lie group,
    $\bl
    \in \calL^\Nat$ and 
    $\bV^\bl$ is stable by group action.
    There holds
    \begin{equation}
    \label{eq:direct_sum_GEPI}
         \bV^{\bl} = \bigoplus_{L \in \calL}  
        \;
        \bigoplus_{k\in \calM_L} [\bV^{\bl,L}]_k,
    \end{equation}
    where $[\bV^{\bl,L}]_k$ denotes the space spanned by the $k$-th component of the functions in $\bV^{\bl,L}$.
\end{proposition}

\begin{proof}
    Since the group action is identical on $V^l$ and $\bV^l$ for $l\in\calL$ using Remark~\ref{rem:group_action}, a similar proof as Proposition~\ref{prop:direct_sum} applies in this case.
\end{proof}


We now provide a recursive construction of GE bases and a recursive formula for the dimensionality of the GE spaces.

\begin{proposition}[GE basis based on recursion]
\label{prop:GE_rec}
Let $N,N_1,N_2\in \N$ with $N_1+N_2=N$.
Let $\bl\in\calL^\Nat$  and $L\in\calL$. Suppose $\bl = (\bl^{(1)}, \bl^{(2)})\in\calL^{\Nat_1}\times\calL^{\Nat_2}$, then a basis for $V^{\bl,L}$ is given by
\[
    b^{\bl,L}_{(L_1,i_1,L_2,i_2,j)} = \left[ \sum_{\substack{k_1\in\calM_{L_1} \\  k_2\in\calM_{L_2}}} 
    c_{(k_1,k_2), k,j}^{(L_1,L_2),L}
    [b^{\bl^{(1)},L_1}_{i_1}]_{k_1} [b^{\bl^{(2)},L_2}_{i_2}]_{k_2} \right]_{k\in \calM_L},\quad 
    \begin{array}{l}
         L_1,L_2\in\calL,   \\
         i_1\in\{1,\ldots, \dimge{\bl^{(1)}}{L_1}\}, \\
         i_2\in\{1,\ldots, \dimge{\bl^{(2)}}{L_2}\},\\
         j \in \{1,\ldots,\dimge{(L_1,L_2)}{L}\},
    \end{array}
\]
where $\{b^{\bl^{(1)},L_1}_{i_1}\}_{i_1 = 1,\ldots, \dimge{\bl^{(1)}}{L_1}}$, 
$\{b^{\bl^{(2)},L_2}_{i_2}\}_{i_2 = 1,\ldots, \dimge{\bl^{(2)}}{L_2}}$ are bases of $V^{\bl^{(1)},L_1}$ and $V^{\bl^{(2)},L_2}$, respectively, and
$\{c_{(k_1,k_2), k,j}^{(L_1,L_2),L}\}_{j=1,\ldots,\dimge{(L_1,L_2)}{L}}$ is a basis of $\ker(M^{(L_1,L_2),L})$ with $M^{(L_1,L_2),L}$ being defined in~\eqref{eq:matM}. 
Moreover, there holds
\begin{equation}\label{eq:dim_GE_rec}
    \dimge{\bl}{L} = \sum_{L_1,L_2\in \calL} \dimge{(L_1,L_2)}{L} \;\dimge{\bl^{(1)}}{L_1}\;\dimge{\bl^{(2)}}{L_2}.
\end{equation}
\end{proposition}

The proof can be found in Section~\ref{sec:proof_equivariant_lie_algebra}.
Note that,
as mentioned in Remark~\ref{rem:CG}, 
the coupling coefficients $c_{(k_1,k_2), k,j}^{(L_1,L_2),L}$ are called Clebsch--Gordan coefficients. 

We then provide a result that can be used to compute a GE-PI basis recursively. 

\begin{proposition}[GE-PI basis based on recursion]
\label{prop:GEPI_rec}
Let $N,N_1,N_2\in \N$ with $N_1+N_2=N$.
Let $\bl\in\calL^\Nat$ and $L\in\calL$ . Suppose $\bl = (\bl^{(1)}, \bl^{(2)})\in\calL^{\Nat_1}\times\calL^{\Nat_2}$, then a spanning set for $\bV^{\bl,L}$ is 
\[
        b^{\bl,L}_{(L_1,i_1,L_2,i_2,j)} = \left[ \sum_{\substack{k_1\in\calM_{L_1} \\  k_2\in\calM_{L_2}}}  
        c_{(k_1,k_2), k,j}^{(L_1,L_2),L}
        [b^{\bl^{(1)},L_1}_{i_1}]_{k_1} [b^{\bl^{(2)},L_2}_{i_2}]_{k_2} \right]_{k\in \calM_L}, \quad
        \begin{array}{l}
         L_1,L_2\in\calL,   \\
         i_1\in\{1,\ldots, \dimpi{\bl^{(1)}}{L_1}\}, \\
         i_2\in\{1,\ldots, \dimpi{\bl^{(2)}}{L_2}\},\\
         j \in \{1,\ldots,\dimge{(L_1,L_2)}{L}\},
    \end{array}
\]
where $\{b^{\bl^{(1)},L_1}_{i_1}\}_{i_1 = 1,\ldots, \dimpi{\bl^{(1)}}{L_1}}$ and 
$\{b^{\bl^{(2)},L_2}_{i_2}\}_{i_2 = 1,\ldots, \dimpi{\bl^{(2)}}{L_2}}$ are bases of $\bV^{\bl^{(1)},L_1}$ and $\bV^{\bl^{(2)},L_2}$
respectively, and
$\{c_{(k_1,k_2), k,j}^{(L_1,L_2),L}\}_{j=1,\ldots,\dimpi{(L_1,L_2)}{L}}$ is a basis of $\ker(M^{(L_1,L_2),L})$ with $M^{(L_1,L_2),L}$ being defined in~\eqref{eq:matM}. 

Moreover, if $\bl^{(1)} \cap \bl^{(2)} = \varnothing$, this family is a basis of $\bV^{\bl,L}$ and there holds
\begin{equation*}
    \dimpi{\bl}{L} = \sum_{L_1,L_2\in\calL}
    \dimge{(L_1,L_2)}{L}
    \;
    \dimpi{\bl^{(1)}}{L_1} \;
    \dimpi{\bl^{(2)}}{L_2}.
\end{equation*}
\end{proposition}
\begin{proof}
    The proof for the non-intersecting case is similar to the proof of Proposition~\ref{prop:GE_rec} 
    noting that if $\bl^{(1)} \cap \bl^{(2)} = \varnothing$,
     then $\bV^\bl = \bV^{\bl^{(1)}} \otimes \bV^{\bl^{(2)}}$, for which a basis is given by
     \[
    ([b^{\bl^{(1)},L_1}_{i_1}]_{k_1} [b^{\bl^{(2)},L_2}_{i_2}]_{k_2})_{L_1,L_2 \in\calL,\; i_1 \in \{ 1, \ldots, \dimpi{\bl^{(1)}}{L_1}\},\; i_2 \in \{1, \ldots, \dimpi{\bl^{(2)}}{L_2}\} },
    \]
    and finally using~\eqref{eq:direct_sum_GEPI} instead of~\eqref{eq:direct_sum_GE}. In that case, the dimensionality is a direct consequence.
    If $\bl^{(1)} \cap \bl^{(2)} \neq \varnothing$, we have $\bV^\bl = \bV^{\bl^{(1)}} \otimes \bV^{\bl^{(2)}}$ but the functions \[
    ([b^{L_1}_{\bl^{(1)} i_1}]_{k_1} [b^{L_2}_{\bl^{(2)} i_2}]_{k_2})_{L_1,L_2 \in \N, i_1 \in \{ 1, \ldots, \dimpi{\bl^{(1)}}{L_1}\}, i_2 \in \{1, \ldots, \dimpi{\bl^{(2)}}{L_2}\} } 
    \]
    are only a spanning set and not a basis of  $\bV^\bl$.
    Therefore, following the steps of the non-intersecting case, we only generate a spanning set of GE-PI functions.
\end{proof}

Note that in the case of intersecting $\bl^{(1)},\bl^{(2)},$ we can numerically obtain a basis by, for example, performing a singular value decomposition on the spanning set given in Proposition~\ref{prop:GEPI_rec}.


\medskip

In practice, it is not always clear which construction is more efficient between the direct and recursive constructions. However, with the result on GE-PI functions for $\bl\in\calL^N$ with minimal partition $\bl=(\bl^{(1)},\bl^{(2)},\ldots,\bl^{(\Nb)})$, we expect it to be efficient to first compute basis functions for the different blocks $\bl^{(j)}$ of identical values, and then assemble them using Proposition~\ref{prop:GEPI_rec}, especially for large $N$, see Section~\ref{sec:num-rec} for an elementary numerical comparison.



\section{Application to some rotation groups}\label{sec:appl}

We now apply our construction to the groups  $SU(2)$ and $SO(3)$, which naturally appear in quantum mechanics, molecular mechanics, and computer graphics. Specifically, the generators of the underlying Lie algebra are taken as the infinitesimal rotations under the ZYZ convention. We show how the matrices~\eqref{eq:matM} and~\eqref{eq:matM-PI} 
simplify for these specific groups, and investigate the algebraic properties of the corresponding GE and GE-PI bases. In particular, we provide their exact dimensionality and prove asymptotic estimates.


\subsection{Representations and Lie algebra}
We first introduce the explicit expression of the Wigner-D matrices with respect to Euler angles, which are irreducible representations of $SU(2)$ and $SO(3)$~\cite{wigner1931gruppentheorie}. As it is standard, we denote the Wigner-D matrices by $D^\ell$ instead of $\rep{\ell}{}$.

\begin{definition}[Wigner-D matrix]\label{def:wigner-D}
Let $(\alpha, \beta, \gamma)\in
[0,4\pi)\times[0,\pi)\times[0,4\pi)$ for $SU(2)$
or 
$[0,2\pi)\times[0,\pi)\times[0,2\pi)$
for $SO(3)$
be the three Euler angles that parametrize an element $Q$ in either the group $SU(2)$ or $SO(3)$ under the ZYZ convention. Then for $\ell = 0,\; \frac12, \;1, \;\frac32, \;\ldots$, the Wigner-D matrix $D^\ell$ is defined as a $(2\ell+1)\times(2\ell+1)$ matrix, whose elements read
\begin{equation}
      D^\ell_{\mu m}(Q) = D^\ell_{\mu m}(\alpha, \beta, \gamma) = e^{-\mathrm{i}m\alpha}d^\ell_{\mu m}(\beta) e^{-\mathrm{i}\mu\gamma},
      \label{eq:wigner-D}
\end{equation}
for $\mu,m \in \{-\ell, -\ell+1, \ldots, \ell-1, \ell\}$, where
\[
    d^\ell_{\mu m}(\beta) = [(\ell+m)!(\ell-m)!(\ell+\mu)!(\ell-\mu)!]^\frac{1}{2}\sum_{s = \max(0, \mu-m)}^{\min(\ell+\mu,\ell-m)} \frac{(-1)^{s}(\cos \frac{\beta}{2})^{2l+\mu-m-2s}(\sin \frac{\beta}{2})^{m-\mu+2s}}{(\ell+\mu-s)!s!(m-\mu+s)!(\ell-m+s)!}.
\]
\end{definition}
In particular, the irreducible representations of $SU(2)$ are the Wigner-D matrices for 
$\ell \in \calL_{SU(2)}$ defined~as
\begin{equation}
    \label{eq:calL}
     \calL_{SU(2)} :=  \left\{\frac{\ell}{2}: \ell\in\N_0 \right\},
\end{equation}
while the irreducible representations of $SO(3)$ correspond to Wigner-D matrices for 
$\ell \in \calL_{SO(3)}$
defined as 
\begin{equation}
\label{eq:calL2}
     \calL_{SO(3)} := \N_0.
\end{equation}
To unify the notation and emphasize the dependency of the index sets for the irreducible representations on the group $G$, we denote the sets as $\calLG$, particularly for $G=SU(2)$ and $G=SO(3)$, in this section. To apply Proposition~\ref{prop:GE_basis} and Theorem~\ref{thm:GEPI_basis}, we need the partial derivatives of the Wigner-D matrices with respect to the Euler angles at the origin, which corresponds to the neutral element of the group (identity denoted by~$I$). For notational convenience, we define for $\ell\in\calLG$ the set
\begin{equation}
    \label{eq:calMU2}
        \calM_\ell = \{-\ell,-\ell+1,\ldots,\ell-1,\ell\}, 
        \quad \text{ with }
        \quad 
         \card{\calM_\ell} = 2\ell+1.
\end{equation}
Then by~\eqref{eq:rep-deriv}, the derivatives of the representations $D^\ell$ evaluated at the chosen generators are given element-wise by the following Proposition~\ref{prop:deri_alpha_gamma}. 

\begin{proposition}[Derivatives of Wigner-D matrices]
\label{prop:deri_alpha_gamma}
Let $\ell\in\calLG$, $\mu,m \in \calM_\ell$, then
\begin{align}
\label{eq:deri_alpha}
    \vrho{\ell,1}{\mu,m} &:=  \frac{\partial D^\ell_{\mu m}(I)}{\partial \alpha} = - \mathrm{i} m \delta_{\mu m}, \\
\label{eq:deri_beta}
     \vrho{\ell,2}{\mu,m} &:= \frac{\partial D^\ell_{\mu m}(I)}{\partial \beta} = 
    \begin{cases}
        \frac{1}{2}[(\ell-m+1)(\ell+m)]^{\frac{1}{2}},  & \quad \textup{if } m = \mu + 1\\
        -\frac{1}{2}[(\ell+m+1)(\ell-m)]^{\frac{1}{2}},  & \quad \textup{if } m = \mu - 1\\
        0, & \quad \mbox{otherwise,} \\ 
    \end{cases} 
    \\
        \vrho{\ell,3}{\mu,m} &:=  
        \frac{\partial D^\ell_{\mu m}(I)}{\partial \gamma} = - \mathrm{i} m \delta_{\mu m}.
    \label{eq:deri_gamma}
\end{align}
\end{proposition}


While this proposition is a standard result, we were not able to find a complete reference for a proof of it in the literature. Whereas some ideas can be found in~\cite{rose1995elementary}, we provide a proof in Appendix~\ref{app:Wigner-Dmatricesproof} for completeness.


\subsection{One-variable compatible bases}
\label{sec:compatible_bases}

We now introduce some possible one-variable compatible bases for $SU(2)$ and $SO(3)$. Although a large freedom exists in the choice of the bases that are stable by group action, we focus on compatible bases with respect to the irreducible representation $D^\ell$ for $\ell\in\calLG$. Specifically, we focus on the bases that are frequently used in real implementations. 
For each case, we need to define the considered spatial domain $\Omega$, the chosen $\F$ and $\K$, the considered group action, and the one-variable compatible bases.

\paragraph{Wigner-D matrices}

First, for all $\ell\in\calLG$, the elements in each column of $D^\ell$ are functions from $\Omega := G$ to $\F := \C$, and form a compatible basis with respect to the representation $D^\ell$ itself. This corresponds to $\K := \C$.
Precisely, let $l=\ell$, $\calM_l=\calM_\ell$, then for any index $\mu\in\calM_\ell$, the space
\[
    \displaystyle V^l_\mu = \textup{Span} 
    \big\{D^\ell_{m\mu}:G\to\C, \; m\in\calM_l \big\}
\]
is stable under the group action defined by matrix multiplication as
\[
    \forall Q, \; \br \in G, \quad 
    Q \cdot \br := Q \br,
\]
since
\[
    \forall Q \in G, \; \br \in G,  \qquad D^\ell_{m \mu }(Q\cdot \br) = [D^\ell(Q)\cdot D^\ell(\br)]_{m \mu}=\sum_{m'\in\calM_\ell}D^\ell_{mm'}(Q)D^\ell_{m' \mu}(\br).
\]
Hence, columns of Wigner-D matrices $\{D^\ell_{m\mu}:G\to\C\}_{m\in\calM_\ell}$ form compatible bases for the representations $D^\ell$. 

\paragraph{Spherical harmonics}

In particular, for $\ell\in\N_0$, $\mu=0\in\calM_l$, and for $Q\in SO(3)$ represented by Euler angles $(\alpha,\beta,\gamma)$, there holds
\[
    D^\ell_{m0}(Q) = D^\ell_{m0}(\alpha,\beta,\gamma) = e^{-\textup{i}m\alpha}\cdot (-1)^m\sqrt{\frac{(l-m)!}{(l+m)!}}P_\ell^m(\cos\beta),
\]
where $P_\ell^m$ is the associated Legendre polynomial. Clearly, $D^\ell_{m0}$ is independent of the angle $\gamma$, and is proportional to the complex spherical harmonics of order $\ell$ and angular momentum quantum number $m$: $Y_\ell^m(\alpha,\beta)$.
In other words, the spherical harmonics $\{Y_\ell^m\}_{m\in \calM_\ell}$ form a compatible basis corresponding to $D^\ell$ for $\ell\in\N_0$, the irreducible representations of $SO(3)$. Thus, defining $\Omega := S^2$ the 2-dimensional unit sphere, $\F := \C$, $\K:= \C$, and the group action as matrix-vector multiplication
\[
    \forall Q \in SO(3), \; \br \in S^2, \quad 
    Q\cdot \br := Q\br,
\]
the space for $l:=\ell\in \N_0$ defined by
\[
    V^l := \textup{Span} \big\{Y_l^m:S^2\to\C, \; 
    m\in\calM_l \big\}
\]
is stable under group action. Note that we could also take $\F := \R$, and $\K :=\R$ by considering real spherical harmonics, although another set of irreducible representations would have to be taken into account in this case.

\paragraph{One-variable basis on $\R^3$}

In practice, it is of particular interest to investigate a compatible basis on $\R^3$. This is often done by padding a 1D radial basis to spherical harmonics. 
Precisely speaking, we take the spatial domain $\Omega$ as $\R^3\backslash \{\mathbf{0}\}$ in this case, 
and let $\F := \C$, and $\K :=\C$. Then for $\br=|\br|\hat{\br}\in\R^3\backslash\{\mathbf{0}\}$ with $|\br|> 0$ and 
$\hat{\br} \in S^2$, we define the group action as
\[
    \forall Q\in SO(3), \;  \br \in \R^3\backslash\{\mathbf{0}\}, 
    \quad 
    Q \cdot \br := |\br| (Q \hat{\br}).
\]
Taking $\{P_n\}_{n\in\N_0}\subset L^2(\R)$, an orthogonal basis of $L^2(\R)$, we
define for $l=(n,\ell)\in\N_0\times\N_0=:\calL$ the set $\calM_l := \calM_\ell$, and for $m\in \calM_l$ 
\[
    \phi^l_m=\phi^{(n,\ell)}_m : \br \in \R^3\backslash\{\mathbf{0}\}  \mapsto P_n(|\br|)Y_\ell^m(\hat{\br}).
\]
Then, $\big\{\phi^l_m \big\}_{m\in\calM_l}$
forms a compatible basis of $SO(3)$ corresponding to representation $D^\ell$, as the radial function $P_n$ remain invariant under $SO(3)$ rotations. The space
\[
    V^{l} = \textup{Span}\big\{ \phi^l_m, \; 
    m\in\calM_l \big\}
\]
is therefore stable under the defined group action. 

\paragraph{One-variable basis on higher-dimensional spaces}

In fact, any invariant transformation with respect to rotations 
has no impact on the stability of the corresponding basis. One can take into account higher-dimensional spaces $\Omega = W \times S^2$, with $W$ independent of rotations. 
For example, a particle in $\R^3$ is often described not only by its position, but also by other properties, such as its chemical species. In this case, one could, e.g., take 
\[
    W = \N\times\big(\R\backslash\{0\}\big), \quad \Omega = \N \times \big(\R^3\backslash\{\mathbf{0}\}\big), \quad \F = \C, \quad \K = \C,
\]
with the group action being defined as 
\[
    \forall Q \in SO(3),\; (z,\br) \in \Omega, 
    \quad Q \cdot (z,\br) := 
    (z,Q\br).
\]
Hence, we may write for $l=(n_1,n_2,\ell)\in \N \times \N_0 \times \N_0=:\calL$ the basis index set $\calM_l:=\calM_\ell$, and the compatible basis 
\[
    \{\phi^l_m: (z,\br)
    \in \N \times \big(\R^3\backslash\{\mathbf{0}\}\big) \mapsto F_{n_1}(z)P_{n_2}(|\br|)Y_\ell^m(\hat\br)\}_{m\in\calM_l},
\]
where $F_{n_1}$ are basis functions depending on the chemical species (and possibly other parameters independent of rotations),
which are usually chosen as $\delta$-functions (also called one-hot embedding in some contexts),
$P_{n_2}$ are radial functions as stated above. 
Thus, the
corresponding function space
\[
    V^l = \textup{Span}
    \big\{\phi^l_m, \; m\in\calM_l \big\}
\]
is stable for the representation $D^\ell$. 

\medskip

To summarize, we present in Table~\ref{tab:comp_bases} the aforementioned compatible bases in a unified form.
Note that $\ell$ is always a component of $l$, and the basis index set for $l\in\calL$ satisfies $\calM_l = \calM_\ell$.
In some cases that should be clear from the context, when it comes to operations on $l\in\calL$, we refer to the corresponding operations on the $\ell$-components.


\begin{remark}
    Note that in general the vector space $\F$ and the field $\K$ do not have to be the same. For example, considering functions with two components, such as one spin up and one spin down, we would have $\F$ = $\C^2$ and $\K=\C$.
\end{remark}

\begin{table}[htb!]
\centering
\caption{Summary of one-variable compatible bases for the irreducible representations of $SU(2)$ and $SO(3)$}
\label{tab:comp_bases}
\begin{tabular}{|m{3cm}|m{3.4cm}|m{2cm}|m{2.5cm}|c|}
    \hline
    \centering Compatible basis functions $\phi^l_m$ & \centering {Applicable groups $G$} & \centering{Domain $\Omega$} & \centering{Index set $\calL$} & 
    {$\;\;\; \displaystyle\bigoplus_{l\in\calL} V^l$}
    \\
    \hline
    \centering $D^l_{m\mu}$ & \centering  $SU(2)$, $SO(3)$ & \centering   $G$   &  \centering    $\calLG$ \eqref{eq:calL}, \eqref{eq:calL2}    & -         \\
    \hline
    \centering $Y_l^m$    &  \centering    $SO(3)$      &  \centering  $S^2$ & \centering    $\N_0$       & $L^2(S^2)$ \\
    \hline
    \centering $P_nY_l^m$ & \centering     $SO(3)$      &  \centering $\R^3\backslash\{\mathbf{0}\}$ & \centering $\N_0\times\N_0$ & $L^2(\R^3\backslash\{\mathbf{0}\})$ \\
    \hline
    \centering $F_{n_1} P_{n_2} Y_l^m$ & \centering     $SO(3)$      & \centering  $\N\times\big(\R^3\backslash\{\mathbf{0}\}\big)$ & \centering $\N\times\N_0\times\N_0$ & $L^2(\Omega)$ \\
    \hline
\end{tabular}
\end{table}


In the following, given compatible bases $\{\phi^l_m\}_{m\in\calM_\ell}$ for $l\in \calL$ depending on the chosen group and representations, such as those summarized in Table~\ref{tab:comp_bases} and the corresponding function spaces $V^l$, our goal is to practically find a basis for the GE space $V^{\bl,L}$ or the GE-PI space $\bV^{\bl,L}$, for given $\bl\in\calL^N$ and $L\in\calL_G$, using the strategy described in Section~\ref{sec:equivariant_lie_algebra}. We present the details in the coming subsections.


\subsection{Construction of multi-variable GE bases}

We first construct the matrix~\eqref{eq:matM} for the groups $SU(2)$ and $SO(3)$.
Instead of considering the whole matrix $\MlL$, we deal with each matrix $M_\dd,\; d=1,2,3$ independently, with the indices $1$, $2$ and $3$ corresponding to the Euler angles $\alpha$, $\beta$ and $\gamma$, respectively.

\subsubsection{Matrices $\MlL_1$ and $\MlL_3$}

Due to~\eqref{eq:deri_alpha} and~\eqref{eq:deri_gamma}, the derivatives of the Wigner-D matrices with respect to the Euler angles $\alpha$ and $\gamma$, $\MlL_1$ and $\MlL_3$, are two identical and diagonal matrices 

\begin{equation}
\label{eq:M1}
        \forall (\bbm,k)\in \calM_\bl\times\calM_L, \qquad \MlL_1[(\bbm,k),(\bbm,k)] = \MlL_3[(\bbm,k),(\bbm,k)] = -\left(\sum_{i=1}^N m_i -k\right)\cdot \mathrm{i},
\end{equation}
with $\calM_L$ defined in~\eqref{eq:calMU2} and $\calM_\bl$ defined in~\eqref{eq:calMbl}.
The following proposition holds trivially. 

\begin{proposition}\label{prop:ker-M1-GE}
    Let $\bl\in \calL^N$, $L\in \calLG$, then for $\MlL_1$ and $\MlL_3$ being defined in~\eqref{eq:M1}, there holds
    \[
        \ker(\MlL_1) = \ker(\MlL_3) = \{\bc=(\bc_{\bbm,k})_{(\bbm,k)\in\calM_\bl\times\calM_L}\in\F^{\card{\calM_\bl\times\calM_L}}: \bc_{\bbm,k} = 0 \textup{ if } \sum_{i=1}^N m_i -k \neq 0\}.
    \]
\end{proposition}


For simplicity, we denote $\sum\bbm := \sum_{i=1}^N m_i$ for any vector of scalars $\bbm$ of length $\Nat$. 
Additionally, for $\bl\in \calL^N$, we define $\sum\bl:=\sum\boldsymbol{\ell}$, with
$\boldsymbol{\ell}$ consisting of the $\ell$-components of $\bl$.
As a direct consequence of Proposition~\ref{prop:ker-M1-GE} and~\eqref{eq:ker_inclusion}, 
\begin{equation}\label{eq:kerM-nullpart}
    \forall\bc=(\bc_{\bbm,k})_{(\bbm,k)\in\calM_\bl\times\calM_L}\in\ker(\MlL), \quad \bc_{\bbm,k} \ne 0 \textup{ only if } k=\sum\bbm\in\calM_L.
\end{equation}
That is, we may restrict the matrix $\MlL$ only to the following column index set 
\begin{equation}\label{nonzero-ker}
    \big\{(\bbm,\sum \bbm): \;  \bbm\in\calM_\bl, \; \sum \bbm \in\calM_L\big\} = \Big\{(\bbm,\sum \bbm):\;\bbm\in\calM_\bl, \;|\sum \bbm|\le L\Big\}.
\end{equation}
Additionally, we have the following result on the dimensionality of the GE space in certain cases. 

\begin{proposition}
Let $\bl\in\calL^N$ and $L \in \calLG$  be such that $(L+\sum\bl)\notin\Z$. There holds
\[ 
    \dimge{\bl}{L} = 0.
\]
\end{proposition}

\begin{proof}
We see from Proposition~\ref{prop:GE_basis} that $\dimge{\bl}{L} = \dim\big(\ker(\MlL)\big)$. 
Note that $k\in\calM_L$ and $\bbm\in\calM_\bl$, we have $k-L\in\Z$ and $\sum\bbm-\sum\bl\in\Z$. When $\big(L+\sum\bl\big)\notin\Z$, the condition $\sum\bbm - k = 0$ is never fulfilled. Thus, due to~\eqref{eq:kerM-nullpart}, $\ker(\MlL)=\{\mathbf{0}\}$. Therefore, $\dim(\ker(\MlL))=0$, which leads to the desired result. 
\end{proof}

\subsubsection{Matrix $\MlL_2$}

We now turn to investigate the matrix $\MlL_2$. 
As suggested in~\eqref{eq:kerM-nullpart} , we consider for $\MlL_2$ only the restricted column index set~\eqref{nonzero-ker}.} In fact, the matrix $\MlL_2$ can be arranged to be a block matrix whose blocks are structured with respect to the values of $\sum \bbm$. To facilitate the writing, we denote
\begin{equation}
\label{eq:calMK}
       \calM_{\bl,K} := \{\bbm\in\calM_\bl:\sum\bbm=K\}.
\end{equation}
Then, it is natural to identify the set~\eqref{nonzero-ker} to $\Big\{\bbm\in\calM_\bl:|\sum\bbm|\le L
\Big\}$, and denote it by
\[
     \calM_{\bl,\le L} := \bigcup_{K\in\calM_L}\calM_{\bl,K}.
\]
As a consequence of~\eqref{eq:deri_beta}, for any column index $\bbm'\in\calM_{\bl,\le L}$ of the matrix $\MlL_2$, the row indices $(\bbm,k)$ that can lead to non-zero elements can either be $\bbm=\bbm'$ with $k = \sum \bbm' \pm 1$, or $\bbm$ satisfying for some $j\in\{1,2,\ldots,N\}$
\[
    \bbm^{\pm}_j := [m_1,m_2,\ldots,m_{j-1},m_j\pm1,m_{j+1},\ldots,m_N]=\bbm',
\]
with $k=\sum\bbm'$.
Precisely, we define for $(\bbm,k)\in \calM_\bl\times\calM_L$, $\bbm'\in \calM_{\bl,\le L}$
\begin{equation}
        \MlL_2[(\bbm,k),\bbm'] = 
    \begin{cases}
        \sqrt{(L+\sum\bbm+1)(L-\sum\bbm)}, & \qquad  \textup{if }\bbm'=\bbm,\;\;\; k = \sum\bbm+1, \\
        -\sqrt{(L-\sum\bbm+1)(L+\sum\bbm)}, & \qquad  \textup{if }\bbm'=\bbm,\;\;\; k = \sum\bbm-1, \\
        -\sqrt{(l_j+m_j+1)(l_j-m_j)}, & \qquad \textup{if }\bbm' = \bbm_j^+, \; k = \sum\bbm+1, \\
        \sqrt{(l_j-m_j+1)(l_j+m_j)}, & \qquad \textup{if }\bbm' = \bbm_j^-, \; k = \sum\bbm-1, \\
        0, & \qquad \textup{otherwise},
    \end{cases}
    \label{eq:M2}
\end{equation}
which is a direct consequence of~\eqref{eq:M_GE}, using~\eqref{eq:deri_beta}, where we have also removed a constant factor 1/2 over the whole matrix for simplicity, as it does not change its kernel.


To characterize the row indices $\bbm$ such that $\bbm^\pm_j\in\calM_\bl$, we prove the following lemma. 

\begin{lemma}\label{lem:row_ind}
Let $\bl\in\calL^N$ and $L \in \calLG$, then there hold
\begin{equation}
    \label{eq:mj+}
    \{\bbm\in\calM_{\bl}:\; \exists j\in \{ 1,\ldots, \Nat\},\; \bbm_j^+ \in \calM_{\bl,K+1}\} = \begin{cases}
        \calM_{\bl,K}, & \;\; \textup{if }K\in\calM_{L+1}\backslash\{L,L+1\}, \; K<\sum\bl, \\
        \varnothing, & \;\; \textup{otherwise,}
    \end{cases}
\end{equation}
and 
\begin{equation}
    \label{eq:mj-}
    \{\bbm\in\calM_{\bl}:\; \exists j\in \{ 1,\ldots, \Nat\},\;\bbm_j^- \in \calM_{\bl,K-1}\} = \begin{cases}
        \calM_{\bl,K}, & \;\; \textup{if }K\in\calM_{L+1}\backslash\{-L,-L-1\}, \; K>-\sum\bl, \\
        \varnothing, & \;\; \textup{otherwise.}
    \end{cases}
\end{equation}
\end{lemma}
\begin{proof}
    First, if $K\ge\sum\bl$, then $\calM_{\bl,K+1}$ is empty, and~\eqref{eq:mj+} holds trivially. 
    Otherwise, assume $K\in\calM_{L+1}\backslash\{L,L+1\}$ and $K<\sum\bl$, then it is straightforward that if $\bbm_j^+\in\calM_{\bl,K+1}$, $\bbm\in\calM_{\bl,K}$. For $\bbm\in\calM_{\bl,K}$,
    if for all $j\in \{1,\ldots,\Nat\}$, $\bbm_j^+ \notin \calM_{\bl,K+1}$, then it means that $\bbm_j^+ \notin \calM_{\bl}$, and 
    $\bbm = (l_1, \ldots, l_\Nat)$.
    Hence, 
    $K=\sum\bbm = \sum\bl$, which is impossible since $K<\sum\bl$. This proves~\eqref{eq:mj+}.
    A similar proof works for~\eqref{eq:mj-}. 
\end{proof}


From the expression of $\MlL_2$ given in~\eqref{eq:M2} and Lemma~\ref{lem:row_ind}, the matrix $\MlL_2$, after removing zero rows, can be arranged to have a block structure with the block 
columns being indexed by $\calM_{\bl,K}$ for $K\in\calM_L$, and 
the block rows being indexed by 
\begin{equation}\label{eq:MlK_minus}
    \big\{ (\bbm, {\sum}\bbm+1), \;  \bbm \in \calM_{\bl,K} \big\},
\quad K \in \calM_{L+1} \backslash \{ L, L+1\}, \quad K<\sum\bl,
\end{equation}

and 
\begin{equation}\label{eq:MlK_plus}
\big\{ (\bbm, {\sum} \bbm-1), 
\; \bbm \in \calM_{\bl,K} \big\},
\quad K \in \calM_{L+1} \backslash \{ -L-1, -L\}, \quad K>-\sum\bl.
\end{equation}
Since the second component of each element in the sets~\eqref{eq:MlK_minus} and~\eqref{eq:MlK_plus} 
is uniquely determined by its first component $\bbm\in\calM_{\bl,K}$, both sets can be identified as $\calM_{\bl,K}$. We note that the restrictions $K<\sum\bl$ and $K>-\sum\bl$ have an impact only when $L>\sum\bl$.

Such a block structure of the matrix $\MlL_{2}$ is illustrated in Figure~\ref{fig:block_structure_of_M}(a), where the yellow ($A^-_K$), pink ($A^+_K$), blue ($\Bm_K$), and purple ($\Bp_K$) blocks respectively correspond to the four non-zero cases in \eqref{eq:M2}. 
The upper half of the matrix corresponds to the row blocks having indices in $\calM_{L+1} \backslash \{ L, L+1\}$ while those for the lower half are in $\calM_{L+1} \backslash \{ -L-1, -L\}$. For $L=0$, in particular, the first two non-zero cases in \eqref{eq:M2} vanish because for any $\bbm\in\calM_{\bl,0}$, 
$\sum\bbm\pm1=\pm1\notin\calM_{0}$. 
As a result, $M^{\bl,0}_{2}$ consists of only two blocks (one blue and one purple).

\begin{figure}[ht]
    \centering
    \includegraphics[width=\linewidth]{Figures/Block_structure_5.pdf}
    \caption{Block structure of matrix $\MlL_2$: (a) $\MlL_2$: original derivation; (b) $\widetilde{\MlL}_2$: Row reordering of the $\MlL_2$ to enable a fast kernel solver; (c) $\MlL^{\textup{up}}$: upper half of $\widetilde{\MlL}_2$; the final matrix for which we find the kernel. In the plots, the red, green and blue boxes correspond to the matrix in the cases $L<\sum\bl,\; L=\sum\bl$, and $L>\sum\bl$, respectively. All the blocks left blank are zeros.
    }
    \label{fig:block_structure_of_M}
\end{figure}


To characterize the kernel of $M_2$, we first switch the ordering of its row blocks as shown in Figure~\ref{fig:block_structure_of_M}(b), which highly simplifies the numerical computation of coupling coefficients (c.f. Section~\ref{sec:num-complexity}).
We denote the rearranged matrix by $\widetilde{\MlL}_2$.
As illustrated in Figure~\ref{fig:block_structure_of_M}(b), the matrix $\widetilde{\MlL}_2$ admits a natural partition into two parts
\[
    \widetilde{\MlL}_2 = \begin{pmatrix}
         \Mup \\
         \Mlow \\
     \end{pmatrix}.
\]

It turns out that either the upper or the lower half of $\widetilde{\MlL}_2$ (Figure~\ref{fig:block_structure_of_M}(b)) is sufficient to determine its kernel. 
Thus, we consider only the upper half $\Mup$ presented in Figure~\ref{fig:block_structure_of_M}(c), which is a block matrix with column blocks indexed by $K'\in\calM_L$ with column indices $\bbm'\in\calM_{\bl,K'}$ and row blocks indexed by 
 $K\in
    \calM_{L+1}\backslash \{L,-L-1\}
$
with row indices $\bbm\in\calM_{\bl,K}$. 
The nonzero blocks of $\Mup$ read 
for $K\in \{-L, \ldots, L-1\}$
\begin{align}
    \Mup_{K,K}[\bbm,\bbm'] &= \sqrt{(L+K+1)(L-K)}\delta_{\bbm,\bbm'} \label{eq:MKK}\\
    \Mup_{K,K+1}[\bbm,\bbm'] &= 
    -\sqrt{(l_j+m_j+1)(l_j-m_j)}\delta_{\bbm',\bbm_j^+},
\end{align}
and for $K = L+1$ 
\begin{equation}\label{eq:MLpL}
    \Mup_{L+1,L}[\bbm,\bbm'] = 
    \sqrt{(l_j- m_j+1)(l_j+  m_j)}\delta_{\bbm',\bbm_j^-}.
\end{equation}

We formalize our claim in the following proposition.

\begin{proposition}\label{prop:ker-M2-GE}
Let $\bl\in\calL^N$, $L \in \calL$ and $\MlL_2$ be defined in~\eqref{eq:M2}. 
The matrix $\Mup$ defined in~\crefrange{eq:MKK}{eq:MLpL} has full rank and satisfies
\[
    \ker(\Mup) =  \ker(\MlL_2).
\]
\end{proposition}
The proof of this result will be presented in Section~\ref{sec:proof-half-mat}. 

\begin{remark}[Size of the matrix $\Mup$]\label{remark:size-M2-GE}
From the expression of the matrix $\Mup$ given in~\crefrange{eq:MKK}{eq:MLpL}, 
the number of columns of the matrix $\Mup$ is always 
$\displaystyle\sum_{K\in\calM_L}\card{\calM_{\bl,K}}$
while the number of rows is always 
$\displaystyle\sum_{K\in\calM_{L+1} \backslash \{ L, L+1\}}\card{\calM_{\bl,K}}$.
In particular, when $L\ge\sum\bl$, the sets $\calM_{\bl,\pm(L+1)}$ are empty and thus $\Mup_{L+1,L}$ vanishes, leaving only the green box as shown in Figure~\ref{fig:block_structure_of_M}(c), 
and when $L>\sum\bl$, all the sets $\calM_{\bl,K}$ with $|K|>\sum\bl$ are empty, leaving only a square matrix (the blue box in Figure~\ref{fig:block_structure_of_M}(c)).
Hence, depending on the relative ordering of $L$ and $\sum\bl$, 
\begin{equation*}
    \displaystyle\sum_{K\in\calM_{L+1} \backslash \{ L, L+1\}}\card{\calM_{\bl,K}} = 
    \begin{cases}
    \displaystyle {\sum_{K\in\calM_L}\card{\calM_{\bl,K}}} + \card{\calM_{\bl,L+1}}-\card{\calM_{\bl,L}}, & \quad \textup{if }L<\sum\bl, \\
    \displaystyle \sum_{K\in\calM_L}\card{\calM_{\bl,K}}-\card{\calM_{\bl,L}}, & \quad \textup{if }L=\sum\bl, \\
    \displaystyle \sum_{K\in\calM_L}\card{\calM_{\bl,K}}, & \quad \textup{if }L>\sum\bl.
\end{cases}
\end{equation*}
\end{remark}

\subsubsection{GE space and its dimensionality}

We now turn to characterize the GE space $V^{\bl,L}$, including its dimensionality and a basis. 

\begin{proposition}\label{prop:ge-su2}
Let $\bl\in\calL^N$ and $L \in \calL$. Let
\[
    \bc^{\bl,L}_i = \{c_{\bbm,i}^{\bl,L}\}_{\bbm\in\calM_{\bl,\le L}},\quad i=1,\ldots,\dim\big(\ker(\Mup)\big)
\]
be a basis of $\ker(\Mup)$, with $\Mup$ given in~\crefrange{eq:MKK}{eq:MLpL}.
Then, a basis of $V^{\bl,L}$ is given by 
\begin{equation}\label{eq:GE-basis-su2}
     b^{\bl,L}_i = 
    \left[
        \sum_{\bbm \in \calM_{\bl,K}}c_{\bbm,i}^{\bl,L}\phi^\bl_\bbm
    \right]_{K\in\calM_L},\qquad i=1,\ldots,\dim\big(\ker(\Mup)\big).
\end{equation}
In addition, 
\begin{equation}\label{eq:dim_su2}
    \dimge{\bl}{L}=\dim\big(\ker(\Mup)\big)=\card{\calM_{\bl,L}} - \card{\calM_{\bl,L+1}}.
\end{equation}
\end{proposition}
\begin{proof}
This is a direct result of Propositions~\ref{prop:GE_basis},~\ref{prop:ker-M1-GE} and~\ref{prop:ker-M2-GE}. Specifically, $\dimge{\bl}{L}$ can be determined by using the fact that $\Mup$ is full rank whose size is given in Remark~\ref{remark:size-M2-GE}.
\end{proof}

In the GE case, 
the dimensionality can be further expressed as follows.

\begin{proposition}[Exact GE dimensionality]
    For $\bl\in\calL^N$ and $L\in\calLG$, we have
    \begin{equation}\label{eq:explicit_dim_GE}
        \dimge{\bl}{L}=\sum_{{\mathbf S}\subseteq \{1,2,\ldots,N\}}
        (-1)^{\left( 1+\sum {\mathbf S} \right)} 
        \;\binom{L+\sum\bl-\sum_{s\in{\mathbf S}}(2l_s+1)+N-1 }{N-2}.
    \end{equation}
\end{proposition}
\begin{proof} By definition of $\calM_{\bl,L}$, $\card{\calM_{\bl,L}}$ is the coefficients of the term $x^L$ in the polynomial
$
    \prod_{i=1}^N \left(\sum_{m=-l_i}^{l_i} x^m \right).
$
Therefore, by the binomial theorem,
\begin{equation*}
    \card{\calM_{\bl,L}}=\sum_{{\mathbf S}\subseteq \{1,2,\ldots,N\}}(-1)^{\sum {\mathbf S}} \binom{L+\sum\bl-\sum_{s\in{\mathbf S}}(2l_s+1)+N-1}{ N-1}.
\end{equation*}
From \eqref{eq:dim_su2}, the dimensionality of the GE basis is,
\begin{align*}
    \dimge{\bl}{L}&=\sum_{{\mathbf S}\subseteq \{1,2,\ldots,N\}}(-1)^{(1+\sum {\mathbf S})} \left(\binom{L+\sum\bl-\sum_{s\in{\mathbf S}}(2l_s+1)+N-1 }{ N-1}-\binom{L+\sum\bl-\sum_{s\in{\mathbf S}}(2l_s+1)+N }{ N-1}\right)\\
    &=\sum_{{\mathbf S}\subseteq \{1,2,\ldots,N\}}(-1)^{(1+\sum {\mathbf S})} \binom{L+\sum\bl-\sum_{s\in{\mathbf S}}(2l_s+1)+N-1 }{ N-2},
\end{align*}
using the binomial formula $\binom{n}{k}+\binom{n}{k+1}=\binom{n+1}{k+1}$.
\end{proof}


\begin{remark}[Link with~\cite{Dusson2022-ie}]
\label{rem:ace}
Note that~\cite[Proposition 5]{Dusson2022-ie}
provides a result on the dimensionality of the space $V^{\bl,0}$ for $G=SO(3)$ as the cardinality of the set
\[
\bigg\{ \bL = (L_2,L_3,\ldots,L_N)\in \N^{N-1}, 
      \quad 
      |l_1-l_2| \le L_2 \le l_1+l_2, \\
       \forall \; 3\le i\le N, \;
      |L_{i-1}-l_i| \le L_i \le {L_{i-1}+l_i}, L_N =0
      \bigg\}.
\]
We note that~\eqref{eq:explicit_dim_GE} is more explicit and more general than what was obtained in~\cite{Dusson2022-ie}.
\end{remark}

\begin{example}\label{exam:dim_cg_su2}
As an immediate result of~\eqref{eq:explicit_dim_GE}, the dimensionality of the Clebsch--Gordan coefficients for $SU(2)$ and $SO(3)$ can at most be 1.
Indeed let $L,L_1,L_2\in\calL_G$, then 
\begin{equation*}
        \dimge{(L_1,L_2)}{L}= \begin{cases}
        1, & \textup{if } L+L_1+L_2\in\Z, \; |L_1-L_2|\le L\le L_1+L_2,\\
        0, & \textup{otherwise}.
    \end{cases}
\end{equation*}
Consequently, for $\bl\in\calL^N$ and $L\in\calLG$, the recursive formula~\eqref{eq:dim_GE_rec} for the dimensions becomes
\begin{equation}\label{eq:dim_recursive_su2}
    \dimge{\bl}{L} = \sum_{L_1=0}^{\sum\bl^{(1)}}\;\sum_{L_2=|L-L_1|}^{\min\{\sum\bl^{(2)},L_1+L\}}\dimge{\bl^{(1)}}{L_1} \; \dimge{\bl^{(2)}}{L_2}.
\end{equation}
\end{example}

With~\eqref{eq:dim_recursive_su2}, we have an easy way to recursively compute GE dimensionality for any $\bl$. 
For more specific values, we provide dimensionalities for different values of $\bl$ and $L$ in 
Appendix~\ref{app:table_of_dim}.

\medskip

While \eqref{eq:explicit_dim_GE} is explicit, the sum contains a total of $2^N$ terms, which becomes increasingly difficult to access as $N$ grows, and when $\bl$ contains different values. 
Asymptotically, we have the following proposition.



\begin{proposition}[Asymptotic dimensionality]\label{prop:asym-dim-ge}
    For $\bl\in\calL^N$, $L\in\calLG$, $L+\sum\bl\in\Z$, define
    \[
        \textup{Var}_\bl = \sum_{i=1}^N \frac{l_i(l_i+1)}{3}.
    \]
    If there exist $l_{\min},l_{\max}> 0$ independent of $\Nat$ such that
    \begin{equation}\label{eq:bounded-ll}
    \forall i=1,\ldots,\Nat, \quad l_{\min} \le l_i\le l_{\max},\;
    \end{equation}
    then for $L\ll\sqrt{\textup{Var}_\bl}$, there holds for $\Nat$ sufficiently large that
    \begin{equation}
        \label{eq:dimge-est}
        \dimge{\bl}{L} = \prod_{i=1}^N\big(2l_i+1\big)
        \left(\frac{2L+1}{2\sqrt{2\pi}(\textup{Var}_\bl)^{3/2}}+O\left(\frac{1}{\Nat^{5/2}}\right)\right).
    \end{equation}
\end{proposition}

Proposition~\ref{prop:asym-dim-ge} suggests that we may approximate $\dimge{\bl}{L}$ by 
\[ \dimgeest{\bl}{L} = \frac{(2L+1)\big(\prod_{i=1}^N(2l_i+1)\big)}{2\sqrt{2\pi}(\textup{Var}_\bl)^{3/2}} \]
in the asymptotic regime, and that the dimensionalities of the equivariant function spaces scale linearly with respect to $L$, the order of equivariance, when $L$ is not too large. A proof of this Proposition is given in Section~\ref{sec:proof-asym-dim-ge}.



\subsection{Construction of multi-variable GE-PI bases}
As illustrated in Section~\ref{sec:ge-pi}, most of the arguments for the GE case can be applied to the GE-PI case directly, with an exception that the set $\calM_\bl$ is replaced by $\oMl$ for given $\bl\in\calL^N$. Hence, before presenting the construction of multi-variable GE-PI bases and the properties of the matrix~\eqref{eq:matM-PI}, 
we first clarify the form of $\oMl$ in this specific case. 

\begin{remark}
For the compatible bases given in Subsection~\ref{sec:compatible_bases}, the index $\bl\in\calL^N$ can either be a tuple of scalars or a tuple of vectors. In the previous case, the definition of $\oMl$ is clear, 
while in the latter case, for example when $\bl = \big((n_i,\ell_i)\big)_{i=1}^N$, the set $\oMl$ depends not only on the $\ell$-component but also on the $n$-component. Specifically, there hold
\[
    \oMl = \big\{\obm:\bbm\in\calM_\bl\big\},
\]
where the equivalent class $\obm$ is defined as
\begin{equation*}
    \obm  = \{\pi\bbm:\pi\in S_\bl\}
    \quad\text{with}\quad 
    S_\bl = \{\pi\in S_N:\pi\bl=\bl\} = \{\pi\in S_N: \pi\bn=\bn \textup{ and } \pi\boldsymbol{\ell}=\boldsymbol{\ell}\},
\end{equation*}
where $\bn=\left(n_i\right)_{i=1}^N$ and $\boldsymbol{\ell}=\left(\ell_i\right)_{i=1}^N$.
\end{remark}

\begin{remark}\label{rem:id-ell-blocks}
    In the case where $\bl$ is a tuple of vectors as above,  while the set $S_\bl$ depends on $\bn$, the set $\oMl = \bigotimes_{j=1}^\Nb \overline{\calM}_{\boldsymbol{\ell}^{(j)}}$ depends only on the $\boldsymbol\ell$-component of the minimal partition of $\bl$, and so do the coupling coefficients. In other words, if $\bl,\;\bl'\in\calL^N$ induce the same $\boldsymbol\ell = (\boldsymbol\ell^{(1)},\ldots,{\boldsymbol\ell}^{(\Nb)})$ in their minimal partitions, they share the same coupling coefficients, despite the potential differences in the underlying compatible bases. 
\end{remark}

With a slight abuse of notation, we still use $\MlL$ to denote the matrix $M^{\bl,L}$ defined in~\eqref{eq:matM-PI}, and use $\MlL_d,\;d=1,2,3$ for the three parts of $\MlL$. 
\subsubsection{Matrices $\MlL_1$ and $\MlL_3$}

Following a similar argument as above, $\MlL_1$ and $\MlL_3$ in the GE-PI case are also two identical and diagonal matrices with diagonal elements
\begin{equation}
\label{eq:M1_PI}
        \forall (\bbm,k)\in {\oMl}\times\calM_L, \qquad \MlL_1[(\bbm,k),(\bbm,k)] = \MlL_3[(\bbm,k),(\bbm,k)] = \left(\sum_{i=1}^N m_i -k\right)\cdot \mathrm{i},
\end{equation}
where $\calM_L$ is defined in~\eqref{eq:calMU2} and  $\oMl$ is defined in~\eqref{eq:GE_class}. 
Hence, the following proposition holds true in the GE-PI case as well.

\begin{proposition}\label{prop:ker-M1-GEPI}
    Let $\bl\in \calL^N$, $L\in \calLG$, then for $\MlL_1$ and $\MlL_3$ being defined in~\eqref{eq:M1_PI}, there holds
    \[
        \ker(\MlL_1) = \ker(\MlL_3) = \{\bc=(\bc_{\bbm,k})_{(\bbm,k)\in\oMl\times\calM_L}\in\F^{\card{\oMl\times\calM_L}}: \bc_{\bbm,k} = 0 \textup{ if } \sum\bbm -k \neq 0\}.
    \]
    Consequently, if $(L+\sum\bl)\notin\Z$, then
    \[ 
        \dimpi{\bl}{L}=
        0.
    \]
\end{proposition}


\subsubsection{Matrix $\MlL_2$}
Using Proposition~\ref{prop:ker-M1-GEPI}, we only need to set up a system for the non-zero coefficients, that is, the column indices of $\MlL_2$ can be restricted to the set
\[
    \big\{(\bbm,\sum \bbm), \quad  \bbm\in\oMl, \quad |\sum \bbm| \le L\big\}\subset\oMl\times\calM_L.
\]
We may identify this set to 
\[
    \overline{\calM}_{\bl,\le L} := \Big\{\bbm\in\oMl:|\sum \bbm| \le L
    \Big\}=\bigcup_{K\in\calM_L}\oMlK, 
\]
where 
\[
    \oMlK = \{\bbm\in\oMl: \; \sum\bbm=K\}.
\]

As a direct consequence of~\eqref{eq:M_GE_PI}, using~\eqref{eq:deri_beta}, the matrix $\MlL_2$ in the GE-PI case is defined for $\bbm'\in \overline{\calM}_{\bl,\le L}$, $(\bbm,k)\in \oMl\times\calM_L$ by
\begin{equation}
        \MlL_2[(\bbm,k),\bbm'] = 
    \begin{cases}
        \sqrt{(L+\sum\bbm+1)(L-\sum\bbm)}, & \qquad  \textup{if }\bbm'=\bbm,\; k = \sum\bbm+1, \\
        -\sqrt{(L-\sum\bbm+1)(L+\sum\bbm)}, & \qquad  \textup{if }\bbm'=\bbm,\; k = \sum\bbm-1, \\
        -\blamb{}{\bbm'^{(j)},p+1}\sqrt{(\bl^{(j)}_1+p+1)(\bl^{(j)}_1-p)}, & \qquad \textup{if }\bbm' = \bbm_{(j,p,p+1)}, \; k = \sum\bbm', \\
        \blamb{}{{\bbm'}^{(j)},p-1}\sqrt{(\bl^{(j)}_1-p+1)(\bl^{(j)}_1+p)}, & \qquad \textup{if }\bbm' = \bbm_{(j,p,p-1)}, \;  k = \sum\bbm', \\
        0, & \qquad \textup{otherwise},
    \end{cases}
    \label{eq:M2-PI}
\end{equation}
where we have also removed the constant factor 1/2 for simplicity.
Following a similar discussion as in Lemma~\ref{lem:row_ind}, we see that the row index $\bbm\in\oMl$ that can lead to nonzero elements in~\eqref{eq:M2-PI} lie in the sets 
\[
\big\{ (\bbm, {\sum}\bbm+1), \;  \bbm \in \overline{\calM}_{\bl,K} \big\},
\quad K \in \calM_{L+1} \backslash \{ L, L+1\},
\]
and 
\[
\big\{ (\bbm, {\sum} \bbm-1), 
\; \bbm \in \overline{\calM}_{\bl,K} \big\},
\quad K \in \calM_{L+1} \backslash \{ -L-1, -L\}.
\]
Hence, we are able to sort the rows and columns using the same convention as that of the GE case. 

Then, after removing the zero rows, the above-defined $\MlL_2$ also has a block structure illustrated in Figure~\ref{fig:block_structure_of_M}(a).
Rearranging the block rows in the manner shown in Figure~\ref{fig:block_structure_of_M}(b), we obtain a reordered matrix $\widetilde{\MlL}_2$, which admits a partition into two parts, that is, 
\[
    \widetilde{\MlL}_2 = \begin{pmatrix}
         \Mup \\
         \Mlow \\
     \end{pmatrix}.
\]
We claim that either half of $\widetilde{\MlL}_2$ is sufficient to determine its kernel, and consider only its upper half $\Mup$, which is a block matrix with column blocks indexed by $K'\in\calM_L$ and row blocks indexed by $K\in\calM_{L+1} \backslash \{ L, L+1\}$, and with nonzero blocks being defined 
for $K\in \{-L, \ldots, L-1\}$ by
\begin{align}
    \Mup_{K,K}[\bbm,\bbm'] &= \sqrt{(L+K+1)(L-K)}\delta_{\bbm',\bbm} \label{eq:MKK-PI}\\
    \Mup_{K,K+1}[\bbm,\bbm'] &= 
    -\blamb{}{{\bbm'}^{(j)},p+1}\sqrt{(\bl^{(j)}_1+p+1)(\bl^{(j)}_1-p)}\;\delta_{\bbm',\bbm_{(j,p,p+1)}}
\end{align}
and for $K = L+1$ by
\begin{equation}\label{eq:MLpL-PI}
    \Mup_{L+1,L}[\bbm,\bbm'] = 
    \blamb{}{{\bbm'}^{(j)},p-1}\sqrt{(\bl^{(j)}_1-p+1)(\bl^{(j)}_1+p)}
    \;\delta_{\bbm' ,\bbm_{(j,p,p-1)}}.
\end{equation}
As in the GE case, the following proposition holds true.
\begin{proposition}\label{prop:ker-M2-GEPI}
Let $\bl\in\calL^N$, $L \in \calL$ and $\MlL_2$ be defined in~\eqref{eq:M2-PI}. 
The matrix $\Mup$ given in~\crefrange{eq:MKK-PI}{eq:MLpL-PI} has full rank and satisfies
\[
    \ker(\Mup) = \ker(\MlL_2).
\]
\end{proposition}
The proof of this result will be presented in Section~\ref{sec:proof-half-mat-pi}. Notably, $\Mup$ in the GE-PI case also has the structure shown in Figure~\ref{fig:block_structure_of_M}(c). 
\begin{remark}[Size of the matrix $\Mup$]\label{remark:size-M2-GEPI}
From the expression of the matrix $\Mup$ given in~\crefrange{eq:MKK-PI}{eq:MLpL-PI}, 
the number of columns of the matrix $\Mup$ is always 
$\displaystyle\sum_{K\in\calM_L}\card{\oMlK}$
while the number of rows is always 
$\sum_{K\in\calM_{L+1} \backslash \{ L, L+1\}}\card{\oMlK}$. 
Hence, depending on the relative ordering of $L$ and $\sum\bl$, 
\begin{equation*}
    \displaystyle\sum_{K\in\calM_{L+1} \backslash \{ L, L+1\}}\card{\oMlK} = 
    \begin{cases}
    \displaystyle {\sum_{K\in\calM_L}\card{\oMlK}} + \card{\overline{\calM}_{\bl,L+1}}-\card{\overline{\calM}_{\bl,L}}, & \quad \textup{if }L<\sum\bl, \\
    \displaystyle \sum_{K\in\calM_L}\card{\oMlK}-\card{\overline{\calM}_{\bl,L}}, & \quad \textup{if }L=\sum\bl, \\
    \displaystyle \sum_{K\in\calM_L}\card{\oMlK}, & \quad \textup{if }L>\sum\bl.
\end{cases}
\end{equation*}
\end{remark}

\subsubsection{GE-PI space and its dimensionality}
Note that the structure of the matrix $M$ remains almost unchanged compared with the GE case, with the only difference being that the index sets $\calM_{\bl,K}$ are replaced by $\oMlK$. 
In addition, the following proposition holds true.

\begin{theorem}\label{thm:gepi-su2}
Let $\bl\in\calL^N$, $L \in \calL$, and let
\[
    \bc^{\bl,L}_i = \{c_{\bbm,i}^{\bl,L}\}_{\bbm\in\overline{\calM}_{\bl,L}},\quad i=1,\ldots,\dim\big(\ker(\Mup)\big)
\]
be a basis of $\ker(\Mup)$, where $\Mup$ is given in~\crefrange{eq:MKK-PI}{eq:MLpL-PI}.
Then, a basis of $\bV^{\bl,L}$
is given by 
\begin{equation*}
     b^{\bl,L}_i = 
    \left[
        \sum_{\bbm \in \oMlK}c_{\bbm,i}^{\bl,L}
        \bphi{\bl}{\bbm}
    \right]_{K\in\calM_L},\quad i=1,\ldots,\dim\big(\ker(\Mup)\big).
\end{equation*}
In addition,
\begin{equation}\label{eq:dim_su2_pi}
    \dimpi{\bl}{L}=\dim\big(\ker(\Mup)\big)=\card{\overline{\calM}_{\bl,L}} - \card{\overline{\calM}_{\bl,L+1}}.
\end{equation}
\end{theorem}
\begin{proof}
This is a direct result of Theorem~\ref{thm:GEPI_basis}, Propositions~\ref{prop:ker-M1-GEPI} and~\ref{prop:ker-M2-GEPI}. Specifically, $\dimpi{\bl}{L}$ can be determined by using the fact that $\Mup$ is full rank, whose size is given in Remark~\ref{remark:size-M2-GEPI}.
\end{proof}

\begin{example}
    In the specific case where $\bl=(1,\ldots,1)\in\N^N$ or $\bl=\left((\boldsymbol{\ell}_i,\bn_i)\right)_{i=1}^{\Nat}\in\calL^N$ with $\boldsymbol{\ell}=(1,\ldots,1)\in\N^N,\; \bn=(n,\ldots,n)\in\N^N$, we have for $L\in\calLG$ that
    \[
        \card{\overline\calM_{\bl,L}} =\left\{
    \begin{array}{ll}
        \lfloor\frac{\Nat-L}{2}\rfloor+1 & \textup{if } 0\le L\le \Nat \\
        0 & \textup{otherwise},
    \end{array}
    \right.
    \]
    and as a consequence, 
    \[
        \dimpi{\bl}{L} = \left\{
    \begin{array}{ll}
        1, & \text{if } L\le \Nat \text{ and } \Nat+L\text{ is even}, \\
        0, & \text{otherwise}.
    \end{array}
\right.
    \]
Indeed, the elements $\bbm\in\oMl$ 
have coordinates either -1, 0 or 1 and are thus uniquely determined by their count vectors $\blamb{}{\bbm}=(\blamb{-1}{\bbm},\blamb{0}{\bbm},\blamb{1}{\bbm})$ with $\blamb{i}{\bbm}\ge 0$ and $\blamb{-1}{\bbm}+\blamb{0}{\bbm}+\blamb{1}{\bbm}=\Nat$. Moreover, $\sum\bbm=L$ if and only if $\blamb{-1}{\bbm}+L=\blamb{1}{\bbm}$. Thus, $\bbm\in\overline\calM_{\bl,L}$ if and only if $2\blamb{-1}{\bbm}+\blamb{0}{\bbm}+L=\Nat$. If $L>\Nat$, this equation has no solution. If $0\le L\le\Nat$, there are $\lfloor\frac{\Nat-L}{2}\rfloor+1 $ possible values of $\blamb{-1}{\bbm}$ such that $N-L-2\blamb{-1}{\bbm}\ge0$ and for each possible value of $\blamb{-1}{\bbm}$, there is only one possible value for $\blamb{0}{\bbm}$, which is $N-L-2\blamb{-1}{\bbm}$. 
The dimensionality of the GE-PI basis is then a direct consequence of Theorem~\ref{thm:gepi-su2}. 
\end{example}

Similar to the discussion in Example~\ref{exam:dim_cg_su2}, we have the following corollary for the recursive formula of GE-PI dimensions. 
\begin{corollary}
    Let $N,N_1,N_2\in \N$ with $N_1+N_2=N$.
    Let $\bl\in\calL^\Nat$ and $L\in\calL$. 
    Suppose $\bl = (\bl^{(1)}, \bl^{(2)})\in\calL^{\Nat_1}\times\calL^{\Nat_2}$ with $\bl^{(1)} \cap \bl^{(2)} = \varnothing$, then
    \begin{equation}\label{eq:rec-su2-gepi}
        \dimpi{\bl}{L} = \sum_{L_1=0}^{\sum\bl^{(1)}}\;\sum_{L_2=|L-L_1|}^{\min\{\sum\bl^{(2)},L_1+L\}}\dimpi{\bl^{(1)}}{L_1} \; \dimpi{\bl^{(2)}}{L_2}.
    \end{equation}
\end{corollary}
With~\eqref{eq:rec-su2-gepi}, the GE-PI dimensionalities for non-identical $\bl\in\calL^N$ can be accessed via the knowledge of the dimensionalities of identical $\bl$'s. We provide in Appendix~\ref{app:table_of_dim} the GE-PI dimensionalities for some typical values of identical $\bl$ and $L$. 

We then turn to investigate the asymptotic dimensions for the GE-PI case. 
Suppose $(\bl^{(j)})_{j=1}^{\Nb}$ is the minimal partition of $\bl$ whose $j$-th block has length $N_j$, then the set $\overline\calM_{\bl,L}$ is isometric to
\begin{equation}\label{eq:Lamb_blL}
    \{\blamb{}{}\in \Lambda_\bl:\sum_{j=1}^{\Nb}\sum_{k=-\bl^{(j)}_1}^{\bl^{(j)}_1}\blamb{(j)}{k}\cdot k = L\},
\end{equation}
which is actually the set of solutions to the equations
\begin{equation}\label{eq:ge-pi-int-part}
    \begin{cases}
        \displaystyle \sum_{k=-\bl^{(j)}_1}^{\bl^{(j)}_1}\blamb{(j)}{k} & = N_j, \quad \forall j=1,\ldots,\Nb,\\ 
        \displaystyle \sum_{j=1}^{\Nb}\sum_{k=-\bl^{(j)}_1}^{\bl^{(j)}_1}\blamb{(j)}{k}\cdot k & = L.
    \end{cases}
\end{equation}
Hence, $\card{\overline\calM_{\bl,L}}$ is the number of solutions to~\eqref{eq:ge-pi-int-part}. When $\Nb=1$ in particular,~\eqref{eq:Lamb_blL}
degenerates to 
\begin{equation*}
    \left\{\blamb{}{}\in \Lambda_\bl:
        \displaystyle \sum_{k=-l}^{l}\blamb{}{k}  = N, 
        \quad
        \displaystyle \sum_{k=-l}^{l}\blamb{}{k}\cdot k  = L
        \right\},
\end{equation*}
where $l$ is the identical value in $\bl$.
In this case, we have that $\card{\overline\calM_{\bl,L}}$ is the coefficients of the term $x^{L+Nl}$ in the Gaussian binomial coefficient 
$
   \displaystyle \left[ \begin{matrix} N+2l \\ 2l \end{matrix} \right]_x = 
        \prod_{i=1}^{2l}\frac{1-x^{N+i}}{1-x^{i}},\quad x\neq 1.
$
With the properties of the Gaussian binomial coefficient being well-studied, an asymptotic estimate can be obtained for identical $\bl$'s. Indeed, denoting 
\[
    \textup{Var}_\bl = \frac{Nl(N+2l+1)}{6}, 
\]
we obtain using~\cite[(45)]{takacs1986asymptotic} that
\[
    \lim_{\Nat\to+\infty} \left(\frac{\card{\overline\calM_{\bl,L}}}{\binom{N+2l }{ 2l }} - \frac{\exp{(-L^2/(2\textup{Var}_\bl))}}{\sqrt{2\pi \textup{Var}_\bl}}\right) = 0,
\]
which yields
\[
  \dimpi{\bl}{L} = \binom{N+2l }{ 2l }\left(\frac{2L+1}{2\sqrt{2\pi} (\textup{Var}_\bl)^{3/2}}+o\left(\frac{1}{\Nat^{3/2}}\right)\right).
\]
Numerically, we observe a faster decay in the remainder of this estimation, which we present as the following conjecture, as its proof goes beyond the scope of this manuscript.
Note that the estimation below seems stronger than the results presented in a recent paper~\cite{melczer2020counting}, which require that both $l$ and $N$ go to infinity.

\begin{conjecture}[Asymptotic dimensionality for identical $\bl$]\label{conj:asym-dim-gepi-1block}
    Let $\bl = (l,\ldots,l)\in\calL^N$ containing identical nonzero values $l$, $L\in\calL_G$ such that $L+\sum\bl\in\Z$. 
    Then for $L\ll\sqrt{\textup{Var}_\bl}$,
    \begin{equation*}
                \dimpi{\bl}{L} = \binom{N+2l }{ 2l }\left(\frac{2L+1}{2\sqrt{2\pi} (\textup{Var}_\bl)^{3/2}}+O\left(\frac{1}{\Nat^{5/2}}\right)\right).
    \end{equation*}
\end{conjecture}

When extending to the multi-block case, we have that for general $\bl$, $\card{\overline\calM_{\bl,L}}$ is the coefficients of the term $x^{L+\sum_{j=1}^\Nb N_{j}\bl^{(j)}_1}$ in the following
\[
    \prod_{j=1}^\Nb\left[ \begin{matrix} N_{j}+2\bl^{(j)}_1 \\ 2\bl^{(j)}_1 \end{matrix} \right]_x = 
    \prod_{j=1}^\Nb\prod_{i=1}^{2\bl^{(j)}_1}\frac{1-x^{N_{j}+i}}{1-x^i},\quad x\neq 1.
\]
In addition, the sums of the blocks can be considered as independent variables. Hence, by using a similar strategy as shown in the proof of Proposition~\ref{prop:asym-dim-ge}, we have the following estimation with the increase of $\Nb$. 
\begin{proposition}[Asymptotic dimensionality for general $\bl$]\label{prop:asym-dim-gepi}
    For $\bl\in\calL^N$, $L\in\calLG$, $L+\sum\bl\in\Z$, where $\bl = (\bl^{(j)})_{j=1}^\Nb$ is the minimal partition of $\bl$ 
    with $N_j$ being the length of its $j$-th block, define
    \[
        \textup{Var}_\bl = \sum_{j=1}^\Nb \frac{N_{j}\bl^{(j)}_1(N_j+2\bl^{(j)}_1+1)}{6}.
    \]
    If there exist $\kappa_{\min},\kappa_{\max}> 0$ independent of $\Nb$ such that
    \begin{equation*}
    \forall j=1,\ldots,\Nb, \quad \kappa_{\min} \le N_j\bl^{(j)}_1\le \kappa_{\max},\;
    \end{equation*}
    then for $L\ll\sqrt{\textup{Var}_\bl}$, there holds for $\Nb$ sufficiently large
    \begin{equation}
        \label{eq:dimpi-est}
        \dimpi{\bl}{L} = \prod_{j=1}^\Nb \binom{N_j+2\bl^{(j)}_1 }{ 2\bl^{(j)}_1}
        \left(\frac{2L+1}{2\sqrt{2\pi}(\textup{Var}_\bl)^{3/2}}+O\left(\frac{1}{\Nb^{5/2}}\right)\right).
    \end{equation}
\end{proposition}


\begin{proof}
    The proof is similar to the proof of Proposition~\ref{prop:asym-dim-ge}, using the variance given in~\cite[(43)]{takacs1986asymptotic}, and the fourth-order cumulants
    \[
    C_{4\bl} = -\sum_{j=1}^\Nb \frac{N_j \bl^{(j)}_1 (N_j+2\bl^{(j)}_1+1)(N_j^2+2N_j
    \bl^{(j)}_1+N_j+4{\bl^{(j)}_1}^2+2\bl^{(j)}_1)}{60}.
    \]
\end{proof}

As Conjecture~\ref{conj:asym-dim-gepi-1block} and Proposition~\ref{prop:asym-dim-gepi} suggest, when the order of equivariance $L$ is not to large, we can estimate $\dimpi{\bl}{L}$ by 
\[ \dimpiest{\bl}{L} = \prod_{j=1}^\Nb \binom{N_j+2\bl^{(j)}_1 }{ 2\bl^{(j)}_1}\left(\frac{2L+1}{2\sqrt{2\pi}(\textup{Var}_\bl)^{3/2}}\right), \]
in the asymptotic regime. In addition, it suggests that the dimensionalities of the GE-PI spaces scale linearly with respect to $L$, for moderate $L$.

Let us mention that when $N_j=1, \;\forall j\in\{1,2,\ldots,\Nb\}$, the GE-PI case is identical to the GE case, and Proposition~\ref{prop:asym-dim-gepi} degenerates to Proposition~\ref{prop:asym-dim-ge}. 


\subsection{Extension to $O(3)$}

In a variety of physical implementations, the $O(3)$ symmetry is of greater interest than the $SO(3)$ one described above, i.e., considering not only rotation equivariance, but also reflection equivariance. Whereas $O(3)$ is not a connected group, our approach can still be applied to its two connected components separately. 

First, since $O(3)=\{\pm1\}\times SO(3)$, each element $R\in O(3)$ can be written as $R=sQ$ where $s:=\det(R)\in\{\pm1\}$ and $Q\in SO(3)$.
The irreducible representations of $O(3)$ are then defined for $R=sQ\in O(3)$ as
\[
    D^{\ell,+}(R) = D^\ell(Q),
    \qquad
    D^{\ell,-}(R) = sD^\ell(Q),
\]
for $\ell\in\N_0$, where $D^\ell$ stands for the Wigner-D matrix of degree $\ell$ (see Definition~\ref{def:wigner-D}). 
Given $l\in\calL$ with component $\ell\in\N_0$, all the compatible bases $\{\phi^l_m\}_{m\in\calM_l}$ mentioned in Section~\ref{sec:compatible_bases} corresponding to the representation $D^\ell$ satisfy
\begin{equation}\label{eq:equiv-1p-basis-O3}
    \phi^l_m(R\cdot\br) = \phi^l_m(sQ\cdot\br) = s^\ell\phi^l_m(Q\cdot\br)=s^\ell\sum_{\mu\in\calM_l}D^\ell_{\mu m}(Q)\phi^l_{\mu}(\br).
\end{equation}
Therefore, they are also compatible bases of the group $O(3)$ corresponding to the representation $D^{\ell,+}$ when $\ell$ is even, and $D^{\ell,-}$ when $\ell$ is odd. While other compatible bases may exist, we restrict ourselves to these specific bases. 

The GE functions in $\left[V^\bl\right]^{2L+1}$ and the GE-PI functions in $\left[\bV^\bl\right]^{2L+1}$ for the representation $D^{L,\pm}$ are strongly connected to the equivariant functions for the representation $D^{L}$. For notation simplicity, we denote by $V^{\bl,L,\pm}$ the space for GE functions in $\left[V^\bl\right]^{2L+1}$ corresponding to $D^{L,\pm}$. Likewise, we write $\bV^{\bl,L,\pm}$ for the spaces of the GE-PI functions in $\left[\bV^\bl\right]^{2L+1}$ for the representations $D^{L,\pm}$. 

\begin{proposition}\label{eq:equiv-basis-o3}
    Let $\bl\in\calL^N$, $L \in \N_0$. If $\sum\bl$ is even, then
    \[
        V^{\bl,L,+} = V^{\bl,L}, \quad \bV^{\bl,L,+} = \bV^{\bl,L}, 
    \]
    and
    \[  
        \dimge{\bl}{L,-} = \dimpi{\bl}{L,-}=0.
    \]
    Similarly, if $\sum\bl$ is odd, then
    \[
        V^{\bl,L,-} = V^{\bl,L}, \quad \bV^{\bl,L,-} = \bV^{\bl,L},
    \]
    and
    \[  
        \dimge{\bl}{L,+} = \dimpi{\bl}{L,+}=0.
    \]
\end{proposition}
\begin{proof}
We prove only for the GE case, and for the case where $\sum\bl$ is even. The proofs for all other situations are similar. 
On the one hand, for all $F\in V^{\bl,L,\pm}$, if we restrict the group action to elements of $SO(3)$ only, there holds 
\[
    \forall \br\in\Omega, \;  Q\in SO(3), \quad F(Q\cdot\br) = D^L(Q) F(\br),
\]
which indicates that $F\in V^{\bl,L}$. Hence, $V^{\bl,L,\pm}\subset V^{\bl,L}$.
On the other hand, for all $F\in V^{\bl,L}$, we have by definition that
\[
    \forall \br\in\Omega, \;  Q\in SO(3), \quad F(Q\cdot\br) = D^L(Q) F(\br).
\]
Furthermore, let $\{b^{\bl,L}_i\}_i$ be a basis of $V^{\bl,L}$ given in Proposition~\ref{prop:ge-su2}, then there exists $\mathbf{f}=\{f_i\}_i\subset\F$ such that 

\[
    F 
    = \sum_i f_i \; b^{\bl,L}_{i} 
    = \sum_i f_i \; \left[
        \sum_{\bbm \in \calM_{\bl,K}}c_{\bbm,i}
        \; \phi^\bl_\bbm
    \right]_{K\in\calM_L}.
\]
The last equality is due to~\eqref{eq:GE-basis-su2}. Using the equivariant property~\eqref{eq:equiv-1p-basis-O3} of the one-variable compatible basis $\phi^\bl_\bbm$, we have
\[
    \forall \br\in\Omega, \quad  F(-I\cdot\br) = (-1)^{\sum\bl} F(\br). 
\]
Thus, when $\sum\bl$ is even, we have immediately
\[
    \forall \br\in\Omega, \;  R=sQ\in O(3), \quad F(R\cdot\br) = s^{\sum\bl} D^L(Q) F(\br) = D^L(Q) F(\br) = D^{L,+}(R) F(\br).
\]
Therefore, $F\in V^{\bl,L,+}$, and thus $V^{\bl,L,+}=V^{\bl,L}$. In addition, $F$ can be in $V^{\bl,L,-}$ only when $\mathbf{f}=\mathbf{0}$. This completes the proof. 
\end{proof}

Proposition~\ref{eq:equiv-basis-o3} implies that the spaces of the GE and GE-PI functions for the $O(3)$ representations $D^{L,\pm}$, $V^{\bl,L,\pm}$ and $\bV^{\bl,L,\pm}$, are either spanned by basis functions given in Proposition~\ref{prop:ge-su2} and Theorem~\ref{thm:gepi-su2}, or are trivial zeros, depending on the parity of $\sum\bl$. 


\section{Numerical results}
\label{sec:num}

In this section, we numerically implement the method described in Section~\ref{sec:equivariant_lie_algebra} on the groups $SO(3)$ and $SU(2)$ to validate our theoretical claims in the preceding Section~\ref{sec:appl}, and to showcase the efficiency of the proposed method, particularly for the GE-PI case. 
For simplicity, we only present numerical results for the group $SO(3)$ as the results for $SU(2)$ are extremely similar.
To ensure reproducibility of the results and facilitate use of the proposed method, the source code implementing the methods described in this work, as well as the code that generates all the figures in this manuscript, have been archived on Zenodo~\cite{zhang2026zenodo}. In addition to the archived repository, the proposed method has also been integrated into \href{https://github.com/ACEsuit/EquivariantTensors.jl}{\emph{EquivariantTensors.jl}}~\cite{ortner2024et}, a Julia package providing the core computational kernels and infrastructure for building equivariant tensor layers.

\subsection{Efficiency}\label{sec:efficiency}


We first illustrate the efficiency gain in obtaining the equivariant bases using our method, which is mainly due to the chosen approach for dealing with the permutations, and the full exploitation of the specific structure of the matrix $\MlL$ defined in~~\eqref{eq:M_GE} and~\eqref{eq:M_GE_PI} or, for $SU(2)$ and $SO(3)$ in particular,~\eqref{eq:M2} and~\eqref{eq:M2-PI}. 

\subsubsection{Complexity of the algorithm}\label{sec:num-complexity}



To compute the coupling coefficients, three main steps have to be performed: (1) constructing the set $\calM_\bl$ (GE case) or the set of classes $\oMl$ (GE-PI case), (2) building the matrix $\Mup$, and (3) finding its kernel. In practice, it appears that the first step is typically negligible, as it amounts to less than 1\% of the overall cost. For instance, in the GE-PI case, when $\bl = (8,8,8,8,8,8)$, $L = 20$, generating the classes takes around 0.3 $ms$, constructing the matrix $\Mup$ takes 16 $ms$, and finding the kernel of $\Mup$ takes 45 $ms$.
We therefore focus below on the complexity of the last two steps only, and for simplicity only in the GE-PI case, the GE case being similar.


First, to build $\Mup$, one needs to loop over the row indices $\bbm\in\overline\calM_{\bl,K}$ for $K\in\calM_{L+1}\backslash\{-L-1,L\}$. Noting that the maximum number of non-zero elements in the row corresponding to $\bbm$ is $\card{\calN_\bbm}+1\le2\sum_{j=1}^\Nb\bl^{(j)}_1+1$, the whole cost of step (2) is bounded by
\begin{equation}\label{eq:cost2}
    \text{\texttt{cost}}_{\rm 2} \lesssim \left(2\sum_{j=1}^\Nb\bl^{(j)}_1+1\right)\sum_{K\in\calM_{L+1}\backslash\{L,L+1\}}\card{\overline{\calM}_{\bl,K}}.
\end{equation}

\begin{figure}[b!]
    \centering
    \includegraphics[width=.8\linewidth]{Figures/cost_vesus_classes_separated_L_full.pdf}
    \caption{Breakdown of construction time (ms) of the GE-PI bases versus the number of classes in the matrix $\Mup$~\eqref{eq:M2-PI}, or the product of the number of classes and the number of basis functions.  
    }
    \label{fig:cost_versus_classes_breakdown}
\end{figure}

The computational cost of step (3) depends on the structure of the matrix $\Mup$. As can be seen on Figure~\ref{fig:block_structure_of_M}(c), $\Mup$ is a supertriangular block matrix, that is, a block matrix with non-zero blocks only on the main diagonal and the super-subdiagonal. Moreover, the first $2L$ diagonal blocks are scaled identities, and the last block $B_L^+$ can in fact be made upper triangular, up to some proper convention of the ordering of its rows and columns. Therefore, the complexity of obtaining the kernel of $\Mup$, using a back-substitution approach, is bounded by
\begin{equation}\label{eq:cost3}
    \text{\texttt{cost}}_{3} \lesssim 
    \dimpi{\bl}{L}
    \left(2\sum_{j=1}^\Nb\bl^{(j)}_1+1\right)
    \sum_{K\in\calM_{L+1}\backslash\{L,L+1\}}\card{\overline{\calM}_{\bl,K}}.
\end{equation}
Consequently, the total cost to generate a basis of GE-PI functions, neglecting step (1), is bounded by
\begin{equation}\label{eq:cost-all}
    \text{\texttt{cost}}^{\rm GE-PI} \lesssim 
    \left(\dimpi{\bl}{L}+1\right)
    \left(2\sum_{j=1}^\Nb\bl^{(j)}_1+1\right)
    \sum_{K\in\calM_{L+1}\backslash\{L,L+1\}}\card{\overline{\calM}_{\bl,K}}.
\end{equation}



 We provide in Figure~\ref{fig:cost_versus_classes_breakdown}  a breakdown of the computational cost as a function of the number of classes (left) for a large variety of choices for parameters $\bl= (l,\ldots,l)$ of length $\Nat$ and equivariance order $L$.  On the right of Figure~\ref{fig:cost_versus_classes_breakdown},
we plot the cost of step (3) for the same parameters as a function of the product of the number of classes and the dimensionality, from which we observe a linear complexity, as expected from~\eqref{eq:cost2} and ~\eqref{eq:cost3}. Notice that for any given $\bl$, the computational cost of our method~\eqref{eq:cost-all} is asymptotically proportional to the size of the coupling coefficients $\bc^{\bl,L}$, and is hence asymptotically optimal.


\subsubsection{Numerical efficiency}



Our focus in this subsection lies in the comparison of the efficiency of constructing the GE-PI spaces $\bV^{\bl,L}$ for a given $\bl$ and for arbitrary $L\in\calLG$, as the computation of GE-PI bases is in general very expensive, and more challenging than GE bases, which can also be computed using generalized Clebsch--Gordan coefficients~\cite{Dusson2022-ie}. Also, while the proposed method and our code work for general $\bl\in\calL^\Nat$, it is impractical to carry out a thorough comparison for all possible $\bl$ due to the large freedom in the parameter choices (e.g. $\bl\in\calL^N$ can be chosen to have arbitrary length, and the elements therein vary freely either). Thus, we restrict ourselves to 
$\bl=(l,l,\ldots,l)\in\calL^\Nat$, especially noting that bases for general $\bl$ can be constructed from such identical $\bl$, using~Proposition~\ref{prop:GEPI_rec}. 
We could also use $\bl=((n,\ell),(n,\ell),\ldots,(n,\ell))\in\calL^\Nat$ with $n\in\N$ without changing the coupling coefficients, but only the underlying basis, as indicated in Remark~\ref{rem:id-ell-blocks}.
Like in some other contexts~\cite{Dusson2022-ie,Batatia2022-zi}, we sometimes refer $\Nat$ to the \emph{correlation order} and $l$ the \emph{polynomial degree} if $l\in\N_0$.



In the following, we benchmark our method against four well-known existing packages that perform similar calculations.
The first package is E3NN~\cite{geiger2022e3nn}, introduced in the context of machine-learning interatomic potentials (MLIPs), which provides a framework for generating equivariant neural networks. 
The second one is Lie-NN~\cite{Batatia2023-us},
which is theoretically presented for any compact Lie group, and tested in practice for a variety of  groups like $SO(3)$, $SU(2)$, and $SO^+(1,3)$. 
The third one is MACE~\cite{Batatia2022-zi},  a widely known foundation model being used to fit MLIPs with equivariant message passing graph neural networks.
Finally, Cuequivariance~\cite{geiger2024accelerate} is released by Nvidia  and provides CUDA-accelerated building blocks for equivariant neural networks.

\begin{figure}[b!]
    \centering
    \includegraphics[width=.89\linewidth]{Figures/cost_fix_l_log_N_4.pdf}
    \caption{Run time for constructing GE-PI bases using different packages for $l=1, 2, 3, 4$ and for correlation orders $N$'s from $1$ to $8$.}
    \label{fig:cost_fixed_l}
\end{figure}

In Figure~\ref{fig:cost_fixed_l}, we set the polynomial degrees $l$ to be different fixed values ($l=1,2,3,4$, respectively), and show the relationship between the run time measured in $ms$ required for constructing bases of $\bV^{\bl,L},$  for $ L\in\{0,1,\ldots,Nl\}$ and the correlation order $N$. Similarly, in Figure~\ref{fig:cost_fixed_N} we fix the correlation order ($N = 3,4,5,6$) and showcase how the run time varies while increasing the polynomial degree $l$. We truncate each method when the construction time for the GE-PI function spaces $\bV^{\bl,L}$ exceeds $10^6\; ms$. 



\begin{figure}[h!]
    \centering
    \includegraphics[width=.89\linewidth]{Figures/cost_fix_N_log_l_4.pdf}
    \caption{Run time for constructing the GE-PI bases using different packages for correlation orders $N = 3, 4, 5, 6$ and for polynomial degrees $l$'s from 1 to 8.}
    \label{fig:cost_fixed_N}
\end{figure}

Both Figure~\ref{fig:cost_fixed_l} and Figure~\ref{fig:cost_fixed_N} are presented on a log-log scale, from which we not only see that the construction cost of our method outperforms others in most of the cases, but also that in all the scenarios, our method appears to be 
of polynomial complexity, while other packages exhibit super-algebraic scaling.
In particular when $l=4$, none of the existing packages succeed in constructing the GE-PI basis for correlation order $N\ge 6$ within $10^6 \; ms$ while our method allows its construction in around $10^2 \; ms$. The advantages become more and more significant with the increase of either the correlation order $N$ or the polynomial degree~$l$.
We mention that all the packages we are comparing are written in Python, whereas ours is implemented using the Julia language. Despite the difference in programming languages, the qualitative scaling of the construction cost remains the same.


To see the scaling difference even clearer, we show in Figure~\ref{fig:cost_versus_classes} the construction time as either a function of the total number of basis functions $ \sum_{L=0}^{\sum\bl}\dimpi{\bl}{L}$ or that of the number of classes $\displaystyle\sum_{K\in\calM_{L+1}\setminus\{L,L+1\}}\card{\overline\calM_{\bl,K}}$. 
Let us mention that we compared the cost of constructing all the equivariant bases for a given $\bl\in\calL^N$ since the tested packages provide all of these equivariant bases simultaneously. In our implementation, each individual $L$-equivariant basis can be constructed independently, allowing us to investigate the cost of each basis individually, as presented on Figure~\ref{fig:cost_versus_classes_breakdown}.
A similar conclusion can be drawn from Figure~\ref{fig:cost_versus_classes}, that is, the cost of the proposed method scales polynomially with respect to the number of basis functions, or the number of classes, while the cost of the existing implementations seems to increase exponentially. 

\begin{figure}[t!]
    \centering
    \includegraphics[width=.88\linewidth]{Figures/cost_all_methods.pdf}
    \caption{Construction time (ms) of the GE-PI bases versus the number of basis functions, or the number of classes in the matrix $\Mup$~\eqref{eq:M2-PI}.
    }
    \label{fig:cost_versus_classes}
\end{figure}

In summary, by using the Lie algebra of the underlying rotation groups, and by exploring the simplified structure of the resulting matrix $\MlL$, we avoid the typical exponential complexity arising from permutations, and achieve a method for constructing GE-PI spaces at an almost linear cost. 

\subsubsection{Recursive construction}\label{sec:num-rec}


\begin{figure}[b!]
    \centering
    \includegraphics[width=.8\linewidth]{Figures/plot_for_recursion_l_8_3_3_latest.pdf}
    \caption{Comparison of recursive and direct construction time (\emph{ms}) for the GE-PI bases for $\boldsymbol{\ell} = (\ell,\ell,\ell,\ell,\ell,\ell)$ and $\bn = (1,1,1,2,2,2)$: (a) Runtime for constructing either all equivariant bases or a specific equivariant basis ($L=0$); (b) Cost breakdown by equivariance order $L$.
    }
    \label{fig:cost_recursive}
\end{figure}

We now evaluate and compare the computational efficiency of the direct and recursive constructions. Due to the large parameter choice in the recursive case, we restrict ourselves to the following setting: we consider $\bl$'s that can be divided into two non-intersecting blocks, each of length $N$, which gives a total of $2N$ variables. The $\ell$ values in both blocks are set to the same value. In particular, we take $N = 3$ with $\boldsymbol{\ell} = (\ell,\ell,\ell,\ell,\ell,\ell)$ and $\bn = (1,1,1,2,2,2)$, and analyze the construction time across varying polynomial degrees $\ell\in\{1,2,\ldots,8\}$ and orders of equivariance $L\in\{0,1,\ldots,N\ell\}$. The direct and the recursive construction time, measured in \emph{ms}, is illustrated in Figure~\ref{fig:cost_recursive}.
The solid lines in Figure~\ref{fig:cost_recursive}(a) represent the total computational cost of all equivariant bases for a given degree $l$, calculated using both the direct and recursive methods. The overall computational costs are comparable in magnitude in all cases, indicating that the recursive approach does not lead to a strict reduction in total complexity. However, a breakdown of the costs with respect to different $L$'s shows that they are distributed differently between the two methods. In particular, the recursive construction is substantially more efficient for small values of $L$, which correspond to the real implementation scenarios, as illustrated in Figure~\ref{fig:cost_recursive}(b). To emphasize this effect, we single out one of the most common cases, $L=0$, shown by dashed lines in Figure~\ref{fig:cost_recursive}(a), which highlights the significant efficiency gains that the recursive method can provide in practical applications.
A deeper comparison is left for future work.







\subsection{Dimensionality comparison}\label{sec:dimensionalities}
In this subsection, we investigate the dimension of GE and GE-PI spaces, or more precisely, 
the trend in the growth of the space dimensions composed of functions with different symmetries. 
Similarly to above, we focus on $\bl\in\calL^N$ that have identical components $l$, unless explicitly specified otherwise. In this setting, the dimensions $\dimge{\bl}{L}$ and $\dimpi{\bl}{L}$ are determined by three parameters: the polynomial degree $l$, the correlation order $N$ and the order of equivariance $L$. 



\paragraph{Linear scaling with respect to $L$}




Propositions~\ref{prop:asym-dim-ge} and Conjecture~\ref{conj:asym-dim-gepi-1block} suggest
that the dimensionalities of $V^{\bl,L}$ and $\bV^{\bl,L}$ scale linearly with respect to the order of equivariance $L$,
as observed on Figure~\ref{fig:Linear_scale_in_L}, 
when $L$ is not too large. In practice, this is indeed the regime of interest.
In the first two panels of Figure~\ref{fig:Linear_scale_in_L}, we fix the correlation order $N=8$, and show the change of GE and GE-PI dimensions with respect to $L$ for several different $l$, while in the last two panels, similar plots for fixed $l=20$ and a variety of $N$ are displayed. 

Owing to the linear scaling of the dimensions in $L$, we consider hereafter only  $L=0$, the invariant case, which is arguably the case of most interest. The numerical results we will present in the rest of the manuscript can naturally be extended to moderate $L$ because of the linear relation verified above. 

\begin{figure}[h!]
    \centering
    \includegraphics[width=.9\linewidth]{Figures/Linear_scaling_in_L.pdf}
    \caption{Relationship between the dimensions of the GE space $V^{\bl,L}$ and the GE-PI space $\bV^{\bl,L}$ and the order of equivariance $L$ with $\bl = (l,l,\ldots,l)\in\calL^N$ for varying polynomial degrees $l$ and correlation orders $N$. Both axes are displayed on a logarithmic scale. 
    }
    
    \label{fig:Linear_scale_in_L}
\end{figure}

\FloatBarrier

\paragraph{Asymptotic estimates}
Despite having explicit formulas for the GE and GE-PI dimensions (c.f. equations~\eqref{eq:dim_su2} and~\eqref{eq:dim_su2_pi}), they become increasingly costly to access as $\Nat$ increases. We rely on the estimations given in~\eqref{eq:dimge-est} and~\eqref{eq:dimpi-est} when $\Nat$ is unaffordably large. To this end, the estimated dimensions and the exact dimensions are compared in Figure~\ref{fig:dim-est}, in which we use solid circles to denote the exact dimensions $\dimge{\bl}{L}$ and $\dimpi{\bl}{L}$ and crosses to denote the estimated values $\dimgeest{\bl}{L}$ and $\dimpiest{\bl}{L}$. Furthermore, the normalized errors with respect to $\Nat$,
\[
  \text{\texttt{err}}_{\rm GE} :=  \frac{\dimge{\bl}{L} - \dimgeest{\bl}{L}}{\left(2l+1\right)^\Nat}, 
  \qquad
    \text{\texttt{err}}_{\rm GE-PI} :=   \frac{\dimpi{\bl}{L} - \dimpiest{\bl}{L}}{\binom{N+2l }{ 2l }},
\]
are presented in the right panels to show the accuracy of the estimations. 
As indicated by Figure~\ref{fig:dim-est}, the estimations align well with the exact dimensions, with the errors decaying at an algebraic rate of $N^{-5/2}$. We also observe a significant dimension reduction by considering PI on top of GE alone, which we will elaborate on in more detail later.

\begin{figure}[htb!]
    \centering
    \includegraphics[width=.9\linewidth]{Figures/dimensionality_estimation_in_N_l2_L0_5.pdf}
    \includegraphics[width=.9\linewidth]{Figures/dimensionality_estimation_in_N_l20_L0_5.pdf}
    \caption{Exact and estimated dimensions of the GE space $V^{\bl,0}$ and the GE-PI space $\bV^{\bl,0}$ for $\bl = (l,l,\ldots,l)\in\calL^N$ with different $l$ and $\Nat$. Axes are displayed on different scales, as indicated in the Figure.
    } 
    \label{fig:dim-est}
\end{figure}



Note that we have estimations in dimensionalities for general $\bl$ as well, as shown in Propositions~\ref{prop:asym-dim-ge} and~\ref{prop:asym-dim-gepi}. To demonstrate that the observed behaviour is not specific to the minimal setting, we present in Figure~\ref{fig:dim-est-Nb} the results for a set of more general $\bl = (\bl^{(1)},\bl^{(2)},\ldots,\bl^{(\Nb)})$, where $\bl^{(j)} = ((n_j,\ell),(n_j,\ell))\in\calL^2$, with $n_j$ all distinct, and $\ell=1,2$. 
The trend suggests a similar conclusion to that of before. Up to this stage, we have all the tools to fully investigate the dimensionalities in all possible cases. 

\begin{figure}[htb!]
    \centering
    \includegraphics[width=.9\linewidth]{Figures/dimensionality_estimation_in_Nb_N2_L0.pdf}
    \caption{Exact and estimated dimensions of the GE-PI space $\bV^{\bl,0}$ for $\bl = (\bl^{(1)},\bl^{(2)},\ldots,\bl^{(\Nb)})$, where $\bl^{(j)} = ((n_j,\ell),(n_j,\ell),\ldots,(n_j,\ell))\in\calL^N$ with $n_j$ distinct, $l=1,2$ and $N_j=2$. Axes are displayed on different scales, as indicated in the Figure.}
    \label{fig:dim-est-Nb}
\end{figure}



\paragraph{GE vs GE-PI}
Finally we compare dimensions of function spaces exhibiting different symmetries. For comparison, the dimension of $V^\bl$ itself is obviously $\prod_{i=1}^N (2l_i+1)$, and hence, 
\[
    \dim(\left[V^\bl\right]^{2L+1}) = (2L+1)\prod_{i=1}^N (2l_i+1).
\]
Besides, we see from Proposition~\ref{prop:basis-pi} that the dimensionality of the PI function space $\left[\bV^\bl\right]^{2L+1}$ is 
\[
    \dim(\left[\bV^\bl\right]^{2L+1}) = (2L+1)\prod_{j=1}^\Nb \binom{N_j+2\bl^{(j)}_1}{ 2\bl^{(j)}_1}.
\]
We are interested in comparing the two dimensions above with $\dimge{\bl}{L}$ and $\dimpi{\bl}{L}$, in both the asymptotic and pre-asymptotic regimes. Owing to the linear scaling in $L$, we fix $L=0$ as before. 
The dimensionalities of the spaces with different symmetries are shown in Figure~\ref{fig:dim_comparison},
for correlation order $N=6,\;20,\;40,\;$ and $80$ respectively, in four separate plots, in each of which the polynomial degree $l$ ranges from $4$ to $20$, covering both the asymptotic and pre-asymptotic regimes. For $N=80$ only, we use the estimation formulas~\eqref{eq:dimge-est} and~\eqref{eq:dimpi-est}. 

\begin{figure}[t!]
    \centering
    \includegraphics[width=.97\linewidth]{Figures/affect_of_sym_asym_preasym.pdf}
    \caption{Dimensionality of spaces of functions in $V^{\bl}$ having different symmetries with $\bl = (l,l,\ldots,l)\in\calL^N$: (a) $N=6$; (b) $N=20$; (c) $N=40$; and (d) $N=80$. For $N=80$, we use the estimated dimensions.}
    \label{fig:dim_comparison}
\end{figure}

As can be seen from Figure~\ref{fig:dim_comparison}, the space without any symmetry and that with both symmetries exhibit the largest and smallest dimensions, respectively. When the correlation order is small, GE turns out to be of equal importance, if not more, than PI for dimension reduction. However, as $\Nat$ increases, the PI dimension becomes increasingly closer to the GE-PI dimension, making PI the dominating source of dimension reduction. 
Tables exhibiting precise dimensionality examples are given in Appendix~\ref{app:table_of_dim}. 
These results suggest that it is in general important for dimensionality reduction to include both permutation and group-equivariance symmetry, starting from permutation invariance.



\section{Conclusion and outlook}
\label{sec:conclusion}

In this article, we have introduced a direct numerical method for efficiently constructing group-equivariant (GE) and permutation-invariant (PI) spaces, valid for arbitrary linear connected Lie groups. Furthermore we have presented an alternative recursive approach valid for compact groups. These methods are based on the use of Lie algebra to derive a linear system whose kernel elements are in one-to-one correspondence with either GE or GE-PI functions. Using the specific form of the Lie algebra for $SO(3)$ and $SU(2)$, the linear system is further simplified, and can be very efficiently solved exploiting sparsity.
Indeed, our method scales almost  linearly with respect to the number of basis functions, largely outperforming existing packages in the literature, which typically scale exponentially. 
On top of this, we manage to provide the explicit dimensionality of the spaces of interest. 
Comparing GE and GE-PI dimensionality shows how important it is in practice to consider the permutation-invariance to reduce the computational cost of approximating GE-PI functions.
It is interesting to note that the structure for the linear systems are very similar in the GE and GE-PI cases, showcasing that imposing PI on top of GE using our method does not require extra computational cost; indeed the GE case can be seen as a particular case of GE-PI, in the extreme scenario that compatible bases remove the permutation-invariance constraint. 

As an immediate application, this method can be used
as a GE-PI feature generator for a broad range of existing architectures in scientific machine learning, such as those for interatomic potentials in chemistry and materials science (see e.g. the equivariant neural networks E3NN~\cite{geiger2022e3nn}, NequIP~\cite{Batzner2022-az}, and MACE~\cite{Batatia2022-zi}). 
At this point we leave for future work the detailed derivation of the linear system structure for groups beyond $SU(2)$ and $SO(3)$, such as $SO^+(1,3)$ and $SU(n)$, but we expect a similar sparsity to arise.
Another natural question is whether this work can be extended beyond permutation invariance.
It is first interesting to note that our derivation for the GE-PI case relies on starting from a known PI basis, and that starting from a basis of permutation antisymmetric functions, e.g. Slater determinants, our strategy directly applies to obtain group-equivariant and permutation anti-invariant functions.
Second, combining the work~\cite{Olver2026-tt} obtaining irreducible representations of tensor representations of the symmetric group, which is similar to what is proposed in this paper but considering the symmetric group instead of a Lie group, 
we could in principle obtain group-equivariant and permutation-equivariant functions for more general representations of the symmetric group, which we also leave for future work.

\FloatBarrier
\section{Proofs}
\label{sec:proofs}

We gather in this section the main proofs of this article.

\subsection{Proofs of Section~\ref{sec:equivariant_lie_algebra}}\label{sec:proof_equivariant_lie_algebra}

\begin{proof}[\bf{Proof of Proposition~\ref{prop:GE_basis}}]
Let $\bl\in \calL^N, L\in \calL.$
We consider a function  $F^{\bl,L}\in [V^\bl]^{\card{\calM_L}}$ as 
\[ 
    F^{\bl,L} := \left[ F^{\bl,L}_{k} \right]_{k\in\calM_L}
     = \left[ \sum_{\bbm \in \calM_\bl} c_{\bbm,k}^{\bl,L} \; \phi^\bl_\bbm \right]_{k\in\calM_L}.
\]
Using~\eqref{eq:1b-basis}, the group action over $F^{\bl,L}$ reads as 
\[
   \forall g\in G,\; \bR\in \Omega^\Nat, \quad F^{\bl,L}(g \cdot \bR) =  \left[ \sum_{\bbm,\bbm' \in \calM_\bl} c_{\bbm,k}^{\bl,L} \rep{\bl}{\bbm,\bbm'}(g) \phi^\bl_{\bbm'}(\bR) \right]_{k\in\calM_L},
\]
where $\displaystyle\rep{\bl}{\bbm,\bbm'}(g) = \prod_{i=1}^\Nat \rep{l_i}{m_i,m'_i}(g)$. Also, by definition, $F^{\bl,L}$ is equivariant with respect to the representation $\rep{L}{}$ if and only if
\[
    \forall g\in G,\; \bR\in\Omega^\Nat, k\in \calM_L, 
    \quad F^{\bl,L}_{k}(g \cdot \bR) = \sum_{k'\in\calM_L} \rep{L}{k,k'}(g)F_{k'}^{\bl,L}(\bR),
\]
that is,
\[ 
\forall g\in G,\; \bR\in\Omega^\Nat,  k\in\calM_L,
\quad
   \sum_{k'\in\calM_L} \sum_{\bbm \in \calM_\bl} c_{\bbm,k'}^{\bl,L} \rep{L}{k,k'}(g) \phi^\bl_\bbm(\bR)
   = 
   \sum_{\bbm,\bbm' \in \calM_\bl} c_{\bbm,k}^{\bl,L} \rep{\bl}{\bbm,\bbm'}(g) \phi^\bl_{\bbm'}(\bR).
\]
Switching $\bbm$ and $\bbm'$ in the right hand side and noting that $(\phi^\bl_\bbm)_{\bbm\in \calM_\bl}$ is a basis of $V^\bl$, we obtain that 
$F^{\bl,L}$ is GE with respect to the representation $\rep{L}{}$ if and only if
\begin{equation}\label{eq:GE-coeff}
\forall \bbm \in \calM_\bl,\; g\in G,\; k\in\calM_L, \qquad
     \sum_{k'\in\calM_L}c_{\bbm,k'}^{\bl,L} 
     \rep{L}{k,k'}(g)  
   = \sum_{\bbm' \in \calM_\bl} c_{\bbm',k}^{\bl,L} \rep{\bl}{\bbm',\bbm}(g). 
\end{equation}
By viewing $\{c_{\bbm,k}^{\bl,L}\}_{\bbm\in\mathcal{M}_\bl,k\in\mathcal{M}_L}$ as a matrix $C^{\bl,L}$ of size $\card{\mathcal{M}_\bl}\times\card{\mathcal{M}_L}$,~\eqref{eq:GE-coeff} is equivalent to
\begin{equation}
\label{eq:matrix_GE_sylvester}
        \forall g\in G, \quad \rep{L}{}(g)\left(C^{\bl,L}\right)^T = \left(C^{\bl,L}\right)^T\rep{\bl}{}(g).
\end{equation}
Taking the derivative of the above equation at the neutral element $e\in G$, we have
\begin{equation}\label{eq:diff-rep-GE}
    \forall X\in\mathfrak g, \quad d\rep{L}{}(X)\left(C^{\bl,L}\right)^T = \left(C^{\bl,L}\right)^Td\rep{\bl}{}(X).
\end{equation}
Since $d\rep{L}{}$ is linear,~\eqref{eq:diff-rep-GE} holds true if and only if
\begin{equation}\label{eq:diff-rep-GE-on-generator}
    \forall\dd=1,\ldots,N_{\rm{dim}},\quad d\rep{L}{}(X_\dd)\left(C^{\bl,L}\right)^T = \left(C^{\bl,L}\right)^Td\rep{\bl}{}(X_\dd).
\end{equation}
Recalling that $d\rep{L}{}(X_\dd) = \vrho{L,\dd}{}$ and noticing that for $\bbm',\bbm\in\calM_{\bl}$, 
\begin{equation}\label{eq:drho}
    d\rep{\bl}{}(X_\dd)_{\bbm',\bbm} = \vrho{\bl,\dd}{\bbm',\bbm}= \sum_{j=1}^\Nat \vrho{l_j,\dd}{m_j',m_j}\prod_{s=1,s\ne j}^\Nat \rep{l_s}{m'_s,m_s}(e)= 
    \sum_{j=1}^\Nat \vrho{l_j,\dd}{m'_j,m_j}\prod_{\substack{s=1 \\ s\ne j }}^N \delta_{m'_s,m_s},
\end{equation}
where $\delta_{ij}$ denotes the Kronecker delta,
we obtain the component-wise form of~\eqref{eq:diff-rep-GE-on-generator} as
\[
    \forall \dd=1,\ldots,N_{\rm dim}, k\in\calM_L, \bbm \in \calM_\bl, \quad
    \sum_{k'\in\calM_L}c_{\bbm,k'}^{\bl,L} \vrho{L,\dd}{k,k'}  
    = \sum_{\bbm' \in \calM_\bl} c_{\bbm',k}^{\bl,L}
    \sum_{j=1}^\Nat \left( \vrho{l_j,\dd}{m'_j,m_j}\prod_{\substack{s=1 \\s\ne j}}^\Nat \delta_{m'_s,m_s} \right),  
\]
or equivalently, for all $\dd=1,\ldots,N_{\rm dim},\; k\in\calM_L,\; \bbm \in \calM_\bl$, 
\begin{align*}
    \sum_{\bbm' \in \calM_\bl} \sum_{k'\in\calM_L}c_{\bbm,k'}^{\bl,L} \vrho{L,\dd}{k,k'}  \delta_{\bbm,\bbm'}
    & - \sum_{\bbm' \in \calM_\bl} \sum_{k'\in\calM_L} c_{\bbm',k}^{\bl,L} \delta_{k,k'}
    \sum_{j=1}^\Nat \left( \vrho{l_j,\dd}{m'_j,m_j}\prod_{\substack{s=1 \\s\ne j}}^\Nat \delta_{m'_s,m_s} \right) = 0,  
\end{align*}
which means that the vector of coupling coefficients $\bc^{\bl,L}:=\left(c_{\bbm,k}^{\bl,L}\right)_{\bbm\in\mathcal{M}_\bl,k\in\mathcal{M}_L}$ belongs to $\ker(M^{\bl,L})$. 

Reciprocally, if $\bc^{\bl,L} \in\ker(M^{\bl,L})$, equation~\eqref{eq:diff-rep-GE} is satisfied.
Multiplying both sides of the latter by $d\rep{L}{}(X)$ from the left and using~\eqref{eq:diff-rep-GE} itself  yields
\[
    \forall X\in\mathfrak g, \quad \left(d\rep{L}{}(X)\right)^2\left(C^{\bl,L}\right)^T = d\rep{L}{}(X)\left(C^{\bl,L}\right)^Td\rep{\bl}{}(X) = \left(C^{\bl,L}\right)^T\left(d\rep{\bl}{}(X)\right)^2.
\]
Iterating this argument, we see that
\[
    \forall X\in\mathfrak g,\; p\in\N_0, \quad \left(d\rep{L}{}(X)\right)^p\left(C^{\bl,L}\right)^T = \left(C^{\bl,L}\right)^T\left(d\rep{\bl}{}(X)\right)^p,
\]
and hence, using~\cite[(3.11)]{Hall2013-af},
\begin{equation}\label{eq:recip-cond}
    \forall X\in\mathfrak g,\quad \rep{L}{}\left(\exp(X)\right)\left(C^{\bl,L}\right)^T = \left(C^{\bl,L}\right)^T\rep{\bl}{}\left(\exp(X)\right).
\end{equation}
Finally, since any element $g\in G$ in the connected linear Lie group $G$ can be expressed as a finite product of exponentials, we easily recover~\eqref{eq:matrix_GE_sylvester} from~\eqref{eq:recip-cond}, using~\eqref{eq:diff-rep-GE} and that $\rep{\bl}{}$ and  $\rep{L}{}$ are representations.
The basis and dimensionality of $V^{\bl,L}$ are straightforward consequences.
\end{proof}

\begin{proof}[\bf{Proof of Proposition~\ref{prop:basis-pi}}]
    Since $\left\{ \phi^\bl_\bbm, \bbm\in \calM_\bl \right\}$  is 
    a basis of $V^{\bl}$, a spanning set for PI functions in $V^{\bl}$ is 
\[
    \left\{ \displaystyle\bpsi{\bl}{\bbm}: =\sum_{\sigma\in S_\Nat}\phi^\bl_\bbm\circ\sigma, \qquad \bbm\in \calM_\bl \right\}.
\]
Noting that $\bpsi{\bl}{\bbm}=\bpsi{\bl}{\bbm'}$ if there exists $\pi\in \calS_{\bl},$ such that $\bbm^{\prime}=\pi\bbm$, a spanning set 
of PI functions in $V^\bl$ is
$\left\{ \displaystyle\bpsi{\bl}{\bbm}, \bbm\in \oMl \right\}$.
Moreover, for $\bbm,\bbm'\in \oMl$,
\begin{align*}
    \left\langle\bpsi{\bl}{\bbm},\bpsi{\bl}{\bbm'}\right\rangle =& 
    \Nat!\left\langle\phi^\bl_\bbm, \sum_{\sigma\in S_\Nat}\phi^\bl_{\bbm'}\circ\sigma\right\rangle.
\end{align*}

Noting that permuting the variables corresponds to an inverse permutation on the basis functions indices, and using~\eqref{eq:orthogonality}, we obtain
\begin{align*}
\left\langle\bpsi{\bl}{\bbm},\bpsi{\bl}{\bbm'}\right\rangle
    = & \Nat !\sum_{\sigma\in S_\Nat}\left\langle\phi^\bl_\bbm, \phi^{\sigma\bl}_{\sigma\bbm'}\right\rangle 
    =  \Nat !\sum_{\sigma\in S_\Nat}\delta_{\bl,\sigma\bl}\delta_{\bbm,\sigma\bbm'} 
    =  \Nat !\sum_{\sigma\in \calS_\bl}\delta_{\bbm,\sigma\bbm'}. 
\end{align*}
Using the definition of $\oMl$, we easily obtain that 
$\left\langle\bpsi{\bl}{\bbm},\bpsi{\bl}{\bbm'}\right\rangle = 0$ if and only if $\bbm \neq \bbm'$, that is $\left\{ \displaystyle\bpsi{\bl}{\bbm},\; \bbm\in \oMl \right\}$ is an orthogonal basis of PI functions in $V^\bl$ and the dimension of the space of PI functions in $V^\bl$ is the cardinality of $\oMl$.
\end{proof}


\begin{proof}[\bf{Proof of Proposition~\ref{prop:basis-pi2}}]
First, since the functions $\bphi{\bl}{\bbm}$ are identical for all $\bbm$ in a class $\overline{\bbm} \in\oMl$
    the functions $\bphi{\bl}{\bbm},\;\bbm\in\oMl$ are clearly a spanning set of $\bV^{\bl}$.
We now need to show that the basis functions $\left\{ \bphi{\bl}{\bbm},\;\bbm\in\oMl\right\}$ are linearly independent.
Assume that
there exists coefficients $\bc \in \K^{|\oMl|}$ such that 
$\displaystyle\sum_{\bbm \in\oMl}c_\bbm\bphi{\bl}{\bbm}=0$. This means that we have for any $\br_1,\ldots,\br_N\in \Omega$,
\begin{equation}
    \label{eq:linearindependence}
    \displaystyle\sum_{\bbm \in\oMl}c_\bbm\bphi{\bl}{\bbm}(\br_1,\ldots,\br_\Nat) = 0.
\end{equation}
We are going to show by induction on $d \in \N$ that for any $\bR = (\br_1,\ldots,\br_\Nat)\in\Omega^\Nat$ and for all $1\le n_1\le\ldots\le n_\Nat\le \Nat$ such that $\card{\{n_1,\ldots,n_\Nat \}  }=d$, there holds
\begin{equation}
        \displaystyle\sum_{\bbm \in\oMl}c_\bbm \sum_{\bm{\alpha}\in \calS_{\{n_1,\ldots,n_\Nat\}}}\phi^\bl_\bbm
    (\br_{\alpha_1}, \ldots, \br_{\alpha_\Nat})=0,
    \label{eq:induction_step}
\end{equation}
where $\calS_{\{n_1,\ldots,n_\Nat\}}=\{({\sigma(n_1)},\ldots,{\sigma(n_\Nat)}),\;\sigma\in S_\Nat\}$ is the set of all different tuples $(n_1,\ldots,n_\Nat)$ in any order.

First, for $d=1$, 
for 
$\bR=(\br,\dots,\br)$, with $\br\in\Omega$,
\begin{align*}
    \displaystyle\sum_{\bbm \in\oMl}c_\bbm\bphi{\bl}{\bbm}(\br,\ldots,\br)
    &=\displaystyle\sum_{\bbm \in\oMl}c_\bbm\prod_{i=1}^\Nat\left(\sum_{j=1}^{\Nat}\phi^{l_i}_{m_i}\left(\br\right)\right)\\
    &=\displaystyle\sum_{\bbm \in\oMl}c_\bbm\prod_{i=1}^\Nat\left(\Nat\phi^{l_i}_{m_i}\left(\br\right)\right)\\
    &=\Nat^\Nat\displaystyle\sum_{\bbm \in\oMl}c_\bbm \phi^\bl_\bbm(\br,\ldots,\br).
\end{align*}
Using~\eqref{eq:linearindependence}, we obtain
\[
    \sum_{\bbm \in\oMl}c_\bbm \phi^\bl_\bbm(\br,\ldots,\br) = 0.
\]
Since $\br$ is arbitrary in $\Omega$, for $1\le n_1=\ldots=n_\Nat=n\le\Nat$,

\[
    \displaystyle\sum_{\bbm \in\oMl}c_\bbm \sum_{\bm{\alpha}\in \calS_{\{n_1,\ldots,n_\Nat\}}}\phi^\bl_\bbm
    (\br_{\alpha_1}, \ldots, \br_{\alpha_\Nat})=\displaystyle\sum_{\bbm \in\oMl}c_\bbm \phi^\bl_\bbm(\br_{n},\ldots,\br_{n})
    =0,
\]
which proves the result for $d=1$.

Now assume that~\eqref{eq:induction_step} holds for all $\bR\in\Omega^\Nat$ and all ordered $\bn\in\{1,\ldots,\Nat\}^\Nat$ having at most $d-1$ distinct values.
Let $\bk = \{k_1,\ldots,k_\Nat\}$ with $1\le k_1\le\ldots\le k_{\Nat}\le d$ such that $\card{\bk}=d$. 
For $\br_1,\ldots, \br_{d} \in \Omega$, denote $\bR_{\bk}=(\br_{k_1},\ldots,\br_{k_{\Nat}})$. Writing explicitly the repetition of variables
\[
  \bphi{\bl}{\bbm}(\bR_{\bk})=\prod_{i=1}^\Nat\left(\sum_{j=1}^{\Nat}\phi^{l_i}_{m_i}\left(\br_{k_j}\right)\right)=\prod_{i=1}^\Nat\left(\sum_{n=1}^{d}\blamb{}{\bk,n}\phi^{l_i}_{m_i}\left(\br_{n}\right)\right),
\]
where $\blamb{}{\bk} \in \N^{d}$ denotes the count vector of $\bk$ for the dictionary $\{1,\ldots,d\}$ (see~\cite[Supplementary lecture H]{Kozen2007-ot}).  
Switching the product and the sum, we obtain 
\begin{align*}
    \bphi{\bl}{\bbm}(\bR_{\bk})&=
    \displaystyle\sum_{1\le n_1,\ldots,n_\Nat\le d}\left(\prod_{i=1}^\Nat \blamb{}{\bk,n_i}\right)\left(\prod_{i=1}^\Nat\phi^{l_i}_{m_i}\left(\br_{{n_i}}\right)\right)\\
    &=
    \displaystyle\sum_{1\le n_1,\ldots,n_\Nat\le d}
    \left(\prod_{i=1}^{d} (\blamb{}{\bk,i})^{\blamb{}{\bn,i}}\right)
    \phi^\bl_\bbm\left(\br_{n_{1}},\ldots,\br_{{n_\Nat}}\right).
\end{align*}
The count vectors $\blamb{}{\bn}$ being independent of the order of the elements in $\bn$, 
we can restrict the first sum on ordered elements, and separate the sum depending on the number of different indices in $\bn$, that is,
\begin{align*}
      \bphi{\bl}{\bbm}(\bR_{\bk})&=
      \displaystyle\sum_{1\le n_1\le\ldots\le n_\Nat\le d}\left(\prod_{i=1}^{d} (\blamb{}{\bk,i})^{\blamb{}{\bn,i}}\right)\sum_{\bm{\alpha}\in \calS_{\{n_1,\ldots,n_\Nat\}}}\phi^\bl_\bbm
    (\br_{\alpha_1}, \ldots, \br_{\alpha_\Nat}) \\
    &= \displaystyle
    \sum_{d' = 1}^d \;
    \sum_{\substack{1\le n_1\le\ldots\le n_\Nat\le d\\\card{\bn}=d'}}\left(\prod_{i=1}^{d} (\blamb{}{\bk,i})^{\blamb{}{\bn,i}}\right)
    \sum_{\bm{\alpha}\in \calS_{\{n_1,\ldots,n_\Nat\}}}\phi^\bl_\bbm
    (\br_{\alpha_1}, \ldots, \br_{\alpha_\Nat}).
\end{align*}
Taking the linear combination of functions $\bphi{\bl}{\bbm}$  and using the induction step~\eqref{eq:induction_step}, the terms involving less than $d$ different values of $\br$ vanish,
\begin{align*}
    \displaystyle\sum_{\bbm \in\oMl}c_\bbm\bphi{\bl}{\bbm}(\bR_{\bk})
    &=\displaystyle\sum_{\substack{1\le n_1\le\ldots\le n_\Nat\le d\\\card{\bn}=d}}\left(\prod_{i=1}^{d} (\blamb{}{\bk,i})^{\blamb{}{\bn,i}}\right)\displaystyle\sum_{\bbm \in\oMl}c_\bbm\sum_{\bm{\alpha}\in \calS_{\{n_1,\ldots,n_\Nat\}}}\phi^\bl_\bbm
    (\br_{\alpha_1}, \ldots, \br_{\alpha_\Nat}) = 0,
\end{align*}
using~\eqref{eq:linearindependence}.
Noting that the Bernstein--Vandermonde matrix $\displaystyle\left(\prod_{i=1}^{d} (\blamb{}{\bk,i})^{\blamb{}{\bn,i}}\right)_{\substack{\text{ordered } \bn,\bk\in\{1,\ldots,\Nat\}^\Nat\\\card{\bk}=\card{\bn}=d}}$ is invertible~\cite[(2.10)]{Ainsworth2011-pk}, 
we obtain multiplying by the inverse of the matrix that, for any ordered $\bn\in\{1,\ldots,\Nat\}^\Nat$, with $\card{\bn}=d$ 
\[
    \displaystyle\sum_{\substack{1\le k_1\le\ldots\le k_\Nat\le d\\\card{\bk}=d}}\displaystyle\sum_{\bbm \in\oMl}
    c_\bbm\sum_{\bm{\alpha}\in \calS_{\{n_1,\ldots,n_\Nat\}}}\phi^\bl_\bbm
    (\br_{\alpha_1}, \ldots, \br_{\alpha_\Nat})=0.
\]
Since the terms inside the sum are independent of $\bk$, we can divide by the number of ordered $\bk\in\{1,\ldots,\Nat\}^\Nat$ having $d$ different coordinates, we finally obtain that for any $1\le n_1\le\ldots\le n_\Nat\le d$ such that $\card{\{n_1,\ldots,n_\Nat \}  }=d$ and any $\br_1,\ldots,\br_d \in \Omega$
\[
  \displaystyle\sum_{\bbm \in\oMl}c_\bbm\sum_{\bm{\alpha}\in \calS_{\{n_1,\ldots,n_\Nat\}}}\phi^\bl_\bbm
    (\br_{\alpha_1}, \ldots, \br_{\alpha_\Nat})=0,
\]
which concludes the induction, using that the labeling of the variables is arbitrary.


In particular for $d=\Nat$ and $\bn=(1,\ldots,\Nat)$, we have

\[
    \displaystyle\sum_{\bbm \in\oMl}c_\bbm  \displaystyle\sum_{\sigma\in S_\Nat}\phi^\bl_\bbm\left(\br_{\sigma_1},\ldots,\br_{\sigma_\Nat}\right)
    = \displaystyle\sum_{\bbm \in\oMl}c_\bbm \bpsi{\bl}{\bbm}(\br_{1},\ldots,\br_{\Nat})=0.
\]
From the linear independence of the $\bpsi{\bl}{\bbm}$, we conclude that the $\bphi{\bl}{\bbm}$ are linearly independent.
\end{proof}

In order to prove Theorem~\ref{thm:GEPI_basis}, we use $\obm$ to denote elements in $\oMl$, to distinguish it from $\bbm$, the elements in $\calM_\bl$. In addition, we need a preliminary lemma. For $\bl\in\calL^\Nat$, define two matrices $S^\bl\in\F^{\card{\oMl}\times\card{\calM_{\bl}}}$ and $T^\bl\in\F^{\card{\calM_{\bl}}\times\card{\oMl}}$ whose elements read, for $\obm\in\oMl$ and $\bbm'\in\calM_\bl$, as
    \[
        S^\bl_{\obm,\bbm'} = \begin{cases}
            1, & \quad\textup{if } \; \obm'= \obm \; \textup{ and }\;  \bbm' \; \textup{ is ordered}, \\
            0, & \quad \textup{otherwise,}
        \end{cases} \quad\quad \textup{and} \quad\quad  T^\bl_{\bbm',\obm} = \begin{cases}
            1, & \quad\textup{if }\;  \obm=\obm', \\
            0, & \quad \textup{otherwise.}
        \end{cases}
    \] 
Note that multiplying by $S^\bl$ on the left selects elements indexed by the ordered class representatives and multiplying by $T^\bl$ on the right sums over the elements of the same class.
\begin{lemma}
    There hold
    \begin{equation}\label{eq:TS_lemma}
        \forall g\in G, \quad T^{\bl}S^{\bl}\rep{\bl}{}(g)T^{\bl}=\rep{\bl}{}(g)T^{\bl},
    \end{equation}
    and
    \begin{equation}\label{eq:TS_lemma_lie_alg}
       \forall X\in \mathfrak g,\quad  T^{\bl}S^{\bl} d\rep{\bl}{}(X)T^{\bl}=d\rep{\bl}{}(X)T^{\bl}.
    \end{equation}
\end{lemma}

\begin{proof}
    Let $g\in G$. From the definition of $T^\bl$ and $S^\bl$, there holds, for $\bbm\in\calM_\bl$ and $\obm'\in\oMl$,
    \[
        \left[\rep{\bl}{}(g)T^\bl\right]_{\bbm,\obm'}=\displaystyle\sum_{\bmu\in\calM_\bl}\rep{\bl}{\bbm,\bmu}(g)T^\bl_{\bmu,\obm'}= \sum_{\bmu\in\obm'}\rep{\bl}{\bbm,\bmu}(g),
    \]
    and for $\bbm,\; \bmu\in\calM_\bl$
    \[
        \left[T^\bl S^\bl\right]_{\bbm,\bmu}= \displaystyle\sum_{\obm'\in\oMl}T^\bl_{\bbm,\obm'} S^\bl_{\obm',\bmu}=\begin{cases}
            1, & \quad\textup{if } \bar\bmu=\obm \textup{ and } \bmu \textup{ is ordered} , \\
            0, & \quad \textup{otherwise}.
        \end{cases} 
    \]
    Thus,
    \begin{equation*}
        \left[T^\bl S^\bl\rep{\bl}{}(g)T^\bl\right]_{\bbm,\obm'} = \displaystyle\sum_{\bmu\in\calM_\bl}\left[T^\bl S^\bl\right]_{\bbm,\bmu}\left[\rep{\bl}{}(g)T^\bl\right]_{\bmu,\obm'} = \sum_{\bmu\in\obm'}\rep{\bl}{\widetilde\bbm,\bmu}(g),
    \end{equation*}
    where $\widetilde\bbm$ is the ordered representative of $\obm$.
    By definition of equivalent classes, there exists ${\pi}\in\calS_\bl$ such that ${\pi}\bbm=\widetilde{\bbm}$. Since ${\pi}\bl=\bl$ for any ${\pi}\in\calS_\bl$, there holds
    \begin{equation*}
        \rep{\bl}{\widetilde\bbm,\bmu}(g) = \rep{\bl}{{\pi}\bbm,\bmu}(g) = \displaystyle\prod_{i=1}^\Nat \rep{l_i}{m_{{\pi}_i},\mu_i}(g) = \displaystyle\prod_{i=1}^\Nat \rep{l_i}{m_i,\mu_{({\pi}^{-1})_i}}(g) = \rep{\bl}{\bbm,{\pi}^{-1}\bmu}(g).
    \end{equation*}
    Hence, since ${\pi}^{-1}$ acts as a bijection on $\calM_\bl$,
    \[
        \sum_{\bmu\in\calM_\bl}\rep{\bl}{\widetilde\bbm,\bmu}(g) = \sum_{\bmu\in\calM_\bl}\rep{\bl}{\bbm,{\pi}^{-1}\bmu}(g) = \sum_{\bmu\in\calM_\bl}\rep{\bl}{\bbm,\bmu}(g),
    \]
    and thus 
    \[
        \left[T^\bl S^\bl\rep{\bl}{}(g)T^\bl\right]_{\bbm,\obm'} = \left[\rep{\bl}{}(g)T^\bl\right]_{\bbm,\obm'},
    \]
    which is exactly~\eqref{eq:TS_lemma} element-wise. Equation~\eqref{eq:TS_lemma_lie_alg} is a direct consequence of differentiating~\eqref{eq:TS_lemma}.
\end{proof}


\begin{proof}[\bf{Proof of Theorem~\ref{thm:GEPI_basis}}]
   Let $\bl\in\calL^N$, $L\in \calL$.
   We consider a function  $F^{\bl,L}\in [\bV^\bl]^{\card{\calM_L}}$ as 
   \begin{equation*}
      F^{\bl,L} := 
    \left[ F^{\bl,L}_{k} \right]_{k\in\calM_L}
    = 
    \left[ \sum_{\obm \in \oMl} c_{\obm,k}^{\bl,L} \; \bphi{\bl}{\obm} \right]_{k\in\calM_L},
\end{equation*}
which is already PI, and where the coupling coefficients are in $\K$.
   Using~\eqref{eq:1b-basis}, the group action over $F^{\bl,L}$ reads as 
\[
   \forall g\in G,\; \bR\in \Omega^\Nat, \quad F^{\bl,L}(g \cdot \bR) =  \left[ \sum_{\obm \in \oMl} c_{\obm,k}^{\bl,L} \sum_{\bbm' \in \calM_\bl} \rep{\bl}{\bbm,\bbm'}(g) \bphi{\bl}{\obm'}(\bR) \right]_{k\in\calM_L}.
\]
Noting that the sum over $\bbm' \in \calM_\bl$ can be decomposed as a sum over $\obm'$ classes and possible permutations, we obtain
\begin{equation*}
\begin{split}
    F^{\bl,L}(g \cdot \bR) & 
     =  \left[ \sum_{\obm,\obm' \in \oMl} c_{\obm,k}^{\bl,L} \left( \sum_{\bmu\in\obm'}\rep{\bl}{\bbm,\bmu}(g) \right) \bphi{\bl}{\obm'}(\bR) \right]_{k\in\calM_L}.
\end{split}
\end{equation*}
Also by definition, $F^{\bl,L}$ is equivariant with respect to the representation $\rep{L}{}$ if and only if
\[
    \forall g\in G, \bR\in\Omega^\Nat, k\in \calM_L, 
    \quad F^{\bl,L}_{k}(g \cdot \bR) = \sum_{k'\in\calM_L} \rep{L}{k,k'}(g)F_{k'}^{\bl,L}(\bR),
\]
that is, for any $g\in G, \bR\in\Omega^\Nat, k\in\calM_L,$
\[ 
   \sum_{k'\in\calM_L} \sum_{\obm \in \oMl} c_{\obm,k'}^{\bl,L} \rep{L}{k,k'}(g) \bphi{\bl}{\obm}(\bR)
   = 
   \sum_{\obm,\obm' \in \oMl} c_{\obm,k}^{\bl,L} 
   \left( \sum_{\bmu\in\obm'}\rep{\bl}{\bbm,\bmu}(g) \right)
   \bphi{\bl}{\obm'}(\bR).
\]
Switching $\obm$ and $\obm'$ in the right hand side and noting that $(\bphi{\bl}{\obm})_{\obm\in \oMl}$ is a basis of PI functions in $\bV^\bl$, we obtain that $F^{\bl,L}$ is GE-PI with respect to the representation $\rep{L}{}$ if and only if
\begin{equation}\label{eq:GEPI-coeff}
\forall \obm \in \oMl, g\in G, k\in\calM_L, \qquad
     \sum_{k'\in\calM_L}c_{\obm,k'}^{\bl,L} 
     \rep{L}{k,k'}(g)  
   = \sum_{\obm' \in \oMl} c_{\obm',k}^{\bl,L} 
   \left( \sum_{\bmu\in\obm}\rep{\bl}{\bbm',\bmu}(g) \right). 
\end{equation}
Viewing $\{c_{\obm,k}^{\bl,L}\}_{\obm\in\oMl,k\in\mathcal{M}_L}$ as a matrix $C^{\bl,L}$ of size $\card{\oMl}\times\card{\mathcal{M}_L}$,
\eqref{eq:GEPI-coeff} is equivalent to
\begin{equation}\label{eq:matrix_GE-PI_sylvester}
    \forall g\in G, \quad \rep{L}{}(g)\left(C^{\bl,L}\right)^T = \left(C^{\bl,L}\right)^TS^\bl\rep{\bl}{}(g)T^\bl.
\end{equation}  
Taking the derivative of the above equation at the neutral element $e\in G$, we have
\begin{equation}\label{eq:diff-rep-GEPI}
    \forall X\in\mathfrak g, \quad d\rep{L}{}(X)\left(C^{\bl,L}\right)^T = \left(C^{\bl,L}\right)^TS^\bl d\rep{\bl}{}(X)T^\bl.
\end{equation}
Since $d\rep{L}{}$ is linear,~\eqref{eq:diff-rep-GEPI} holds true if and only if
\begin{equation}\label{eq:diff-rep-GEPI-on-generator}
    \forall\dd=1,\ldots,N_{\rm dim},\quad d\rep{L}{}(X_\dd)\left(C^{\bl,L}\right)^T = \left(C^{\bl,L}\right)^TS^\bl d\rep{\bl}{}(X_\dd)T^\bl.
\end{equation}
As in the GE case, we obtain the component-wise form of~\eqref{eq:diff-rep-GEPI-on-generator} using~\eqref{eq:drho}, that is
\begin{equation}
\label{eq:system_GEPI}
    \forall \dd =1,\ldots,N_{\rm{dim}}, k\in\calM_L, \obm \in \oMl, \;
    \sum_{k'\in\calM_L}c_{\obm,k'}^{\bl,L} \vrho{L,\dd}{k,k'}  
    = \sum_{\obm' \in \oMl} c_{\obm',k}^{\bl,L}
    \sum_{\bmu\in\obm}
    \sum_{j=1}^\Nat \left( \vrho{l_j,\dd}{m'_j,\mu_j}\prod_{\substack{s=1 \\s\ne j}}^\Nat \delta_{m'_s,\mu_s} \right).
\end{equation}
The above terms in the right-hand side can only be non-zero if $\bmu$ and $\bbm'$ differ by at most one element.
Once the class $\obm$ is given, the classes $\obm'\in \oMl$ that can contribute can be characterized using the count vectors. Indeed, changing one element in a vector $\bbm$ corresponds to modifying the count vector $\blamb{}{\bbm}$ as follows: one element in $\blamb{}{\bbm}$ is decreased by one, another is increased by one, and this can only affect a single block of indices corresponding to identical $\bl$'s at a time. 
Hence, the set of classes possibly leading to non-zero terms in the sum over $\obm'$ in~\eqref{eq:system_GEPI}, apart from $\obm$ itself, belongs to the set
$\calN_\obm$ defined in~\eqref{eq:interacting_classes}.
Therefore,
 \begin{align*}
     \forall \dd =1,\ldots,N_{\rm{dim}}, k\in\calM_L, \obm \in \oMl,
     \quad   \sum_{k'\in\calM_L}c_{\obm,k'}^{\bl,L} \vrho{L,\dd}{k,k'}  
    &= \sum_{j=1}^\Nat \vrho{l_j,\dd}{m_j,m_j}c_{\obm,k}^{\bl,L} + \sum_{\obm' \in \calN_\obm} c_{\obm',k}^{\bl,L}
    \sum_{\bmu\in\obm}
    \sum_{j=1}^\Nat
    \vrho{l_j,\dd}{m'_j,\mu_j}.
 \end{align*}
 Now given $\obm\in \oMl$ with corresponding 
 $\blamb{}{\obm} = (\blamb{}{\obm^{(1)}},\ldots,\blamb{}{\obm^{(\Nb)}})$, for any $\obm'\in \calN_\obm$, $\obm' = \obm_{(i,p,q)}$ for some $i\in\{ 1,\ldots, \Nb\}$, $p,q \in \calM_{\bl^{(i)}_1},\;\; p\ne q
     $ defined in~\eqref{eq:interacting_classes}, a short calculation noticing that $\bbm'$ is ordered shows that
\[
\sum_{\bmu\in\obm}
    \sum_{j=1}^\Nat
    \vrho{l_j,\dd}{m'_j,\mu_j} = 
    \blamb{}{{\obm'}^{(i)},q}
    \; \vrho{\bl^{(i)}_1,\dd}{q,p}.
\]
Hence, for all $\dd =1,\ldots,N_{\rm{dim}},\; k\in\calM_L,\; \obm \in \oMl,$ there holds
\begin{align*}
    \sum_{k'\in\calM_L}c_{\obm,k'}^{\bl,L} \vrho{L,\dd}{k,k'}  
    &- \sum_{j=1}^\Nat \vrho{l_j,\dd}{m_j,m_j}c_{\obm,k}^{\bl,L} - 
    \sum_{i=1}^\Nb 
    \sum_{\substack{p,q=1 \\ p \neq q}}^{\card{\calM_{l_i}}}
     c_{\obm_{(i,p,q)},k}^{\bl,L}
    \;\blamb{}{{\obm'}^{(i)},q}
    \; \vrho{\bl^{(i)}_1,\dd}{q,p}=0, 
\end{align*}
which indicates that the vector of coupling coefficients $\bc^{\bl,L}:=\left(c_{\bbm,k}^{\bl,L}\right)_{\obm\in\oMl,k\in\mathcal{M}_L}$ belongs to $\ker(M^{\bl,L})$. 

Reciprocally, assume that $\bc^{\bl,L} \in\ker(M^{\bl,L})$, i.e., that~\eqref{eq:diff-rep-GEPI}holds true.
Multiplying on the left by $d\rep{L}{}(X)$, using~\eqref{eq:diff-rep-GEPI} and~\eqref{eq:TS_lemma_lie_alg} yields
\begin{align*}
     \forall X\in\mathfrak g, \quad \left(d\rep{L}{}(X)\right)^2\left(C^{\bl,L}\right)^T &= d\rep{L}{}(X)\left(C^{\bl,L}\right)^TS^\bl d\rep{\bl}{}(X)T^\bl, \\
     & = \left(C^{\bl,L}\right)^TS^\bl d\rep{\bl}{}(X) T^\bl S^\bl d\rep{\bl}{}(X)T^\bl, \\
     & = \left(C^{\bl,L}\right)^TS^\bl\left(d\rep{\bl}{}(X)\right)^2 T^\bl.
\end{align*}
Iterating this argument as in the GE case and using~\cite[Equation (3.11)]{Hall2013-af}, we see that
\begin{equation}\label{eq:recip-cond-GEPI}
    \forall X\in\mathfrak g,\quad \rep{L}{}\left(\exp(X)\right)\left(C^{\bl,L}\right)^T = \left(C^{\bl,L}\right)^TS^\bl\rep{\bl}{}\left(\exp(X)\right)T^\bl.
\end{equation}
Finally, expressing any element $g\in G$ in the connected linear Lie group $G$ as a finite product of exponentials, we recover~\eqref{eq:matrix_GE-PI_sylvester} from~\eqref{eq:recip-cond-GEPI} using~\eqref{eq:TS_lemma} and that $\rep{\bl}{}$ and  $\rep{L}{}$ are representations.
The basis and dimensionality of $V^{\bl,L}$ are straightforward consequences.
\end{proof}

\begin{proof}[\bf{Proof of Proposition~\ref{prop:GE_rec}}]
First, since $V^\bl = V^{\bl^{(1)}} \otimes V^{\bl^{(2)}}$,~\eqref{eq:direct_sum_GE} yields that
\[ 
([b^{L_1}_{\bl^{(1)} i_1}]_{k_1} [b^{L_2}_{\bl^{(2)} i_2}]_{k_2})_{L_1,L_2 \in \N, i_1 \in \{ 1, \ldots, \dim{}{}(\bl^{(1)},L_1)\}, i_2 \in \{1, \ldots, \dimge{\bl^{(2)}}{L_2}\} }
\]
is a basis of $V^\bl$. Thus, 
any function $F^{\bl,L}\in [V^\bl]^{\card{\calM_L}}$ can be decomposed as 
\[
F^{\bl,L}
   = \left[ \sum_{k_1 \in \calM_{L_1}} \sum_{k_2 \in \calM_{L_2}}
    c^{\bl,L}_{(k_1,k_2), k}
    [b^{L_1}_{\bl^{(1)} i_1}]_{k_1} [b^{L_2}_{\bl^{(2)} i_2}]_{k_2} \right]_{k\in \calM_L},
\]
for some coefficients $\bc^{\bl,L}_{(k_1,k_2), k} \in \K$.
The group action over $F^{\bl,L}$ reads as 
\begin{equation}\label{eq:F-rec}
   \forall g\in G, \; \bR\in \Omega^\Nat, \quad F^{\bl,L}(g \cdot \bR) = 
   \left[ \sum_{k_1, k_1'\in \calM_{L_1}} \sum_{k_2 k_2' \in \calM_{L_2}} 
    c^{\bl,L}_{(k_1,k_2), k}
    \rep{L_1}{k_1,k_1'}(g) \rep{L_2}{k_2,k_2'}(g)
    [b^{L_1}_{\bl^{(1)} i_1}]_{k_1'} [b^{L_2}_{\bl^{(2)} i_2}]_{k_2'} \right]_{k\in \calM_L}.
\end{equation}
Moreover, $F^{\bl,L}$ is equivariant with respect to the representation $\rep{L}{}$ if and only if
\[
    \forall g\in G,\; \bR\in\Omega^\Nat,\; k\in \calM_L, 
    \quad F^{\bl,L}_{k}(g \cdot \bR) = \sum_{k'\in\calM_L} \rep{L}{k,k'}(g)F_{k'}^{\bl,L}(\bR).
\]
Switching $k_1,k_2$ and $k_1',k_2'$ in the right-hand side of~\eqref{eq:F-rec} and writing the equality on the considered basis, we obtain that 
$F^{\bl,L}$ is GE with respect to the representation $\rep{L}{}$ if and only if
\begin{equation*}
    \forall k_1 \in \calM_{L_1},\; k_2 \in \calM_{L_2},\; k\in\calM_L,\; g\in G, \quad \sum_{k'\in\calM_L} c^{\bl,L}_{(k_1,k_2), k'} \rep{L}{k,k'}(g)  
   = \hspace{-3mm} \sum_{k_1 \in \calM_{L_1}} \sum_{k_2 \in \calM_{L_2}}
   \hspace{-3mm}
   c^{\bl,L}_{(k_1',k_2'), k} \rep{L_1}{k_1,k_1'}(g) \rep{L_2}{k_2,k_2'}(g). 
\end{equation*}
These are the same equations as in Proposition~\ref{prop:GE_basis}, equation~\eqref{eq:GE-coeff}, with $(L_1,L_2)$ replacing $\bl$, hence the solutions are the same, as stated in the proposition.
The dimensionality \eqref{eq:dim_GE_rec} is then a direct consequence.
\end{proof}


\subsection{Proofs of Section~\ref{sec:appl}}\label{sec:proof_appl}

\subsubsection{Proof of Proposition~\ref{prop:ker-M2-GE}}\label{sec:proof-half-mat}


Before proving Proposition~\ref{prop:ker-M2-GE}, we first show that either of the upper half or the lower half of the matrix $\MlL_{2}$ defined in~\eqref{eq:M2} (c.f. Figure~\ref{fig:block_structure_of_M}(a)) is sufficient to determine its kernel, and therefore the coupling coefficients. 
Thus,  
we decompose 
\begin{equation}\label{eq:split-M}
    \MlL_{2} = \begin{pmatrix}
         \Mm_{2} \\
         \Mp_{2} \\
     \end{pmatrix},
\end{equation}
with 
$\Mm_{2}=\big[\Mm_{K,K'}\big]_{K\in \calM_{L+1}\backslash\{L,L+1\},~K'\in\calM_L}$ 
and
$\Mp_{2}=\big[\Mp_{K,K'}\big]_{K\in\calM_{L+1}\backslash\{-L,-L-1\},~K'\in\calM_L}$ 
with each block being defined as
\begin{equation*}
\Mpm_{K,K'} = \begin{cases}
    \Apm_{K}, 
    \quad & \textup{if } K'=K, \\
    \Bpm_{K\mp1}, 
    \quad & \textup{if } K'=K\mp 1, \\
    \mathbf{0}, \quad &\textup{otherwise.}
    \end{cases}
\end{equation*}
The matrices $\Apm_{K}\in\R^{\card{\calM_{\bl,K}}\times\card{\calM_{\bl,K}}}$ are diagonal matrices 
\[
    \Apm_{K} = a^\pm_K I_{\card{\calM_{\bl,K}}},
\]
with 
\begin{equation}
\label{eq:akpm2}
    a^\pm_K = \mp\sqrt{(L\mp K+1)(L\pm K)}.
\end{equation}
In particular, there hold
\begin{equation}
\label{eq:akpm}
    a^+_{K} = -a^-_{K-1},
\end{equation}
and
\begin{equation}\label{eq:A-diff}  
    (a^+_K)^2 - (a^-_K)^2= 2K.
\end{equation}
The matrices $\Bpm_{K}$ are defined as $\Bpm_{K}=\sum_{j=1}^N \Bpm_{K,j}$, where $\Bpm_{K,j}\in\R^{\card{\calM_{\bl,K\pm1}}\times\card{\calM_{\bl,K}}}$ is defined for $\bbm\in\calM_{\bl,K\pm1}$ and $\bbm' \in \calM_{\bl,K}$ as
\begin{equation}
\label{eq:Bpm}
        \Bpm_{K,j}[\bbm,\bbm'] 
        = \pm\sqrt{(l_j\mp m_j+1)(l_j\pm m_j)}\delta_{\bbm',\bbm_j^\mp} 
        = \pm\sqrt{(l_j\pm m_j'+1)(l_j\mp m_j')}\delta_{\bbm,({\bbm'}_j)^\pm}.
\end{equation}
While in the matrix $\MlL_2$, there appear only $B^\pm_K, \; K\in\calM_L$, we mention that the definition of the $B^\pm_K$ matrices applies to general $K\in\calM_\infty$, where $\calM_\infty=\{\pm L \;,\; L\in\calL_G\}$. Even when the set $\calM_{\bl,K \pm 1}$ (or $\calM_{\bl,K}$) is empty, the corresponding $B^\pm_K$ can be viewed as a generalized matrix with 0 rows (or 0 columns). We collect below some useful properties of the $B^\pm_K$ matrices. 

\begin{lemma}[Properties of the $B_K$ matrices]\label{lemma:laddar}
    Let $\bl\in\calL^N$, $L\in\calL_G,\;K\in\calM_\infty$, 
    and let the matrices $B_K^\pm$ be defined in~\eqref{eq:Bpm}, then
    \begin{equation}\label{eq:Bequality}
        (B_{K}^-)^T = -B_{{K-1}}^+.
    \end{equation}
    and
    \begin{equation}
        \big(\Bp_K\big)^T\Bp_K  - \big(\Bm_K\big)^T\Bm_K 
     = -2K\cdot I_{\card{\calM_{\bl,K}}}\label{eq:ladder2}.
    \end{equation}
\end{lemma}

\begin{proof}
From~\eqref{eq:Bpm}, for $K\in \calM_\infty,\;j\in \{1,\ldots, \Nat\}$, and $(\bbm,\bbm')\in\calM_{\bl,K-1}\times\calM_{\bl,K}$, there holds
\[
    \Bm_{K,j}[\bbm,\bbm'] = -\Bp_{K-1,j}[\bbm',\bbm],
\]
i.e.
    $(\Bm_{K,j})^T=-\Bp_{K-1,j}$.
We obtain~\eqref{eq:Bequality} by summing over $j\in\{1,\ldots, \Nat\}$.
Moreover, for $j,\;j'\in \{1,\ldots, \Nat\}$, $K\in \calM_\infty$, and $\bbm, \;\bbm'\in\calM_{\bl,K}$, 
\begin{equation}\label{eq:BTB}
\begin{split}
    [\big(\Bpm_{K,j}\big)^T \Bpm_{K,j'}][\bbm,\bbm']
    &= 
    \sum_{\bbm''\in\calM_{\bl,K\pm 1}}\Bpm_{K,j}[\bbm'',\bbm]\Bpm_{K,j'}[\bbm'',\bbm'] \\
    &= 
    \sum_{\bbm''\in\calM_{\bl,K\pm 1}}
    \sqrt{(l_j\pm m_j+1)(l_j\mp m_j)(l_{j'}\pm m'_{j'}+1)(l_{j'}\mp m'_{j'})}\delta_{\bbm'',\bbm_j^\pm}\delta_{\bbm'',{\bbm'}_{j'}^\pm}\\
    &= \begin{cases}
        (l_j\pm m_j+1)(l_j\mp m_j),
        \quad &j=j', ~\bbm' = \bbm, \\
        \sqrt{(l_j\pm m_j+1)(l_j\mp m_j)(l_{j'}\pm m'_{j'}+1)(l_{j'}\mp m'_{j'})},
        \quad &j\ne j', ~\bbm'=\bbm\mp{\bf e}_j\pm{\bf e}_{j'}, \\
        0, \qquad &\textup{otherwise.}
    \end{cases}
\end{split}
\end{equation}
By swapping $j$ and $j'$ in~\eqref{eq:BTB} and comparing it with the resulting equation, we see that
    \[
        \forall j\ne j', \qquad \big(\Bp_{K,j})^T \Bp_{K,j'} = \big(\Bm_{K,j'}\big)^T \Bm_{K,j},
    \]
    and
    \[
        \big(\Bp_{K,j}\big)^T\Bp_{K,j} - \big(\Bm_{K,j}\big)^T \Bm_{K,j} = \textup{Diag}\{-2m_j\}_{\bbm\in\calM_{\bl,K}} = -2m_j \cdot I_{\card{\calM_{\bl,K}}}.
    \]
Therefore, combining the two previous expressions gives
\begin{align*}
    \big(\Bp_K\big)^T\Bp_K  - \big(\Bm_K\big)^T\Bm_K 
    & = 
 \sum_{j,j'=1}^N 
            \Big( \big(\Bp_{K,j}\big)^T
            \Bp_{K,j'}
             - \big(\Bm_{K,j'}\big)^T \Bm_{K,j} \Big)
             \nonumber \\
        &= \sum_{j=1}^N 
            \Big( \big(\Bp_{K,j}\big)^T \Bp_{K,j}
             - \big(\Bm_{K,j}\big)^T \Bm_{K,j} \Big) \nonumber \\
        &= -\sum_{j=1}^N 2m_j\cdot I_{\card{\calM_{\bl,K}}} = -2\sum\bbm \cdot I_{\card{\calM_{\bl,K}}} = -2K\cdot I_{\card{\calM_{\bl,K}}},
\end{align*}
which proves~\eqref{eq:ladder2}.
\end{proof}

We can then prove the following proposition.
\begin{proposition}\label{prop:matrix_final}
    For $\bl\in \calL^N$, $L\in \calLG$, there holds
    \[
        \ker(\MlL_{2}) = \ker(\Mm_{2}) = \ker(\Mp_{2}).
    \]
\end{proposition}
\begin{proof}
We first show that $\big(\Mm_{2}\big)^T\Mm_{2} = 
    \big(\Mp_{2}\big)^T\Mp_{2}.$
Due to the block structure of $\Mm_{2}$ (respectively $\Mp_{2}$), the matrix $\big(\Mm_{2}\big)^T\Mm_{2}$ (resp. $\big(\Mp_{2}\big)^T\Mp_{2}$) is also a block matrix with row and column indices taking values from $\calM_L$. Noting that 
$a^+_{-L} = a^-_{L} =0$, 
we have for all $K,K'\in\calM_L$ that
\[
    \Big[\big(\Mm_{2}\big)^T\Mm_{2}\Big]_{K,K'} = 
    \begin{cases}
        \big(\Bm_K\big)^T\Bm_K + \big(\Am_K\big)^T\Am_K, \quad & \textup{if } K'=K, \\
        \big(\Bm_{K}\big)^T\Am_{K-1}, \quad & \textup{if } K'=K-1, \\
        \big(\Am_{K}\big)^T\Bm_{K+1}, \quad & \textup{if } K'=K+1, \\
        \mathbf{0}, \quad & \textup{otherwise,}
    \end{cases}
\]
and 
\[
    \Big[\big(\Mp_{2}\big)^T\Mp_{2}\Big]_{K,K'} = 
    \begin{cases}
        \big(\Bp_K\big)^T\Bp_K + \big(\Ap_K\big)^T\Ap_K, \quad & \textup{if } K'=K, \\
        \big(\Ap_{K}\big)^T\Bp_{K-1}, \quad & \textup{if } K'=K-1, \\
        \big(\Bp_{K}\big)^T\Ap_{K+1}, \quad & \textup{if } K'=K+1, \\
        \mathbf{0}, \quad & \textup{otherwise.}
    \end{cases}
\]
Using~\eqref{eq:akpm} and~\eqref{eq:Bequality}, we can easily show for $K\in \calM_L$, when the corresponding blocks are defined, that
\begin{equation}
    \Big[\big(\Mm_{2}\big)^T\Mm_{2}\Big]_{K,K\pm 1} 
= 
\Big[\big(\Mp_{2}\big)^T\Mp_{2}\Big]_{K,K \pm 1},
\label{eq:offdiagblocks}
\end{equation}
and using~\eqref{eq:A-diff} and ~\eqref{eq:ladder2}, we obtain that for $K\in\calM_L$,
\begin{equation}
    \Big[\big(\Mm_{2}\big)^T\Mm_{2}\Big]_{K,K} 
= 
\Big[\big(\Mp_{2}\big)^T\Mp_{2}\Big]_{K,K }.
\label{eq:diagblocks}
\end{equation}
Combining~\eqref{eq:offdiagblocks} and~\eqref{eq:diagblocks}, we have
$\big(\Mm_{2}\big)^T\Mm_{2} = 
    \big(\Mp_{2}\big)^T\Mp_{2}.$
Thus, for any $\bc\in\ker(\Mm_{2})$,
\[
    \|\Mm_{2}\bc\|^2 = \bc^T(\Mm_{2})^T\Mm_{2}\bc=\bc^T(\Mp_{2})^T\Mp_{2}\bc=\|\Mp_{2}\bc\|^2=0,
\]
which leads to $\bc\in\ker(\Mp_{2})$, and vice versa. Consequently, $\ker(\Mm_{2})=\ker(\Mp_{2})$. We complete the proof by using \eqref{eq:split-M}.
\end{proof}

We are now in the position of proving Proposition~\ref{prop:ker-M2-GE}. 

\begin{proof}[\bf{Proof of Proposition~\ref{prop:ker-M2-GE}}]
    We may write
    $\Mup=\big[\Mup_{K,K'}\big]_{K\in\calM_{L+1}\backslash \{L,-L-1\},~K'\in\calM_L}$, with each block being defined as
    \begin{equation*}
    \Mup_{K,K'} = \begin{cases}
        \Am_{K}, 
        \quad & \textup{if } K'=K, \\
        \Bm_{K\mp1}, 
        \quad & \textup{if } K'=K\mp1\ne L, \\
        \Bp_{L}, 
        \quad & \textup{if } K'=K-1=L, \\
        \mathbf{0}, \quad &\textup{otherwise.}
        \end{cases}
    \end{equation*}
    Using \eqref{eq:ladder2} with $K=L+1$, we have
    \begin{equation*}
        \big(\Bm_{L+1}\big)^T\Bm_{L+1} = \big(\Bp_{L+1}\big)^T\Bp_{L+1} +2(L+1) \cdot I_{\card{\calM_{\bl,L+1}}},
    \end{equation*}
    which yields that $\big(\Bm_{L+1}\big)^T\Bm_{L+1}$ is positive definite, so $\Bm_{L+1}$ has full rank. We obtain using~\eqref{eq:Bequality} that $\Bp_{L}$ has full rank.
    Thus, since the matrices $\Am_{K}$ 
    (yellow blocks on Figure~\ref{fig:block_structure_of_M}(a)) are scaled identities with non-zero diagonal elements for $K\in\calM_L\setminus \{L\}$, the matrix $\Mup$ has full rank.
    Since $\Mup$ consists of a subset of the rows of $\MlL_2$, we have $\ker(\MlL_2)\subset\ker(\Mup)$, 
    and using Proposition~\ref{prop:matrix_final}
    we obtain
    \begin{equation}\label{eq:dim-bound}
        \dim\big(\ker(\Mup)\big)\ge \dim\big(\ker(\MlL_2)\big)=\dim\big(\ker(\Mm_2)\big).
    \end{equation}
    Noting that $\Mup$ and $\Mm_2$ have the same size, equation~\eqref{eq:dim-bound} together with the rank–nullity theorem indicate that $\text{rank}(\Mup)\le\text{rank}(\Mm_2)$. However, as $\Mup$ has full rank, $\Mm_2$ must have full rank, i.e., $\text{rank}(\Mup)=\text{rank}(\Mm_2)$. Therefore
    \[
        \dim\big(\ker(\Mup)\big)=\dim\big(\ker(\Mm_2)\big) = \dim\big(\ker(\MlL_2)\big).
    \]
    Recalling that $\ker(\MlL_2)\subset\ker(\Mup)$, we easily obtain the result.
\end{proof}


\subsubsection{Proof of Proposition~\ref{prop:asym-dim-ge}}\label{sec:proof-asym-dim-ge}

\begin{proof}[\bf{Proof of Proposition~\ref{prop:asym-dim-ge}}]
By definition, $\card{\calM_{\bl,L}}$ is the number of $\bbm\in\calM_\bl$ such that $\sum\bbm=L$.
For $\bbm\in\calM_\bl$, we may view each $(m_i)_{i=1}^N$ as independent random variables uniformly distributed on $\calM_{l_i}$. Hence, $\sum\bbm = \sum_{i=1}^N m_i$ is a random variable with mean 0,  variance $\textup{Var}_{\bl} = \sum_{i=1}^N l_i(l_i+1)/3$ and fourth order cumulant 
\[
\textup{C}_{4\bl} = \sum_{i=1}^\Nat \left[ \frac{1}{|\calM_{l_i}|}\sum_{k\in\calM_{l_i}} k^4 - \frac{3}{|\calM_{l_i}|^2} \left( \sum_{k\in\calM_{l_i}} k^2 \right)^2 \right] = -\sum_{i=1}^N \frac{l_i (l_i+1) (2l_i^2 + 2l_i+1)}{15}.
\]
Under the condition~\eqref{eq:bounded-ll}, we have that $O(\textup{Var}_{\bl})=O(\textup{C}_{4\bl})=O(N)$ for $N$ sufficiently large. 
As a consequence of~\cite[Ch.~VII, Theorem~12]{petrov1975-si}, we have
\begin{equation*}
    P
    \left(\sum\bbm=L{\Big|}\bbm\in\calM_{\bl}
    \right) = 
    \frac{1}{\sqrt{2\pi \textup{Var}_{\bl}}}
    \exp\left(\frac{-L^2}{2\textup{Var}_{\bl}}\right)
        \Big(1 + \frac{\textup{C}_{4\bl}}{24 \textup{Var}_{\bl}^{2}}H_4
        \left(\frac{L}{\sqrt{\textup{Var}_{\bl}}}\right) + O\left(\frac{1}{\Nat^{2}}\right)\Big),
\end{equation*}
where $H_4(x) = x^4-6x^2+3$ stands for the Hermite polynomial of order $4$.
Since
$    |\calM_\bl|=\prod_{i=1}^N\big(2l_i+1\big), $
there holds
\begin{equation*}
    \card{\calM_{\bl,L}} = 
    \frac{\prod_{i=1}^N\big(2l_i+1\big)}{\sqrt{2\pi \textup{Var}_{\bl}}}
    \exp\left(\frac{-L^2}{2\textup{Var}_{\bl}}\right)
        \left(1 + \frac{\textup{C}_{4\bl}}{24 \textup{Var}_{\bl}^{2}}H_4\left(\frac{L}{\sqrt{\textup{Var}_{\bl}}}\right) + O\left(\frac{1}{\Nat^{2}}\right)\right).
\end{equation*}
Noting that for $L\ll \sqrt{\textup{Var}_{\bl}}$
\[
    H_4 \left(\frac{L}{\sqrt{\textup{Var}_{\bl}}} \right) - H_4 \left(\frac{L+1}{\sqrt{\textup{Var}_{\bl}}}\right) = O\left(\frac{1}{N}\right),
\]
and 
\[
    \exp\left(\frac{-L^2}{2\textup{Var}_{\bl}}\right) -  \exp\left(\frac{-(L+1)^2}{2\textup{Var}_{\bl}}\right)
    = 
    \frac{2L+1}{2\textup{Var}_{\bl}} + O\left(\frac{1}{N^2}\right), 
\]
we easily obtain~\eqref{eq:dimge-est}. 
\end{proof}

\subsubsection{Proof of Proposition~\ref{prop:ker-M2-GEPI}}\label{sec:proof-half-mat-pi}
Proposition~\ref{prop:ker-M2-GEPI} can be proven using a similar protocol as the proof of Proposition~\ref{prop:ker-M2-GE}. Concretely, we first show that either of the upper half or the lower half of the matrix $\MlL_{2}$ defined in~\eqref{eq:M2-PI} (Figure~\ref{fig:block_structure_of_M}(a)) is sufficient to determine its kernel. With a slight abuse of notation, we denote the two halves of $\MlL_2$ as $\MlL_2^\mp$, respectively. Then,
\begin{equation}\label{eq:split-M-PI}
    \MlL_2 = \begin{pmatrix}
        \MlL_2^- \\
        \MlL_2^+
    \end{pmatrix},
\end{equation}
where 
$\Mm_{2}=\big[\Mm_{K,K'}\big]_{K\in \calM_{L+1}\backslash\{L,L+1\},~K'\in\calM_L}$ 
and
$\Mp_{2}=\big[\Mp_{K,K'}\big]_{K\in\calM_{L+1}\backslash\{-L,-L-1\},~K'\in\calM_L}$ 
whose blocks are defined as
\begin{equation*}
\Mpm_{K,K'} = \begin{cases}
    \Apm_{K}, 
    \quad & \textup{if } K'=K, \\
    \Bpm_{K\mp1}, 
    \quad & \textup{if } K'=K\mp 1, \\
    \mathbf{0}, \quad &\textup{otherwise.}
    \end{cases}
\end{equation*}

Similar to the GE case, the matrices $\Apm_{K}$ are scaled identity matrices
\begin{equation*}
    \Apm_{K} = a^\pm_K I_{\card{\overline{\calM}_{\bl,K}}},
\end{equation*}
with $a^\pm_K$ defined in~\eqref{eq:akpm2}. In particular, ~\eqref{eq:akpm} and~\eqref{eq:A-diff} hold true in this case as well. 
The matrices $\Bpm_{K}$ are defined as
\[  
    \Bpm_{K}=\sum_{j=1}^\Nb\sum_{p\in\calM_{\bl^{(j)}_1}} \Bpm_{K,j,p}, 
\]
where $\Bpm_{K,j,p}\in\R^{\card{\overline{\calM}_{\bl,K\pm1}}\times\card{\oMlK}}$ is defined for $\bbm\in \overline{\calM}_{\bl,K\pm1}$ and $\bbm'\in\oMlK$ as
\begin{equation}\label{eq:Bpm-pi}
\begin{split}
    \Bp_{K,j,p}[\bbm,\bbm'] & = \blamb{}{{\bbm'}^{(j)},p-1}\sqrt{(\bl^{(j)}_1-p+1)(\bl^{(j)}_1+p)}\delta_{\bbm',\bbm_{(j,p,p-1)}} \\
    & = \big(\blamb{}{\bbm^{(j)},p-1}+1\big)\sqrt{(\bl^{(j)}_1-p+1)(\bl^{(j)}_1+p)}\delta_{\bbm,\bbm'_{(j,p-1,p)}}, \\
     \Bm_{K,j,p}[\bbm,\bbm'] & = 
    -\blamb{}{\bbm'^{(j)},p+1}\sqrt{(\bl^{(j)}_1+p+1)(\bl^{(j)}_1-p)}\delta_{\bbm',\bbm_{(j,p,p+1)}} \\
    & = -\big(\blamb{}{\bbm^{(j)},p+1}+1\big)\sqrt{(\bl^{(j)}_1+p+1)(\bl^{(j)}_1-p)}\delta_{\bbm,\bbm'_{(j,p+1,p)}}, 
\end{split}
\end{equation}
from which we obtain, noting that above definition for $B^\pm_K$ works for all $K\in\calM_\infty$, that for $K\in\calM_\infty$,
\begin{equation}
\label{eq:Bequality-pi}
    \forall\bbm\in\overline{\calM}_{\bl,K-1}, \;\bbm'\in\oMlK, \quad \big(\blamb{}{{\bbm}^{(j)},p+1}+1\big)\Bp_{K-1,j,p+1}[\bbm',\bbm] = -\blamb{}{\bbm^{(j)},p}\;\Bm_{{K,j,p}}[\bbm,\bbm'],
\end{equation}
or equivalently,
\begin{equation*}
    (\Bm_{{K,j,p}})^TW_{K-1,j,p} = -W_{K,j,p+1}\Bp_{K-1,j,p+1}, 
\end{equation*}
where 
\[
    W_{K,j,p} = \textup{Diag}\big(\blamb{}{{\bbm}^{(j)},p}\big)_{\bbm\in\overline{\calM}_{\bl,K}}.
\]
We have the following proposition for the matrices $\Bpm_{K}$.
\begin{lemma}[Properties of the $B_K$ matrices]\label{lem:laddar-pi}
    Let $\bl\in\calL^N$, $L\in\calL_G,$ $K\in\calM_\infty$, 
    and let the matrices $B_K^\pm$ be defined in~\eqref{eq:Bpm-pi}, then
    \begin{equation}\label{eq:per-sym-B-pi}
        \big(\Bm_{K}\big)^TW_{K-1} = - W_{K}\Bp_{K-1}, 
    \end{equation}
    and 
    \begin{equation}\label{eq:ladder-pi}
        \big(\Bp_{K}\big)^TW_{K+1}\Bp_{K} - \big(\Bm_{K}\big)^TW_{K-1}\Bm_{K} = -2K\cdot W_K,
    \end{equation}
    with $W_K$ being a diagonal matrix defined as
    \begin{equation}
    \label{eq:WK}
                W_K = \prod_{j=1}^\Nb\prod_{p\in\calM_{\bl^{(j)}_1}}
        (W_{K,j,p}) !,
    \end{equation}
    where the factorial is applied pointwise.
\end{lemma}
\begin{proof}
    We evaluate elements on both sides of~\eqref{eq:per-sym-B-pi} for $\bbm'\in\overline{\calM}_{\bl,K}$ and $\bbm\in\overline{\calM}_{\bl,K-1}$. Note that there exist at most a pair of $(j,p)$ such that $\bbm'=\bbm_{(j,p,p-1)}$, in which case 
    the consistency of elements in both sides of~\eqref{eq:per-sym-B-pi} is verified with
    ~\eqref{eq:Bequality-pi}. 
    In the case of the absence of such a pair, both sides are $0$. Hence,~\eqref{eq:per-sym-B-pi} is valid. 

    Additionally, following a similar calculation to that in the proof of Lemma~\ref{lemma:laddar} and using~\eqref{eq:Lamb_blL}, $\Bm_{K+1}\Bp_{K} - \Bp_{K-1}\Bm_{K}$ is a diagonal matrix with diagonal elements being defined for $\bbm\in\overline{\calM}_{\bl,K}$ and satisfying
    \[
    \begin{split}
        (\Bm_{K+1}\Bp_{K} - \Bp_{K-1}\Bm_{K})[\bbm,\bbm] & = \sum_{j=1}^\Nb \sum_{p\in\calM_{\bl^{(j)}_1}} 2p\cdot \blamb{}{\bbm^{(j)},p} \\
        & = 2\sum\bbm = 2K.
    \end{split}
    \]
    Thus, 
    \begin{equation}\label{eq:ladder-pi-inner}  
        \Bm_{K+1}\Bp_{K} - \Bp_{K-1}\Bm_{K} = 2K\cdot I_{\card{\overline{\calM}_{\bl,K}}}.
    \end{equation}
    We complete the proof by multiplying both sides of~\eqref{eq:ladder-pi-inner} by $W_K$ on the left, and using~\eqref{eq:per-sym-B-pi}.
    \end{proof}
We then see that the following proposition holds true. 
\begin{proposition}
\label{prop:matrix_final_pi}
    For $\bl\in \calL^N$, $L\in \calLG$, there hold
    \[
        \ker(\MlL_{2}) = \ker(\Mm_{2}) = \ker(\Mp_{2}).
    \]
\end{proposition}
\begin{proof}
We show that $\big(\Mm_{2}\big)^TW^-\Mm_{2} = 
    \big(\Mp_{2}\big)^TW^+\Mp_{2}$,
where 
\[
\begin{split}
    W^- = \textup{Diag}(W_K)_{K\in\{-L-1,\ldots,L-1\}}, \\
    W^+ = \textup{Diag}(W_K)_{K\in\{-L+1,\ldots,L+1\}}.
\end{split}
\]
Clearly, $W^\pm$ are positive definite. Due to the block structure of $W^\pm$ and $\Mm_{2}$ (respectively $\Mp_{2}$), the matrix $\big(\Mm_{2}\big)^TW^-\Mm_{2}$ (resp. $\big(\Mp_{2}\big)^TW^+\Mp_{2}$) is also a block matrix with row and column indices taking values from $\calM_L$. Specifically, $\forall K,K'\in\calM_L$, noting that $a^-_{L} = a^+_{-L} =0,$
\[
    \Big[\big(\Mm_{2}\big)^TW^-\Mm_{2}\Big]_{K,K'} = 
    \begin{cases}
        \big(\Bm_K\big)^TW_{K-1}\Bm_K + \big(\Am_K\big)^TW_{K}\Am_K, \quad & \textup{if } K'=K, \\
        \big(\Bm_{K}\big)^TW_{K-1}\Am_{K-1}, \quad & \textup{if } K'=K-1, \\
        \big(\Am_{K}\big)^TW_{K}\Bm_{K+1}, \quad & \textup{if } K'=K+1, \\
        \mathbf{0}, \quad & \textup{otherwise,}
    \end{cases}
\]
and 
\[
    \Big[\big(\Mp_{2}\big)^TW^+\Mp_{2}\Big]_{K,K'} = 
    \begin{cases}
        \big(\Bp_K\big)^TW_{K+1}\Bp_K + \big(\Ap_K\big)^TW_{K}\Ap_K, \quad & \textup{if } K'=K, \\
        \big(\Ap_{K}\big)^TW_{K}\Bp_{K-1}, \quad & \textup{if } K'=K-1, \\
        \big(\Bp_{K}\big)^TW_{K+1}\Ap_{K+1}, \quad & \textup{if } K'=K+1, \\
        \mathbf{0}, \quad & \textup{otherwise.}
    \end{cases}
\]
Then, we see from equations~\eqref{eq:akpm} and~\eqref{eq:per-sym-B-pi} that their corresponding off-diagonal blocks are identical, and from equations~\eqref{eq:akpm} and~\eqref{eq:ladder-pi} that their corresponding diagonal blocks also coincide.

Thus, we have shown
$\big(\Mm_{2}\big)^TW^-\Mm_{2} = \big(\Mp_{2}\big)^TW^+\Mp_{2}$. 
As a result, for any $\bc\in\ker(\Mm_{2})$,
\[
    \|\Mp_{2}\bc\|_{W^+}^2 = \bc^T(\Mp_{2})^TW^+\Mp_{2}\bc=\bc^T(\Mm_{2})^TW^-\Mm_{2}\bc=\|\Mm_{2}\bc\|_{W^-}^2=0,
\]
which means $\bc\in\ker(\Mp_{2})$, and vice versa. Consequently, $\ker(\Mm_{2})=\ker(\Mp_{2})$. We complete the proof by using \eqref{eq:split-M-PI}.
\end{proof}

We are now in the position of proving Proposition~\ref{prop:ker-M2-GEPI}. 


\begin{proof}[\bf{Proof of Proposition~\ref{prop:ker-M2-GEPI}}]
    We may write
    $\Mup=\big[\Mup_{K,K'}\big]_{K\in\calM_{L+1}\backslash \{L,-L-1\},~K'\in\calM_L}$, with each block being defined as
        \begin{equation*}
        \Mup_{K,K'} = \begin{cases}
            \Am_{K}, 
            \quad & \textup{if } K'=K, \\
            \Bm_{K\mp1}, 
            \quad & \textup{if } K'=K\mp1\ne L, \\
            \Bp_{L}, 
            \quad & \textup{if } K'=K-1=L, \\
            \mathbf{0}, \quad &\textup{otherwise.}
            \end{cases}
        \end{equation*}
    Using~\eqref{eq:ladder-pi} with $K=L+1$, 
    we obtain 
    \begin{equation*}
        \big(\Bm_{L+1}\big)^TW_{L}\Bm_{L+1} = \big(\Bp_{L+1}\big)^TW_{L+2}\Bp_{L+1} + 2(L+1)\cdot W_{L+1},
    \end{equation*}
    which yields that $\big(\Bm_{L+1}\big)^TW_{L}\Bm_{L+1}$
    is positive definite noting that $W_L,W_{L+1},W_{L+2}$ defined in~\eqref{eq:WK} are diagonal matrices with strictly positive elements. Hence, $\Bm_{L+1}$ has full rank. We obtain using~\eqref{eq:Bequality-pi} that $\Bp_{L}$ has full rank.
    Thus, since the matrices $\Am_{K}$ 
    (yellow blocks on Figure~\ref{fig:block_structure_of_M}(a)) are scaled identities with non-zero diagonal elements for $K\in\calM_L\setminus \{L\}$, the matrix $\Mup$ has full rank.
    Since $\Mup$ consists of a subset of the rows of $\MlL_2$, we have $\ker(\MlL_2)\subset\ker(\Mup)$, which combined with Proposition~\ref{prop:matrix_final_pi} leads to
    \begin{equation}\label{eq:dim-bound-pi}
        \dim\big(\ker(\Mup)\big)\ge \dim\big(\ker(\MlL_2)\big)=\dim\big(\ker(\Mm_2)\big).
    \end{equation}
    Noting that $\Mup$ and $\Mm_2$ have the same size, equation~\eqref{eq:dim-bound-pi}, together with the rank–nullity theorem, indicate that $\text{rank}(\Mup)\le\text{rank}(\Mm_2)$. However, as $\Mup$ has full rank, $\Mm_2$ must have full rank, i.e., $\text{rank}(\Mup)=\text{rank}(\Mm_2)$. Therefore
    \[
        \dim\big(\ker(\Mup)\big)=\dim\big(\ker(\Mm_2)\big) = \dim\big(\ker(\MlL_2)\big).
    \]
    We conclude recalling that $\ker(\MlL_2)\subset\ker(\Mup)$.
\end{proof}




\section*{Acknowledgements}

We would like to thank Christoph Ortner for insightful discussions. 
This work has received funding from the ANR NUMERIQ, project number ANR-24-CE46-2255. 
This work has been supported by the EIPHI Graduate school (contract ANR-17-EURE-0002) and by the Région Bourgogne Franche-Comté.
L.Z. acknowledges funding by  the Deutsche Forschungsgemeinschaft (DFG, German Research Foundation) - Project number 442047500 through the Collaborative Research Center “Sparsity and Singular Structures” (SFB 1481). 

\bibliographystyle{siam}
\bibliography{biblio}

\appendix



\section{Proof of Proposition~\ref{prop:deri_alpha_gamma}}\label{app:Wigner-Dmatricesproof}
\begin{proof}[Proof of Proposition~\ref{prop:deri_alpha_gamma}]
Taking the derivative of $D^l_{\mu m}$ with respect to $\alpha$ and $\gamma$, respectively, we have that
\[
    \frac{\partial D^l_{\mu m}}{\partial \alpha}(\alpha, \beta, \gamma)  = -\mathrm{i}\cdot me^{-\mathrm{i}m\alpha}d^l_{\mu m}(\beta)e^{-\mathrm{i}\mu\gamma},
\]
and
\[
    \frac{\partial D^l_{\mu m}}{\partial \gamma}(\alpha, \beta, \gamma)  = -\mathrm{i}\cdot \mu e^{-\mathrm{i}m\alpha}d^l_{\mu m}(\beta)e^{-\mathrm{i}\mu\gamma}.
\]
Inserting $(\alpha,\beta,\gamma) = (0,0,0)$ and noticing that $d_{\mu m}^l(0) = \delta_{\mu m}$, we arrive at~\eqref{eq:deri_alpha} and~\eqref{eq:deri_gamma}.

Differentiating the formula for the Wigner-D matrix with respect to $\beta$ gives
\begin{align*}
    \frac{\partial D^l_{\mu m}}{\partial \beta} & (\alpha, \beta, \gamma) = e^{-\mathrm{i}(m\alpha+ \mu\gamma)} [(l+m)!(l-m)!(l+\mu)!(l-\mu)!]^{\frac12} \\ &\sum_s \frac{(-1)^{s}}{2}\frac{(\cos\frac{\beta}{2})^{2l + \mu - m - 2s - 1}(\sin\frac{\beta}{2})^{m-\mu+2s-1}[(-2l-\mu+m+2s)(\sin\frac{\beta}{2})^2 + (m-\mu+2s)(\cos\frac{\beta}{2})^2]}{(l+\mu-s)!s! (m-\mu+s)!(l-m-s)!}.
\end{align*}
Let $(\alpha,\beta,\gamma) = (0,0,0)$, we see that the above derivative can be non-zero only when $m-\mu+2s-1=0$ or $m-\mu+2s+1=0$. Since $s\in \{\max(0, \mu-m),\ldots,\min(l+\mu,l-m)\}$, the latter is never true since its left hand side is always positive. In addition, the previous conditions can be fulfilled only when $|m-\mu| = 1$. Otherwise, the derivative with respect to $\beta$ at the origin will be zero.

If $m-\mu=1$, then only $s=0$ can contribute to the summation. Thus 
\begin{align*}
    \frac{\partial D^l_{\mu m}}{\partial \beta}(I)
    & = [(l+m)!(l-m)!(l+\mu)!(l-\mu)!]^{\frac{1}{2}}\frac{(-1)^{0}}{2}\frac{1}{(l+\mu)!0!1!(l-m)!} \\
    & = \frac{1}{2}\frac{[(l+m)!(l-m)!(l+m-1)!(l-m+1)!]^{\frac{1}{2}}}{(l+m-1)!(l-m)!} \\
    & = \frac{1}{2}\frac{[(l+m-1)!(l+m)(l-m)!(l+m-1)!(l-m)!(l-m+1)]^{\frac{1}{2}}}{(l+m-1)!(l-m)!} \\
    & = \frac{1}{2}\frac{[(l+m)(l-m+1)]^{\frac{1}{2}}(l+m-1)!(l-m'-1)!}{(l+m-1)!(l-m'-1)!} \\
    & = \frac{1}{2}[(l+m)(l-m+1)]^{\frac{1}{2}}. \\
\end{align*}

Similarly, when $m-\mu=-1$, only $s=1$ contributes to the summation and 
\begin{align*}
    \frac{\partial D^l_{\mu m}}{\partial \beta}(I)
    & = [(l+m)!(l-m)!(l+\mu)!(l-\mu)!]^{\frac{1}{2}}\frac{(-1)^{1}}{2}\frac{1}{(l+\mu-1)!1! 0!(l-m-1)!} \\
    & = - \frac{1}{2}\frac{[(l+m)!(l-m)!(l+m+1)!(l-m-1)!]^{\frac{1}{2}}}{(l+m)!(l-m-1)!} \\
    & = - \frac{1}{2}\frac{[(l+m)!(l-m-1)!(l-m)(l+m)!(l+m+1)(l-m-1)!]^{\frac{1}{2}}}{(l+m)!(l-m-1)!} \\
    & = - \frac{1}{2}[(l-m)(l+m+1)]^{\frac{1}{2}}, \\
\end{align*}
which completes the proof of~\eqref{eq:deri_beta}.
\end{proof} 

\section{Table of dimensionalities}\label{app:table_of_dim}
In this appendix, we provide tables that contain exact values of $\dimge{\bl}{L}$ (denoted by GE) and $\dimpi{\bl}{L}$ (denoted by GE-PI) for some typical $\bl=(l,l,\ldots,l)\in\calL^N$, which is the minimal unit of constructing the equivariant bases (c.f. Propositions~\ref{prop:GE_rec} and ~\ref{prop:GEPI_rec}). 
In this setting, the dimensions $\dimge{\bl}{L}$ and $\dimpi{\bl}{L}$ are determined by three parameters: the degree $l$, the correlation order $N$ and the order of equivariance $L$. We show tables for fixed $N$, or fixed $l$.

\begin{table}[h]
\resizebox{\textwidth}{!}{%
\begin{tabular}{|c|cc|cc|cc|cc|cc|cc|cc|cc|}
\hline
\multirow{2}{*}{$L/\Nat$} &
 \multicolumn{2}{c|}{1} &
 \multicolumn{2}{c|}{2} &
 \multicolumn{2}{c|}{3} &
 \multicolumn{2}{c|}{4} &
 \multicolumn{2}{c|}{5} &
 \multicolumn{2}{c|}{6} &
 \multicolumn{2}{c|}{7} &
 \multicolumn{2}{c|}{8} \\ \cline{2-17} 
&
 \multicolumn{1}{c|}{GE-PI} &
 GE &
 \multicolumn{1}{c|}{GE-PI} &
 GE &
 \multicolumn{1}{c|}{GE-PI} &
 GE &
 \multicolumn{1}{c|}{GE-PI} &
 GE &
 \multicolumn{1}{c|}{GE-PI} &
 GE &
 \multicolumn{1}{c|}{GE-PI} &
 GE &
 \multicolumn{1}{c|}{GE-PI} &
 GE &
 \multicolumn{1}{c|}{GE-PI} &
 GE \\ \hline
0 &
 \multicolumn{1}{c|}{0} &
 0 &
 \multicolumn{1}{c|}{1} &
 1 &
 \multicolumn{1}{c|}{0} &
 1 &
 \multicolumn{1}{c|}{1} &
 3 &
 \multicolumn{1}{c|}{0} &
 6 &
 \multicolumn{1}{c|}{1} &
 15 &
 \multicolumn{1}{c|}{0} &
 36 &
 \multicolumn{1}{c|}{1} &
 91 \\ \hline
1 &
 \multicolumn{1}{c|}{1} &
 1 &
 \multicolumn{1}{c|}{0} &
 1 &
 \multicolumn{1}{c|}{1} &
 3 &
 \multicolumn{1}{c|}{0} &
 6 &
 \multicolumn{1}{c|}{1} &
 15 &
 \multicolumn{1}{c|}{0} &
 36 &
 \multicolumn{1}{c|}{1} &
 91 &
 \multicolumn{1}{c|}{0} &
 232 \\ \hline
2 &
 \multicolumn{1}{c|}{} &
  &
 \multicolumn{1}{c|}{1} &
 1 &
 \multicolumn{1}{c|}{0} &
 2 &
 \multicolumn{1}{c|}{1} &
 6 &
 \multicolumn{1}{c|}{0} &
 15 &
 \multicolumn{1}{c|}{1} &
 40 &
 \multicolumn{1}{c|}{0} &
 105 &
 \multicolumn{1}{c|}{1} &
 280 \\ \hline
3 &
 \multicolumn{1}{c|}{} &
  &
 \multicolumn{1}{c|}{} &
  &
 \multicolumn{1}{c|}{1} &
 1 &
 \multicolumn{1}{c|}{0} &
 3 &
 \multicolumn{1}{c|}{1} &
 10 &
 \multicolumn{1}{c|}{0} &
 29 &
 \multicolumn{1}{c|}{1} &
 84 &
 \multicolumn{1}{c|}{0} &
 238 \\ \hline
4 &
 \multicolumn{1}{c|}{} &
  &
 \multicolumn{1}{c|}{} &
  &
 \multicolumn{1}{c|}{} &
  &
 \multicolumn{1}{c|}{1} &
 1 &
 \multicolumn{1}{c|}{0} &
 4 &
 \multicolumn{1}{c|}{1} &
 15 &
 \multicolumn{1}{c|}{0} &
 49 &
 \multicolumn{1}{c|}{1} &
 154 \\ \hline
5 &
 \multicolumn{1}{c|}{} &
  &
 \multicolumn{1}{c|}{} &
  &
 \multicolumn{1}{c|}{} &
  &
 \multicolumn{1}{c|}{} &
  &
 \multicolumn{1}{c|}{1} &
 1 &
 \multicolumn{1}{c|}{0} &
 5 &
 \multicolumn{1}{c|}{1} &
 21 &
 \multicolumn{1}{c|}{0} &
 76 \\ \hline
6 &
 \multicolumn{1}{c|}{} &
  &
 \multicolumn{1}{c|}{} &
  &
 \multicolumn{1}{c|}{} &
  &
 \multicolumn{1}{c|}{} &
  &
 \multicolumn{1}{c|}{} &
  &
 \multicolumn{1}{c|}{1} &
 1 &
 \multicolumn{1}{c|}{0} &
 6 &
 \multicolumn{1}{c|}{1} &
 28 \\ \hline
7 &
 \multicolumn{1}{c|}{} &
  &
 \multicolumn{1}{c|}{} &
  &
 \multicolumn{1}{c|}{} &
  &
 \multicolumn{1}{c|}{} &
  &
 \multicolumn{1}{c|}{} &
  &
 \multicolumn{1}{c|}{} &
  &
 \multicolumn{1}{c|}{1} &
 1 &
 \multicolumn{1}{c|}{0} &
 7 \\ \hline
8 &
 \multicolumn{1}{c|}{} &
  &
 \multicolumn{1}{c|}{} &
  &
 \multicolumn{1}{c|}{} &
  &
 \multicolumn{1}{c|}{} &
  &
 \multicolumn{1}{c|}{} &
  &
 \multicolumn{1}{c|}{} &
  &
 \multicolumn{1}{c|}{} &
  &
 \multicolumn{1}{c|}{1} &
 1 \\ \hline
\end{tabular}
}
\caption{Dimensions of $V^{\bl,L}$ (GE) and $\bV^{\bl,L}$ (GE-PI) for $\bl = (1,1,\ldots,1)$ with length $\Nat$, and for $L\in\{0,1,\ldots,\Nat\}$.}
\label{tab:l1}
\end{table}

\begin{table}[htb!]
\resizebox{\textwidth}{!}{%
\begin{tabular}{|c|cc|cc|cc|cc|cc|cc|cc|cc|}
\hline
\multirow{2}{*}{$L$/$\Nat$} & \multicolumn{2}{c|}{1} & \multicolumn{2}{c|}{2} & \multicolumn{2}{c|}{3} & \multicolumn{2}{c|}{4} & \multicolumn{2}{c|}{5} & \multicolumn{2}{c|}{6} & \multicolumn{2}{c|}{7} & \multicolumn{2}{c|}{8} \\ \cline{2-17} 
 & \multicolumn{1}{c|}{GE-PI} & GE & \multicolumn{1}{c|}{GE-PI} & GE & \multicolumn{1}{c|}{GE-PI} & GE & \multicolumn{1}{c|}{GE-PI} & GE & \multicolumn{1}{c|}{GE-PI} & GE & \multicolumn{1}{c|}{GE-PI} & GE & \multicolumn{1}{c|}{GE-PI} & GE & \multicolumn{1}{c|}{GE-PI} & GE \\ \hline
0 & \multicolumn{1}{c|}{0} & 0 & \multicolumn{1}{c|}{1} & 1 & \multicolumn{1}{c|}{0} & 1 & \multicolumn{1}{c|}{2} & 7 & \multicolumn{1}{c|}{0} & 31 & \multicolumn{1}{c|}{3} & 175 & \multicolumn{1}{c|}{0} & 981 & \multicolumn{1}{c|}{4} & 5719 \\ \hline
1 & \multicolumn{1}{c|}{0} & 0 & \multicolumn{1}{c|}{0} & 1 & \multicolumn{1}{c|}{1} & 3 & \multicolumn{1}{c|}{0} & 18 & \multicolumn{1}{c|}{2} & 90 & \multicolumn{1}{c|}{0} & 504 & \multicolumn{1}{c|}{4} & 2856 & \multicolumn{1}{c|}{1} & 16688 \\ \hline
2 & \multicolumn{1}{c|}{0} & 0 & \multicolumn{1}{c|}{1} & 1 & \multicolumn{1}{c|}{0} & 5 & \multicolumn{1}{c|}{2} & 26 & \multicolumn{1}{c|}{1} & 140 & \multicolumn{1}{c|}{4} & 780 & \multicolumn{1}{c|}{2} & 4480 & \multicolumn{1}{c|}{7} & 26320 \\ \hline
3 & \multicolumn{1}{c|}{1} & 1 & \multicolumn{1}{c|}{0} & 1 & \multicolumn{1}{c|}{2} & 7 & \multicolumn{1}{c|}{1} & 13 & \multicolumn{1}{c|}{4} & 175 & \multicolumn{1}{c|}{3} & 981 & \multicolumn{1}{c|}{7} & 5719 & \multicolumn{1}{c|}{5} & 33922 \\ \hline
4 & \multicolumn{1}{c|}{} &  & \multicolumn{1}{c|}{1} & 1 & \multicolumn{1}{c|}{1} & 6 & \multicolumn{1}{c|}{3} & 33 & \multicolumn{1}{c|}{2} & 189 & \multicolumn{1}{c|}{6} & 1095 & \multicolumn{1}{c|}{5} & 6489 & \multicolumn{1}{c|}{11} & 39046 \\ \hline
5 & \multicolumn{1}{c|}{} &  & \multicolumn{1}{c|}{0} & 1 & \multicolumn{1}{c|}{1} & 5 & \multicolumn{1}{c|}{1} & 32 & \multicolumn{1}{c|}{4} & 186 & \multicolumn{1}{c|}{3} & 1120 & \multicolumn{1}{c|}{8} & 6776 & \multicolumn{1}{c|}{7} & 41524 \\ \hline
6 & \multicolumn{1}{c|}{} &  & \multicolumn{1}{c|}{1} & 1 & \multicolumn{1}{c|}{1} & 4 & \multicolumn{1}{c|}{3} & 28 & \multicolumn{1}{c|}{3} & 170 & \multicolumn{1}{c|}{7} & 1064 & \multicolumn{1}{c|}{7} & 6621 & \multicolumn{1}{c|}{13} & 41468 \\ \hline
7 & \multicolumn{1}{c|}{} &  & \multicolumn{1}{c|}{} &  & \multicolumn{1}{c|}{1} & 3 & \multicolumn{1}{c|}{1} & 21 & \multicolumn{1}{c|}{4} & 145 & \multicolumn{1}{c|}{4} & 945 & \multicolumn{1}{c|}{9} & 6105 & \multicolumn{1}{c|}{9} & 39235 \\ \hline
8 & \multicolumn{1}{c|}{} &  & \multicolumn{1}{c|}{} &  & \multicolumn{1}{c|}{0} & 2 & \multicolumn{1}{c|}{2} & 15 & \multicolumn{1}{c|}{2} & 115 & \multicolumn{1}{c|}{6} & 791 & \multicolumn{1}{c|}{6} & 5334 & \multicolumn{1}{c|}{13} & 35357 \\ \hline
9 & \multicolumn{1}{c|}{} &  & \multicolumn{1}{c|}{} &  & \multicolumn{1}{c|}{1} & 1 & \multicolumn{1}{c|}{1} & 10 & \multicolumn{1}{c|}{3} & 84 & \multicolumn{1}{c|}{4} & 625 & \multicolumn{1}{c|}{9} & 4424 & \multicolumn{1}{c|}{10} & 30436 \\ \hline
10 & \multicolumn{1}{c|}{} &  & \multicolumn{1}{c|}{} &  & \multicolumn{1}{c|}{} &  & \multicolumn{1}{c|}{1} & 6 & \multicolumn{1}{c|}{2} & 56 & \multicolumn{1}{c|}{5} & 465 & \multicolumn{1}{c|}{6} & 3486 & \multicolumn{1}{c|}{12} & 25060 \\ \hline
11 & \multicolumn{1}{c|}{} &  & \multicolumn{1}{c|}{} &  & \multicolumn{1}{c|}{} &  & \multicolumn{1}{c|}{0} & 3 & \multicolumn{1}{c|}{2} & 35 & \multicolumn{1}{c|}{2} & 324 & \multicolumn{1}{c|}{7} & 2611 & \multicolumn{1}{c|}{8} & 19740 \\ \hline
12 & \multicolumn{1}{c|}{} &  & \multicolumn{1}{c|}{} &  & \multicolumn{1}{c|}{} &  & \multicolumn{1}{c|}{1} & 1 & \multicolumn{1}{c|}{1} & 20 & \multicolumn{1}{c|}{4} & 210 & \multicolumn{1}{c|}{5} & 1855 & \multicolumn{1}{c|}{11} & 14868 \\ \hline
13 & \multicolumn{1}{c|}{} &  & \multicolumn{1}{c|}{} &  & \multicolumn{1}{c|}{} &  & \multicolumn{1}{c|}{} &  & \multicolumn{1}{c|}{1} & 10 & \multicolumn{1}{c|}{2} & 126 & \multicolumn{1}{c|}{5} & 1245 & \multicolumn{1}{c|}{7} & 10696 \\ \hline
14 & \multicolumn{1}{c|}{} &  & \multicolumn{1}{c|}{} &  & \multicolumn{1}{c|}{} &  & \multicolumn{1}{c|}{} &  & \multicolumn{1}{c|}{0} & 4 & \multicolumn{1}{c|}{2} & 70 & \multicolumn{1}{c|}{3} & 785 & \multicolumn{1}{c|}{8} & 7336 \\ \hline
15 & \multicolumn{1}{c|}{} &  & \multicolumn{1}{c|}{} &  & \multicolumn{1}{c|}{} &  & \multicolumn{1}{c|}{} &  & \multicolumn{1}{c|}{1} & 1 & \multicolumn{1}{c|}{1} & 35 & \multicolumn{1}{c|}{4} & 462 & \multicolumn{1}{c|}{5} & 4781 \\ \hline
16 & \multicolumn{1}{c|}{} &  & \multicolumn{1}{c|}{} &  & \multicolumn{1}{c|}{} &  & \multicolumn{1}{c|}{} &  & \multicolumn{1}{c|}{} &  & \multicolumn{1}{c|}{1} & 15 & \multicolumn{1}{c|}{2} & 252 & \multicolumn{1}{c|}{6} & 2947 \\ \hline
17 & \multicolumn{1}{c|}{} &  & \multicolumn{1}{c|}{} &  & \multicolumn{1}{c|}{} &  & \multicolumn{1}{c|}{} &  & \multicolumn{1}{c|}{} &  & \multicolumn{1}{c|}{0} & 5 & \multicolumn{1}{c|}{2} & 126 & \multicolumn{1}{c|}{3} & 1708 \\ \hline
18 & \multicolumn{1}{c|}{} &  & \multicolumn{1}{c|}{} &  & \multicolumn{1}{c|}{} &  & \multicolumn{1}{c|}{} &  & \multicolumn{1}{c|}{} &  & \multicolumn{1}{c|}{1} & 1 & \multicolumn{1}{c|}{1} & 56 & \multicolumn{1}{c|}{4} & 924 \\ \hline
19 & \multicolumn{1}{c|}{} &  & \multicolumn{1}{c|}{} &  & \multicolumn{1}{c|}{} &  & \multicolumn{1}{c|}{} &  & \multicolumn{1}{c|}{} &  & \multicolumn{1}{c|}{} &  & \multicolumn{1}{c|}{1} & 21 & \multicolumn{1}{c|}{2} & 462 \\ \hline
20 & \multicolumn{1}{c|}{} &  & \multicolumn{1}{c|}{} &  & \multicolumn{1}{c|}{} &  & \multicolumn{1}{c|}{} &  & \multicolumn{1}{c|}{} &  & \multicolumn{1}{c|}{} &  & \multicolumn{1}{c|}{0} & 6 & \multicolumn{1}{c|}{2} & 210 \\ \hline
21 & \multicolumn{1}{c|}{} &  & \multicolumn{1}{c|}{} &  & \multicolumn{1}{c|}{} &  & \multicolumn{1}{c|}{} &  & \multicolumn{1}{c|}{} &  & \multicolumn{1}{c|}{} &  & \multicolumn{1}{c|}{1} & 1 & \multicolumn{1}{c|}{1} & 84 \\ \hline
22 & \multicolumn{1}{c|}{} &  & \multicolumn{1}{c|}{} &  & \multicolumn{1}{c|}{} &  & \multicolumn{1}{c|}{} &  & \multicolumn{1}{c|}{} &  & \multicolumn{1}{c|}{} &  & \multicolumn{1}{c|}{} &  & \multicolumn{1}{c|}{1} & 28 \\ \hline
23 & \multicolumn{1}{c|}{} &  & \multicolumn{1}{c|}{} &  & \multicolumn{1}{c|}{} &  & \multicolumn{1}{c|}{} &  & \multicolumn{1}{c|}{} &  & \multicolumn{1}{c|}{} &  & \multicolumn{1}{c|}{} &  & \multicolumn{1}{c|}{0} & 7 \\ \hline
24 & \multicolumn{1}{c|}{} &  & \multicolumn{1}{c|}{} &  & \multicolumn{1}{c|}{} &  & \multicolumn{1}{c|}{} &  & \multicolumn{1}{c|}{} &  & \multicolumn{1}{c|}{} &  & \multicolumn{1}{c|}{} &  & \multicolumn{1}{c|}{1} & 1 \\ \hline
\end{tabular}%
}
\caption{Dimensions of $V^{\bl,L}$ (GE) and $\bV^{\bl,L}$ (GE-PI) for $\bl = (3,3,\ldots,3)$ with length $\Nat$, and for $L\in\{0,1,\ldots,3\Nat\}$.}
\label{tab:l3}
\end{table}


\begin{table}[htb!]
\setlength{\tabcolsep}{7pt}
\resizebox{\textwidth}{!}{%
\begin{tabular}{|c|cc|cc|cc|cc|cc|cc|cc|cc|}
\hline
\multirow{2}{*}{$L/l$} & \multicolumn{2}{c|}{0} & \multicolumn{2}{c|}{1} & \multicolumn{2}{c|}{2} & \multicolumn{2}{c|}{3} & \multicolumn{2}{c|}{4} & \multicolumn{2}{c|}{5} & \multicolumn{2}{c|}{6} & \multicolumn{2}{c|}{7} \\ \cline{2-17} 
 & \multicolumn{1}{c|}{GE-PI} & GE & \multicolumn{1}{c|}{GE-PI} & GE & \multicolumn{1}{c|}{GE-PI} & GE & \multicolumn{1}{c|}{GE-PI} & GE & \multicolumn{1}{c|}{GE-PI} & GE & \multicolumn{1}{c|}{GE-PI} & GE & \multicolumn{1}{c|}{GE-PI} & GE & \multicolumn{1}{c|}{GE-PI} & GE \\ \hline
0 & \multicolumn{1}{c|}{1} & 1 & \multicolumn{1}{c|}{0} & 1 & \multicolumn{1}{c|}{1} & 1 & \multicolumn{1}{c|}{0} & 1 & \multicolumn{1}{c|}{1} & 1 & \multicolumn{1}{c|}{0} & 1 & \multicolumn{1}{c|}{1} & 1 & \multicolumn{1}{c|}{0} & 1 \\ \hline
1 & \multicolumn{1}{c|}{} &  & \multicolumn{1}{c|}{1} & 3 & \multicolumn{1}{c|}{0} & 3 & \multicolumn{1}{c|}{1} & 3 & \multicolumn{1}{c|}{0} & 3 & \multicolumn{1}{c|}{1} & 3 & \multicolumn{1}{c|}{0} & 3 & \multicolumn{1}{c|}{1} & 3 \\ \hline
2 & \multicolumn{1}{c|}{} &  & \multicolumn{1}{c|}{0} & 2 & \multicolumn{1}{c|}{1} & 5 & \multicolumn{1}{c|}{0} & 5 & \multicolumn{1}{c|}{1} & 5 & \multicolumn{1}{c|}{0} & 5 & \multicolumn{1}{c|}{1} & 5 & \multicolumn{1}{c|}{0} & 5 \\ \hline
3 & \multicolumn{1}{c|}{} &  & \multicolumn{1}{c|}{1} & 1 & \multicolumn{1}{c|}{1} & 4 & \multicolumn{1}{c|}{2} & 7 & \multicolumn{1}{c|}{1} & 7 & \multicolumn{1}{c|}{2} & 7 & \multicolumn{1}{c|}{1} & 7 & \multicolumn{1}{c|}{2} & 7 \\ \hline
4 & \multicolumn{1}{c|}{} &  & \multicolumn{1}{c|}{} &  & \multicolumn{1}{c|}{1} & 3 & \multicolumn{1}{c|}{1} & 6 & \multicolumn{1}{c|}{2} & 9 & \multicolumn{1}{c|}{1} & 9 & \multicolumn{1}{c|}{2} & 9 & \multicolumn{1}{c|}{1} & 9 \\ \hline
5 & \multicolumn{1}{c|}{} &  & \multicolumn{1}{c|}{} &  & \multicolumn{1}{c|}{0} & 2 & \multicolumn{1}{c|}{1} & 5 & \multicolumn{1}{c|}{1} & 8 & \multicolumn{1}{c|}{2} & 11 & \multicolumn{1}{c|}{1} & 11 & \multicolumn{1}{c|}{2} & 11 \\ \hline
6 & \multicolumn{1}{c|}{} &  & \multicolumn{1}{c|}{} &  & \multicolumn{1}{c|}{1} & 1 & \multicolumn{1}{c|}{1} & 4 & \multicolumn{1}{c|}{2} & 7 & \multicolumn{1}{c|}{2} & 10 & \multicolumn{1}{c|}{3} & 13 & \multicolumn{1}{c|}{2} & 13 \\ \hline
7 & \multicolumn{1}{c|}{} &  & \multicolumn{1}{c|}{} &  & \multicolumn{1}{c|}{} &  & \multicolumn{1}{c|}{1} & 3 & \multicolumn{1}{c|}{1} & 6 & \multicolumn{1}{c|}{2} & 9 & \multicolumn{1}{c|}{2} & 12 & \multicolumn{1}{c|}{3} & 15 \\ \hline
8 & \multicolumn{1}{c|}{} &  & \multicolumn{1}{c|}{} &  & \multicolumn{1}{c|}{} &  & \multicolumn{1}{c|}{0} & 2 & \multicolumn{1}{c|}{1} & 5 & \multicolumn{1}{c|}{1} & 8 & \multicolumn{1}{c|}{2} & 11 & \multicolumn{1}{c|}{2} & 14 \\ \hline
9 & \multicolumn{1}{c|}{} &  & \multicolumn{1}{c|}{} &  & \multicolumn{1}{c|}{} &  & \multicolumn{1}{c|}{1} & 1 & \multicolumn{1}{c|}{1} & 4 & \multicolumn{1}{c|}{2} & 7 & \multicolumn{1}{c|}{2} & 10 & \multicolumn{1}{c|}{3} & 13 \\ \hline
10 & \multicolumn{1}{c|}{} &  & \multicolumn{1}{c|}{} &  & \multicolumn{1}{c|}{} &  & \multicolumn{1}{c|}{} &  & \multicolumn{1}{c|}{1} & 3 & \multicolumn{1}{c|}{1} & 6 & \multicolumn{1}{c|}{2} & 9 & \multicolumn{1}{c|}{2} & 12 \\ \hline
11 & \multicolumn{1}{c|}{} &  & \multicolumn{1}{c|}{} &  & \multicolumn{1}{c|}{} &  & \multicolumn{1}{c|}{} &  & \multicolumn{1}{c|}{0} & 2 & \multicolumn{1}{c|}{1} & 5 & \multicolumn{1}{c|}{1} & 8 & \multicolumn{1}{c|}{2} & 11 \\ \hline
12 & \multicolumn{1}{c|}{} &  & \multicolumn{1}{c|}{} &  & \multicolumn{1}{c|}{} &  & \multicolumn{1}{c|}{} &  & \multicolumn{1}{c|}{1} & 1 & \multicolumn{1}{c|}{1} & 4 & \multicolumn{1}{c|}{2} & 7 & \multicolumn{1}{c|}{2} & 10 \\ \hline
13 & \multicolumn{1}{c|}{} &  & \multicolumn{1}{c|}{} &  & \multicolumn{1}{c|}{} &  & \multicolumn{1}{c|}{} &  & \multicolumn{1}{c|}{} &  & \multicolumn{1}{c|}{1} & 3 & \multicolumn{1}{c|}{1} & 6 & \multicolumn{1}{c|}{2} & 9 \\ \hline
14 & \multicolumn{1}{c|}{} &  & \multicolumn{1}{c|}{} &  & \multicolumn{1}{c|}{} &  & \multicolumn{1}{c|}{} &  & \multicolumn{1}{c|}{} &  & \multicolumn{1}{c|}{0} & 2 & \multicolumn{1}{c|}{1} & 5 & \multicolumn{1}{c|}{1} & 8 \\ \hline
15 & \multicolumn{1}{c|}{} &  & \multicolumn{1}{c|}{} &  & \multicolumn{1}{c|}{} &  & \multicolumn{1}{c|}{} &  & \multicolumn{1}{c|}{} &  & \multicolumn{1}{c|}{1} & 1 & \multicolumn{1}{c|}{1} & 4 & \multicolumn{1}{c|}{2} & 7 \\ \hline
16 & \multicolumn{1}{c|}{} &  & \multicolumn{1}{c|}{} &  & \multicolumn{1}{c|}{} &  & \multicolumn{1}{c|}{} &  & \multicolumn{1}{c|}{} &  & \multicolumn{1}{c|}{} &  & \multicolumn{1}{c|}{1} & 3 & \multicolumn{1}{c|}{1} & 6 \\ \hline
17 & \multicolumn{1}{c|}{} &  & \multicolumn{1}{c|}{} &  & \multicolumn{1}{c|}{} &  & \multicolumn{1}{c|}{} &  & \multicolumn{1}{c|}{} &  & \multicolumn{1}{c|}{} &  & \multicolumn{1}{c|}{0} & 2 & \multicolumn{1}{c|}{1} & 5 \\ \hline
18 & \multicolumn{1}{c|}{} &  & \multicolumn{1}{c|}{} &  & \multicolumn{1}{c|}{} &  & \multicolumn{1}{c|}{} &  & \multicolumn{1}{c|}{} &  & \multicolumn{1}{c|}{} &  & \multicolumn{1}{c|}{1} & 1 & \multicolumn{1}{c|}{1} & 4 \\ \hline
19 & \multicolumn{1}{c|}{} &  & \multicolumn{1}{c|}{} &  & \multicolumn{1}{c|}{} &  & \multicolumn{1}{c|}{} &  & \multicolumn{1}{c|}{} &  & \multicolumn{1}{c|}{} &  & \multicolumn{1}{c|}{} &  & \multicolumn{1}{c|}{1} & 3 \\ \hline
20 & \multicolumn{1}{c|}{} &  & \multicolumn{1}{c|}{} &  & \multicolumn{1}{c|}{} &  & \multicolumn{1}{c|}{} &  & \multicolumn{1}{c|}{} &  & \multicolumn{1}{c|}{} &  & \multicolumn{1}{c|}{} &  & \multicolumn{1}{c|}{0} & 2 \\ \hline
21 & \multicolumn{1}{c|}{} &  & \multicolumn{1}{c|}{} &  & \multicolumn{1}{c|}{} &  & \multicolumn{1}{c|}{} &  & \multicolumn{1}{c|}{} &  & \multicolumn{1}{c|}{} &  & \multicolumn{1}{c|}{} &  & \multicolumn{1}{c|}{1} & 1 \\ \hline
\end{tabular}%
}
\caption{Dimensions of $V^{\bl,L}$ (GE) and $\bV^{\bl,L}$ (GE-PI) for $\bl=(l,l,l)$ with different $l$ and $L$.}
\label{tab:n3}
\end{table}


\begin{table}[htb!]
\setlength{\tabcolsep}{7pt}
\centering
\resizebox{\textwidth}{!}{%
\begin{tabular}{|c|cc|cc|cc|cc|cc|cc|cc|cc|}
\hline
\multirow{2}{*}{$L/l$} & \multicolumn{2}{c|}{0} & \multicolumn{2}{c|}{1} & \multicolumn{2}{c|}{2} & \multicolumn{2}{c|}{3} & \multicolumn{2}{c|}{4} & \multicolumn{2}{c|}{5} & \multicolumn{2}{c|}{6} & \multicolumn{2}{c|}{7} \\ \cline{2-17} 
 & \multicolumn{1}{c|}{GE-PI} & GE & \multicolumn{1}{c|}{GE-PI} & GE & \multicolumn{1}{c|}{GE-PI} & GE & \multicolumn{1}{c|}{GE-PI} & GE & \multicolumn{1}{c|}{GE-PI} & GE & \multicolumn{1}{c|}{GE-PI} & GE & \multicolumn{1}{c|}{GE-PI} & GE & \multicolumn{1}{c|}{GE-PI} & GE \\ \hline
0 & \multicolumn{1}{c|}{1} & 1 & \multicolumn{1}{c|}{1} & 3 & \multicolumn{1}{c|}{1} & 5 & \multicolumn{1}{c|}{2} & 7 & \multicolumn{1}{c|}{2} & 9 & \multicolumn{1}{c|}{2} & 11 & \multicolumn{1}{c|}{3} & 13 & \multicolumn{1}{c|}{3} & 15 \\ \hline
1 & \multicolumn{1}{c|}{} &  & \multicolumn{1}{c|}{0} & 6 & \multicolumn{1}{c|}{0} & 12 & \multicolumn{1}{c|}{0} & 18 & \multicolumn{1}{c|}{0} & 24 & \multicolumn{1}{c|}{0} & 30 & \multicolumn{1}{c|}{0} & 36 & \multicolumn{1}{c|}{0} & 42 \\ \hline
2 & \multicolumn{1}{c|}{} &  & \multicolumn{1}{c|}{1} & 6 & \multicolumn{1}{c|}{2} & 16 & \multicolumn{1}{c|}{2} & 26 & \multicolumn{1}{c|}{3} & 36 & \multicolumn{1}{c|}{4} & 46 & \multicolumn{1}{c|}{4} & 56 & \multicolumn{1}{c|}{5} & 66 \\ \hline
3 & \multicolumn{1}{c|}{} &  & \multicolumn{1}{c|}{0} & 3 & \multicolumn{1}{c|}{0} & 17 & \multicolumn{1}{c|}{1} & 31 & \multicolumn{1}{c|}{1} & 45 & \multicolumn{1}{c|}{1} & 59 & \multicolumn{1}{c|}{2} & 73 & \multicolumn{1}{c|}{2} & 87 \\ \hline
4 & \multicolumn{1}{c|}{} &  & \multicolumn{1}{c|}{1} & 1 & \multicolumn{1}{c|}{2} & 15 & \multicolumn{1}{c|}{3} & 33 & \multicolumn{1}{c|}{4} & 51 & \multicolumn{1}{c|}{5} & 69 & \multicolumn{1}{c|}{6} & 87 & \multicolumn{1}{c|}{7} & 105 \\ \hline
5 & \multicolumn{1}{c|}{} &  & \multicolumn{1}{c|}{} &  & \multicolumn{1}{c|}{1} & 10 & \multicolumn{1}{c|}{1} & 32 & \multicolumn{1}{c|}{2} & 54 & \multicolumn{1}{c|}{3} & 76 & \multicolumn{1}{c|}{3} & 98 & \multicolumn{1}{c|}{4} & 120 \\ \hline
6 & \multicolumn{1}{c|}{} &  & \multicolumn{1}{c|}{} &  & \multicolumn{1}{c|}{1} & 6 & \multicolumn{1}{c|}{3} & 28 & \multicolumn{1}{c|}{4} & 54 & \multicolumn{1}{c|}{5} & 80 & \multicolumn{1}{c|}{7} & 106 & \multicolumn{1}{c|}{8} & 132 \\ \hline
7 & \multicolumn{1}{c|}{} &  & \multicolumn{1}{c|}{} &  & \multicolumn{1}{c|}{0} & 3 & \multicolumn{1}{c|}{1} & 21 & \multicolumn{1}{c|}{2} & 51 & \multicolumn{1}{c|}{3} & 81 & \multicolumn{1}{c|}{4} & 111 & \multicolumn{1}{c|}{5} & 141 \\ \hline
8 & \multicolumn{1}{c|}{} &  & \multicolumn{1}{c|}{} &  & \multicolumn{1}{c|}{1} & 1 & \multicolumn{1}{c|}{2} & 15 & \multicolumn{1}{c|}{4} & 45 & \multicolumn{1}{c|}{6} & 79 & \multicolumn{1}{c|}{7} & 113 & \multicolumn{1}{c|}{9} & 147 \\ \hline
9 & \multicolumn{1}{c|}{} &  & \multicolumn{1}{c|}{} &  & \multicolumn{1}{c|}{} &  & \multicolumn{1}{c|}{1} & 10 & \multicolumn{1}{c|}{2} & 36 & \multicolumn{1}{c|}{3} & 74 & \multicolumn{1}{c|}{5} & 112 & \multicolumn{1}{c|}{6} & 150 \\ \hline
10 & \multicolumn{1}{c|}{} &  & \multicolumn{1}{c|}{} &  & \multicolumn{1}{c|}{} &  & \multicolumn{1}{c|}{1} & 6 & \multicolumn{1}{c|}{3} & 28 & \multicolumn{1}{c|}{5} & 66 & \multicolumn{1}{c|}{7} & 108 & \multicolumn{1}{c|}{9} & 150 \\ \hline
11 & \multicolumn{1}{c|}{} &  & \multicolumn{1}{c|}{} &  & \multicolumn{1}{c|}{} &  & \multicolumn{1}{c|}{0} & 3 & \multicolumn{1}{c|}{1} & 21 & \multicolumn{1}{c|}{3} & 55 & \multicolumn{1}{c|}{4} & 101 & \multicolumn{1}{c|}{6} & 147 \\ \hline
12 & \multicolumn{1}{c|}{} &  & \multicolumn{1}{c|}{} &  & \multicolumn{1}{c|}{} &  & \multicolumn{1}{c|}{1} & 1 & \multicolumn{1}{c|}{2} & 15 & \multicolumn{1}{c|}{4} & 45 & \multicolumn{1}{c|}{7} & 91 & \multicolumn{1}{c|}{9} & 141 \\ \hline
13 & \multicolumn{1}{c|}{} &  & \multicolumn{1}{c|}{} &  & \multicolumn{1}{c|}{} &  & \multicolumn{1}{c|}{} &  & \multicolumn{1}{c|}{1} & 10 & \multicolumn{1}{c|}{2} & 36 & \multicolumn{1}{c|}{4} & 78 & \multicolumn{1}{c|}{6} & 132 \\ \hline
14 & \multicolumn{1}{c|}{} &  & \multicolumn{1}{c|}{} &  & \multicolumn{1}{c|}{} &  & \multicolumn{1}{c|}{} &  & \multicolumn{1}{c|}{1} & 6 & \multicolumn{1}{c|}{3} & 28 & \multicolumn{1}{c|}{5} & 66 & \multicolumn{1}{c|}{8} & 120 \\ \hline
15 & \multicolumn{1}{c|}{} &  & \multicolumn{1}{c|}{} &  & \multicolumn{1}{c|}{} &  & \multicolumn{1}{c|}{} &  & \multicolumn{1}{c|}{0} & 3 & \multicolumn{1}{c|}{1} & 21 & \multicolumn{1}{c|}{3} & 55 & \multicolumn{1}{c|}{5} & 105 \\ \hline
16 & \multicolumn{1}{c|}{} &  & \multicolumn{1}{c|}{} &  & \multicolumn{1}{c|}{} &  & \multicolumn{1}{c|}{} &  & \multicolumn{1}{c|}{1} & 1 & \multicolumn{1}{c|}{2} & 15 & \multicolumn{1}{c|}{4} & 45 & \multicolumn{1}{c|}{7} & 91 \\ \hline
17 & \multicolumn{1}{c|}{} &  & \multicolumn{1}{c|}{} &  & \multicolumn{1}{c|}{} &  & \multicolumn{1}{c|}{} &  & \multicolumn{1}{c|}{} &  & \multicolumn{1}{c|}{1} & 10 & \multicolumn{1}{c|}{2} & 36 & \multicolumn{1}{c|}{4} & 78 \\ \hline
18 & \multicolumn{1}{c|}{} &  & \multicolumn{1}{c|}{} &  & \multicolumn{1}{c|}{} &  & \multicolumn{1}{c|}{} &  & \multicolumn{1}{c|}{} &  & \multicolumn{1}{c|}{1} & 6 & \multicolumn{1}{c|}{3} & 28 & \multicolumn{1}{c|}{5} & 66 \\ \hline
19 & \multicolumn{1}{c|}{} &  & \multicolumn{1}{c|}{} &  & \multicolumn{1}{c|}{} &  & \multicolumn{1}{c|}{} &  & \multicolumn{1}{c|}{} &  & \multicolumn{1}{c|}{0} & 3 & \multicolumn{1}{c|}{1} & 21 & \multicolumn{1}{c|}{3} & 55 \\ \hline
20 & \multicolumn{1}{c|}{} &  & \multicolumn{1}{c|}{} &  & \multicolumn{1}{c|}{} &  & \multicolumn{1}{c|}{} &  & \multicolumn{1}{c|}{} &  & \multicolumn{1}{c|}{1} & 1 & \multicolumn{1}{c|}{2} & 15 & \multicolumn{1}{c|}{4} & 45 \\ \hline
21 & \multicolumn{1}{c|}{} &  & \multicolumn{1}{c|}{} &  & \multicolumn{1}{c|}{} &  & \multicolumn{1}{c|}{} &  & \multicolumn{1}{c|}{} &  & \multicolumn{1}{c|}{} &  & \multicolumn{1}{c|}{1} & 10 & \multicolumn{1}{c|}{2} & 36 \\ \hline
22 & \multicolumn{1}{c|}{} &  & \multicolumn{1}{c|}{} &  & \multicolumn{1}{c|}{} &  & \multicolumn{1}{c|}{} &  & \multicolumn{1}{c|}{} &  & \multicolumn{1}{c|}{} &  & \multicolumn{1}{c|}{1} & 6 & \multicolumn{1}{c|}{3} & 28 \\ \hline
23 & \multicolumn{1}{c|}{} &  & \multicolumn{1}{c|}{} &  & \multicolumn{1}{c|}{} &  & \multicolumn{1}{c|}{} &  & \multicolumn{1}{c|}{} &  & \multicolumn{1}{c|}{} &  & \multicolumn{1}{c|}{0} & 3 & \multicolumn{1}{c|}{1} & 21 \\ \hline
24 & \multicolumn{1}{c|}{} &  & \multicolumn{1}{c|}{} &  & \multicolumn{1}{c|}{} &  & \multicolumn{1}{c|}{} &  & \multicolumn{1}{c|}{} &  & \multicolumn{1}{c|}{} &  & \multicolumn{1}{c|}{1} & 1 & \multicolumn{1}{c|}{2} & 15 \\ \hline
25 & \multicolumn{1}{c|}{} &  & \multicolumn{1}{c|}{} &  & \multicolumn{1}{c|}{} &  & \multicolumn{1}{c|}{} &  & \multicolumn{1}{c|}{} &  & \multicolumn{1}{c|}{} &  & \multicolumn{1}{c|}{} &  & \multicolumn{1}{c|}{1} & 10 \\ \hline
26 & \multicolumn{1}{c|}{} &  & \multicolumn{1}{c|}{} &  & \multicolumn{1}{c|}{} &  & \multicolumn{1}{c|}{} &  & \multicolumn{1}{c|}{} &  & \multicolumn{1}{c|}{} &  & \multicolumn{1}{c|}{} &  & \multicolumn{1}{c|}{1} & 6 \\ \hline
27 & \multicolumn{1}{c|}{} &  & \multicolumn{1}{c|}{} &  & \multicolumn{1}{c|}{} &  & \multicolumn{1}{c|}{} &  & \multicolumn{1}{c|}{} &  & \multicolumn{1}{c|}{} &  & \multicolumn{1}{c|}{} &  & \multicolumn{1}{c|}{0} & 3 \\ \hline
28 & \multicolumn{1}{c|}{} &  & \multicolumn{1}{c|}{} &  & \multicolumn{1}{c|}{} &  & \multicolumn{1}{c|}{} &  & \multicolumn{1}{c|}{} &  & \multicolumn{1}{c|}{} &  & \multicolumn{1}{c|}{} &  & \multicolumn{1}{c|}{1} & 1 \\ \hline
\end{tabular}%
}
\caption{Dimensions of $V^{\bl,L}$ (GE) and $\bV^{\bl,L}$ (GE-PI) for $\bl=(l,l,l,l)$ with different $l$ and $L$.}
\label{tab:n4}
\end{table}

\end{document}
